\documentclass[11pt]{article}

\usepackage[margin=1in]{geometry}
\usepackage{amsmath,amssymb,amsthm}
\usepackage{enumitem}

\newtheorem{theorem}{Theorem}
\newtheorem{lemma}{Lemma}
\newtheorem{proposition}{Proposition}
\newtheorem{corollary}{Corollary}
\newtheorem{conjecture}{Conjecture}
\newtheorem{definition}{Definition}
\newtheorem{remark}{Remark}

\newcommand{\Deck}{\operatorname{Deck}}
\newcommand{\PDeck}{\operatorname{PDeck}}
\newcommand{\Ext}{\operatorname{Ext}}
\newcommand{\Aut}{\operatorname{Aut}}
\newcommand{\Iso}{\operatorname{Iso}}
\newcommand{\defect}{\operatorname{def}}
\newcommand{\Err}{\operatorname{Err}}
\newcommand{\V}{V}
\newcommand{\E}{E}
\newcommand{\one}{\mathbf{1}}

\title{A Centered-Extension and Defect-Minimization Strategy for Graph Reconstruction}
\author{Research sketch}
\date{\today}

\begin{document}
\maketitle

\begin{abstract}
This note records an English proof strategy for the graph reconstruction
conjecture.  The main reframing is to fix one common card and view two
reconstructions as one-vertex extensions of that common card.  The problem
then becomes: does the deck determine the attachment set of the missing
vertex?  The proposed proof uses a defect set between two attachment sets and
tries to show that any nonzero defect can be repaired by moving the deleted
``hole'' to a better card alignment.  The sharp local form of the problem is
that a matched rooted card gives a color-preserving bijection whose graph
error is concentrated in one star; the desired descent says that some such
one-star error is smaller than the original color defect.
\end{abstract}

\section{The centered-extension reduction}

Let \(C\) be a finite simple graph with vertex set \(W\).  For a subset
\(A \subseteq W\), define \(\Ext(C,A)\) to be the graph obtained from \(C\)
by adding one new vertex \(\ast\), called the root, adjacent exactly to the
vertices of \(A\).

Thus
\[
  V(\Ext(C,A)) = W \sqcup \{\ast\},
\]
and
\[
  E(\Ext(C,A)) =
    E(C) \cup \{\{\ast,u\} : u \in A\}.
\]

For \(u \in W\), deleting the old vertex \(u\) gives
\[
  \Ext(C,A)-u \cong \Ext(C-u, A \setminus \{u\}).
\]
Deleting the root gives
\[
  \Ext(C,A)-\ast \cong C.
\]

\begin{definition}[Punctured extension deck]
The punctured extension deck of the attachment set \(A\) over \(C\) is the
multiset
\[
  \PDeck_C(A)
  =
  \{\!\{\, \Ext(C-u, A\setminus\{u\}) : u \in W \,\}\!\},
\]
where entries are taken up to graph isomorphism.
\end{definition}

Then the full vertex deck of \(\Ext(C,A)\) is
\[
  \Deck(\Ext(C,A)) = \{\!\{ C \}\!\} + \PDeck_C(A).
\]

\begin{proposition}[Centered reduction]
Suppose \(G\) and \(H\) are graphs on \(n \geq 3\) vertices with the same
deck.  Choose a matched card \(G-x \cong H-y\), and identify both cards with a
single graph \(C\).  Then there are subsets \(A,B \subseteq V(C)\) such that
\[
  G \cong \Ext(C,A),
  \qquad
  H \cong \Ext(C,B),
\]
and
\[
  \PDeck_C(A) = \PDeck_C(B).
\]
\end{proposition}

\begin{proof}[Proof sketch]
The subset \(A\) is the neighborhood of \(x\), transported into the card
\(G-x \cong C\).  Similarly \(B\) is the neighborhood of \(y\), transported
into \(H-y \cong C\).  Since the full decks of \(G\) and \(H\) agree and both
contain the chosen matched card \(C\), cancel one occurrence of \(C\) from the
two deck multisets.
\end{proof}

Thus the reconstruction conjecture follows from the following centered theorem.

\begin{conjecture}[Centered extension theorem]
For every finite graph \(C\) and subsets \(A,B \subseteq V(C)\), if
\[
  \PDeck_C(A) = \PDeck_C(B),
\]
then
\[
  \Ext(C,A) \cong \Ext(C,B).
\]
\end{conjecture}

\begin{conjecture}[Root-preserving centered theorem]
If
\[
  \PDeck_C(A) = \PDeck_C(B),
\]
then \(A\) and \(B\) differ by
an automorphism of \(C\):
\[
  \exists \alpha \in \Aut(C), \quad \alpha(A)=B.
\]
\end{conjecture}

\begin{remark}
This root-preserving version implies the centered extension theorem.  It is a
stronger target because an isomorphism
\(\Ext(C,A) \cong \Ext(C,B)\) might in principle send the root of one extension
to a non-root vertex of the other.  Computationally, however, the stronger
statement survives the small cases below.  For now we treat it as the primary
target, while retaining root-moving isomorphisms as a possible escape hatch in
the moving-hole argument.
\end{remark}

\begin{remark}[Computational evidence]
A brute-force search found no counterexample to the root-preserving centered
theorem for all labeled cards \(C\) on at most five vertices, and no
counterexample for the 156 unlabeled cards \(C\) on six vertices.  The search
buckets all attachment sets \(A\subseteq V(C)\) by \(\PDeck_C(A)\) and checks
that every bucket is a single \(\Aut(C)\)-orbit.  This is only evidence, but it
is useful: the stronger theorem survives the first small cases where nontrivial
attachment-set buckets are abundant.
\end{remark}

\begin{remark}[Computational evidence for descent]
In the same search range, a stronger one-step phenomenon held.  Whenever
\(\PDeck_C(A)=\PDeck_C(B)\) and \(A\neq B\), there was a defect vertex
\(u\in A\triangle B\), a vertex \(v\in V(C)\), and a root-preserving
isomorphism
\[
  \Ext(C-u,A\setminus\{u\})
  \cong
  \Ext(C-v,B\setminus\{v\})
\]
whose recentering strictly reduced \(|A\triangle B|\).  For the 156 unlabeled
six-vertex cards, this covered 11,008 nontrivial same-punctured-deck pairs.
This suggests that the raw defect \(|A\triangle B|\) may itself be the right
induction parameter.
\end{remark}

\begin{remark}[Computational evidence for the single-switch base case]
A separate check of the single-switch situation
\[
  A=S\cup\{p\},\qquad B=S\cup\{q\}
\]
found no failure of the zero-error-match prediction for all unlabeled cards
\(C\) on at most six vertices.  Among the 156 unlabeled six-vertex cards,
there were 37,440 single-switch pairs, 7,328 of which had equal punctured
extension decks.  Every same-punctured-deck pair had a root-preserving
direct \(p\leftrightarrow q\) match with zero star error.  Exact two-error
matches also occurred, in 1,240 of the same-punctured-deck pairs, but always
alongside this direct zero-error companion.  This is useful evidence for the
corrected base-case target: exact two-error matches are allowed, but they
should never be the only available defect-rooted matches.
A follow-up check showed that detouring matches also occur frequently: among
the six-vertex same-punctured-deck single switches, 3,000 admitted a
root-preserving match from a defect vertex to a vertex of \(S\).  Thus detours
are not obstructions by themselves; the evidence is instead that a direct
zero-error representative is always also available.
A classification of the same six-vertex data found that all 7,328
same-punctured-deck single switches satisfy
\(\deg_S(p)=\deg_S(q)\); 5,568 are settled by twin collapse, while 1,760 are
non-twin and therefore represent the genuinely interesting balanced witness
case.
\end{remark}

\section{Defect}

Fix \(C\) and two attachment sets \(A,B \subseteq V(C)\).  Define the defect
set
\[
  D(A,B) = A \triangle B.
\]
Equivalently,
\[
  D_+ = A \setminus B,
  \qquad
  D_- = B \setminus A,
  \qquad
  D(A,B) = D_+ \sqcup D_-.
\]
If \(D(A,B)=\varnothing\), then the identity on \(C\), extended by root to
root, is an isomorphism
\[
  \Ext(C,A) \cong \Ext(C,B).
\]

The plan is to prove that a nonempty defect cannot survive equality of
punctured extension decks.

\begin{conjecture}[No nontrivial centered trade]
If \(\PDeck_C(A)=\PDeck_C(B)\), then either
\[
  \Ext(C,A) \cong \Ext(C,B),
\]
or there is another centered presentation of the same two extension graphs
with strictly smaller defect.
\end{conjecture}

Since defect size is finite, repeated repair would terminate at defect zero or
at an isomorphism that moves the root.  Either outcome proves reconstruction.

\section{First shadow: edge counts}

The defect is visible in aligned punctured shadows.

\begin{lemma}[Aligned edge-count defect]
For each \(u \in V(C)\),
\[
  |E(\Ext(C-u,A\setminus\{u\}))|
  -
  |E(\Ext(C-u,B\setminus\{u\}))|
  =
  \one_{u \in B} - \one_{u \in A}.
\]
\end{lemma}

\begin{proof}
Both graphs have the same base \(C-u\).  The first root has
\(|A|-\one_{u\in A}\) incident edges, while the second root has
\(|B|-\one_{u\in B}\) incident edges.  Since equality of the full decks of
\(\Ext(C,A)\) and \(\Ext(C,B)\) implies equality of their total edge counts,
we have \(|A|=|B|\).  Subtracting gives the displayed formula.
\end{proof}

Thus vertices in \(A\setminus B\) and vertices in \(B\setminus A\) create
opposite edge-count errors in the aligned shadows.  Equality of punctured
decks can hide these errors only by rematching cards.

\begin{remark}
This is the first conservation law.  It cannot by itself prove the theorem:
the deck is an unlabelled multiset, so the card indexed by \(u\) on the
\(A\)-side may be matched to a card indexed by some different \(v\) on the
\(B\)-side.  The proof must exploit the entire isomorphism type of each
punctured extension, not just its number of edges.
\end{remark}

\section{Higher shadows}

Let \(F\) be any finite graph.  Consider the difference
\[
  s_F(\Ext(C-u,A\setminus\{u\}))
  -
  s_F(\Ext(C-u,B\setminus\{u\})),
\]
where \(s_F(X)\) denotes the number of induced copies of \(F\) in \(X\).
Copies of \(F\) that avoid the root depend only on \(C-u\) and cancel.  The
entire difference is caused by copies of \(F\) using the root.  Hence it is a
function of the rooted neighborhood data \(A\setminus\{u\}\) versus
\(B\setminus\{u\}\) inside \(C-u\).

This gives a family of constraints:
\[
  \begin{aligned}
  &\{\!\{\,
      \text{rooted pattern profile of }(C-u,A\setminus\{u\})
      : u \in V(C)
    \,\}\!\}
  \\
  &\qquad =
    \{\!\{\,
      \text{rooted pattern profile of }(C-u,B\setminus\{u\})
      : u \in V(C)
    \,\}\!\}.
  \end{aligned}
\]
Equivalently, for every finite pattern \(F\), let
\(r_F(C-u,A\setminus\{u\})\) be the number of induced copies of \(F\) in
\(\Ext(C-u,A\setminus\{u\})\) that use the root.  Then
\[
  \{\!\{\, r_F(C-u,A\setminus\{u\}) : u \in V(C) \,\}\!\}
  =
  \{\!\{\, r_F(C-u,B\setminus\{u\}) : u \in V(C) \,\}\!\}.
\]

The centered theorem says that no non-isomorphic pair of attachment sets can
satisfy all of these constraints simultaneously.

\section{The moving-hole repair lemma}

The most promising proof mechanism is to use an isomorphism between two
matched punctured shadows to move the deleted vertex, or ``hole''.

Suppose
\[
  \theta :
  \Ext(C-u,A\setminus\{u\})
  \cong
  \Ext(C-v,B\setminus\{v\})
\]
is one of the card isomorphisms supplied by
\(\PDeck_C(A)=\PDeck_C(B)\).

Recenter both full extensions at the deleted old vertices \(u\) and \(v\).
The new common card is the left punctured extension
\[
  D = \Ext(C-u,A\setminus\{u\}),
\]
and the right punctured extension is identified with \(D\) by
\(\theta^{-1}\).

In the left full graph \(\Ext(C,A)\), the missing vertex \(u\) is adjacent in
the new card \(D\) to
\[
  A_u^{\mathrm{new}}
  =
  N_C(u)\setminus\{u\}
  \;\cup\;
  \begin{cases}
    \{\ast_A\}, & u \in A,\\
    \varnothing, & u \notin A,
  \end{cases}
\]
where \(\ast_A\) is the root of the left extension.  In the right full graph,
the missing vertex \(v\) is adjacent in its punctured card to
\[
  B_v^{\mathrm{raw}}
  =
  N_C(v)\setminus\{v\}
  \;\cup\;
  \begin{cases}
    \{\ast_B\}, & v \in B,\\
    \varnothing, & v \notin B.
  \end{cases}
\]
After transporting to the common card \(D\), the right attachment set is
\[
  B_v^{\mathrm{new}} = \theta^{-1}(B_v^{\mathrm{raw}}).
\]
Thus the recentered defect is
\[
  \Delta_{\theta}(u,v)
  =
  A_u^{\mathrm{new}} \triangle B_v^{\mathrm{new}}.
\]

The desired repair statement is therefore concrete:
\[
  |\Delta_{\theta}(u,v)| < |A\triangle B|
\]
unless the recentering already exhibits an isomorphism of the full extensions.

There are two cases.

\subsection{Case 1: the root maps to the root}

Then \(\theta\) restricts to an isomorphism
\[
  \psi : C-u \cong C-v
\]
and carries \(A\setminus\{u\}\) to \(B\setminus\{v\}\).
The new defect formula simplifies to
\[
\begin{aligned}
  \Delta_{\theta}(u,v)
  &=
  \bigl(N_C(u)\setminus\{u\}\bigr)
  \triangle
  \psi^{-1}\bigl(N_C(v)\setminus\{v\}\bigr).
\end{aligned}
\]
Indeed, because \(\theta\) maps root to root, it preserves root degree:
\[
  |A|-\one_{u\in A} = |B|-\one_{v\in B}.
\]
Since the full extensions have equal edge count, \(|A|=|B|\), and hence
\(\one_{u\in A}=\one_{v\in B}\).  Thus the possible root term cancels.

For an arbitrary unrooted card isomorphism, without assuming root maps to root,
we only get the weaker conservation law
\[
  \deg_C(u)+\one_{u\in A}
  =
  \deg_C(v)+\one_{v\in B}.
\]
Equivalently, if root membership changes in an unrooted match, it is
compensated by a one-unit change in the old-card degree.

\begin{lemma}[Empty recentered defect extends]
In the root-preserving case, if
\[
  \Delta_{\theta}(u,v)=\varnothing,
\]
then
\[
  \Ext(C,A)\cong \Ext(C,B).
\]
\end{lemma}

\begin{proof}[Proof sketch]
The equality \(\Delta_{\theta}(u,v)=\varnothing\) says that \(\psi\) carries
the neighborhood of \(u\) in \(C\) to the neighborhood of \(v\) in \(C\).
Therefore \(\psi\) extends over \(u\mapsto v\) to an isomorphism \(C\cong C\).
The root-preserving card isomorphism also gives
\(\psi(A\setminus\{u\})=B\setminus\{v\}\), and root-degree preservation gives
\(u\in A\) if and only if \(v\in B\).  Hence the extended isomorphism carries
\(A\) to \(B\), and then extends root-to-root to an isomorphism of
\(\Ext(C,A)\) and \(\Ext(C,B)\).
\end{proof}

\subsection{Star-error reformulation}

The root-preserving case can be stated more cleanly without mentioning the
new centered graph \(D\).  Given a root-preserving card isomorphism
\[
  \theta:
  \Ext(C-u,A\setminus\{u\})
  \cong
  \Ext(C-v,B\setminus\{v\}),
\]
let
\[
  \psi:C-u\cong C-v
\]
be its restriction to old vertices, and extend \(\psi\) to a bijection
\[
  \widehat\psi:V(C)\to V(C)
\]
by setting \(\widehat\psi(u)=v\).

\begin{lemma}[A rooted card match is a one-star graph error]
In the root-preserving case, \(\widehat\psi(A)=B\).  Moreover
\(\widehat\psi\) preserves every edge and non-edge of \(C\) except possibly
pairs incident to \(u\).  Its failure to be an automorphism is exactly
\[
  \Err_u(\widehat\psi)
  =
  \{\,w\in V(C)\setminus\{u\} :
      (uw\in E(C))
      \not\leftrightarrow
      (\widehat\psi(u)\widehat\psi(w)\in E(C))
  \,\}.
\]
Under the recentering identification,
\[
  |\Err_u(\widehat\psi)| = |\Delta_\theta(u,v)|.
\]
\end{lemma}

\begin{proof}[Proof sketch]
The card isomorphism gives
\(\psi(A\setminus\{u\})=B\setminus\{v\}\).  Root degree preservation, together
with \(|A|=|B|\), gives \(u\in A\) if and only if \(v\in B\).  Hence the
extension \(\widehat\psi\) carries all of \(A\) to all of \(B\).  Since
\(\psi\) is an isomorphism \(C-u\cong C-v\), every adjacency relation not
involving \(u\) is preserved.  The only possible failures are the adjacencies
from \(u\), and the set of such failures is precisely the neighborhood
symmetric difference appearing in \(\Delta_\theta(u,v)\).
\end{proof}

\begin{lemma}[Parity of star error]
In the same root-preserving situation,
\[
  \deg_C(u)=\deg_C(v),
\]
and therefore
\[
  |\Err_u(\widehat\psi)|
  =
  \left|
    \bigl(N_C(u)\setminus\{u\}\bigr)
    \triangle
    \psi^{-1}\bigl(N_C(v)\setminus\{v\}\bigr)
  \right|
\]
is even.
\end{lemma}

\begin{proof}[Proof sketch]
The restricted map \(\psi:C-u\cong C-v\) gives
\[
  |E(C-u)|=|E(C-v)|.
\]
Since
\[
  |E(C-u)|=|E(C)|-\deg_C(u)
  \quad\text{and}\quad
  |E(C-v)|=|E(C)|-\deg_C(v),
\]
the deleted vertices have equal degree.  The two neighborhoods compared in
the symmetric difference therefore have equal cardinality, so their symmetric
difference has even size.
\end{proof}

The parity lemma exposes a genuine base case.  If
\[
  |A\triangle B|=2,
\]
then a strict descent move must have star error \(0\), not merely a smaller
positive error.  Thus the first nontrivial local theorem is the following.

\begin{conjecture}[Single-switch rigidity]
Suppose
\[
  A=S\cup\{p\},
  \qquad
  B=S\cup\{q\},
  \qquad
  p\neq q,
\]
with \(p,q\notin S\), and suppose
\[
  \PDeck_C(A)=\PDeck_C(B).
\]
Then \(\Ext(C,A)\cong\Ext(C,B)\).  In the root-preserving version, there is
an automorphism of \(C\) carrying the set \(S\cup\{p\}\) to the set
\(S\cup\{q\}\).
\end{conjecture}

This should be treated as the base case for the whole descent strategy.
Computationally it is always resolved by a zero-error card match in the tested
range.  Conceptually, it says that the deck cannot hide the movement of a
single colored vertex unless the card already has the corresponding symmetry.

\subsection{Single-switch constraints}

For the rest of the single-switch analysis, work in the root-preserving
regime: either root-moving matches have already given success/ambiguity
descent, or we have extracted a root-preserving trade.  The rooted card-match
graph and all root-degree arguments below are statements in that regime.

In the single-switch situation, the two exceptional aligned shadows are
visible even before using the full isomorphism type of the cards.  For
\[
  A=S\cup\{p\},
  \qquad
  B=S\cup\{q\},
\]
the aligned edge-count defect is
\[
  |E(\Ext(C-p,S))|-|E(\Ext(C-p,S\cup\{q\}))|=-1
\]
and
\[
  |E(\Ext(C-q,S\cup\{p\}))|-|E(\Ext(C-q,S))|=+1.
\]
All other aligned shadows have edge-count defect \(0\).  Therefore any deck
matching must compensate for the missing root edge at \(p\) by rematching
cards; it cannot simply align every old vertex with itself.

\begin{lemma}[Single-switch degree balance]
If \(\PDeck_C(S\cup\{p\})=\PDeck_C(S\cup\{q\})\), then
\[
  \deg_C(p)=\deg_C(q).
\]
\end{lemma}

\begin{proof}[Proof sketch]
The full extensions have the same vertex deck, so they have the same degree
multiset.  The root has the same degree on both sides, namely \(|S|+1\).
Every old vertex except \(p\) and \(q\) has the same full-extension degree on
both sides.  At \(p\) and \(q\), the two multisets are
\[
  \{\!\{\, \deg_C(p)+1,\deg_C(q) \,\}\!\}
  \quad\text{and}\quad
  \{\!\{\, \deg_C(p),\deg_C(q)+1 \,\}\!\}.
\]
Equality of these two multisets forces \(\deg_C(p)=\deg_C(q)\).
\end{proof}

\begin{lemma}[Single-switch \(S\)-degree balance]
If \(\PDeck_C(S\cup\{p\})=\PDeck_C(S\cup\{q\})\), then
\[
  \deg_S(p)=\deg_S(q).
\]
\end{lemma}

\begin{proof}[Proof sketch]
The full decks of
\[
  \Ext(C,S\cup\{p\})
  \quad\text{and}\quad
  \Ext(C,S\cup\{q\})
\]
are equal, because both are obtained from the common root-deleted card \(C\)
plus the equal punctured extension decks.  Therefore they have the same
number of triangles, since triangle counts are reconstructible from the
vertex deck by Kelly's lemma.

Triangles avoiding the root are just triangles of \(C\), hence cancel.  The
root-using triangles are exactly the edges induced by the attachment set.  On
the \(A\)-side this number is
\[
  |E(C[S])|+\deg_S(p),
\]
and on the \(B\)-side it is
\[
  |E(C[S])|+\deg_S(q).
\]
Equality of triangle counts gives \(\deg_S(p)=\deg_S(q)\).
\end{proof}

\begin{lemma}[Root-link profile balance]
In the same single-switch setting, let \(F\) be any finite graph with fewer
vertices than \(\Ext(C,S\cup\{p\})\).  Then the number of induced copies of
\(F\) in \(\Ext(C,S\cup\{p\})\) that use the added root is equal to the
number of induced copies of \(F\) in \(\Ext(C,S\cup\{q\})\) that use the
added root.
\end{lemma}

\begin{proof}[Proof sketch]
The full decks of the two extensions are equal, so Kelly's lemma gives equal
total induced \(F\)-counts.  Copies of \(F\) that avoid the added root lie
entirely in the common card \(C\), so their counts are identical on the two
sides.  Subtracting these common root-avoiding counts gives equality of the
root-using counts.
\end{proof}

This is the real source of the low-order balances.  Taking \(F=K_3\) gives
\(\deg_S(p)=\deg_S(q)\).  Larger choices of \(F\) say that the way \(p\)
interacts with \(S\), and more generally with all of \(C\), has the same
root-visible profile as the way \(q\) interacts with \(S\).  A non-twin
single-switch state must therefore hide its witnesses from every proper
root-link count; the only remaining distinction is how those witnesses are
arranged globally in the card matching.

Now suppose the \(p\)-shadow on the \(A\)-side is matched root-preservingly to
the \(v\)-shadow on the \(B\)-side:
\[
  \Ext(C-p,S)
  \cong
  \Ext(C-v,B\setminus\{v\}).
\]
Root degree preservation forces
\[
  |S|=|B|-\one_{v\in B}=|S|+1-\one_{v\in B},
\]
so \(v\in B\).  Thus the \(A\)-side deleted defect vertex \(p\) can only be
matched to a \(B\)-colored vertex \(v\in S\cup\{q\}\).  Symmetrically, the
\(B\)-side deleted defect vertex \(q\) can only be matched to an
\(A\)-colored vertex.

This is a useful first step: in defect \(2\), the trade must move deleted
defect vertices into the colored set on the other side.  If one of these
matches has zero star error, we are done.  If not, parity says its star error
is at least \(2\), which is exactly the original defect size.  Hence a
minimal single-switch counterexample would have the following rigid form:
every defect-rooted root-preserving card match has star error exactly \(2\) or
larger, and no zero-error match exists.

\begin{conjecture}[No exact two-error single switch]
The rigid form just described cannot occur.  More concretely, if
\[
  A=S\cup\{p\},
  \qquad
  B=S\cup\{q\},
\]
and \(\PDeck_C(A)=\PDeck_C(B)\), then some root-preserving card match based at
\(p\) or \(q\) has zero star error.
\end{conjecture}

\begin{conjecture}[Direct single-switch card match]
In the same single-switch situation, deck equality already forces a
root-preserving isomorphism
\[
  \Ext(C-p,S)\cong \Ext(C-q,S)
\]
whose recentered star error is zero.  Equivalently, there is an isomorphism
\[
  \psi:C-p\cong C-q
\]
with \(\psi(S)=S\) that extends by \(p\mapsto q\) to an automorphism of \(C\).
\end{conjecture}

This is stronger than merely asking for some zero-error defect-rooted match.
It says that the two natural defect-deleted shadows already match each other.
Computationally this stronger form holds in the tested range.  If true, it
proves the single-switch base case in one shot.

\subsection{Detours in the single-switch trade cycle}

The direct card-match conjecture isolates the only way the base case can
fail.  The natural \(p\)-deleted shadow
\[
  \Ext(C-p,S)
\]
has root degree \(|S|\).  On the \(B\)-side, a root-preserving match with the
same root degree must delete a vertex \(v\in B=S\cup\{q\}\).  If \(v=q\), we
are in the direct case.  Thus a failure of the direct conjecture means that
every trade sends \(p\) first to some vertex of \(S\), creating a detour
around the unique defect cycle before it reaches \(q\).

Choose a root-preserving trade and let
\[
  p=x_0,\ x_1,\ldots,x_k=q
\]
be the directed arc of the \(\tau\)-cycle from \(p\) to \(q\), where
\(\tau(x_i)=x_{i+1}\).  In the non-direct case, \(k\geq 2\), and the internal
vertices \(x_1,\ldots,x_{k-1}\) lie in \(S\).  Along this arc, the rooted
cards give isomorphisms
\[
  \Ext(C-x_i,A\setminus\{x_i\})
  \cong
  \Ext(C-x_{i+1},B\setminus\{x_{i+1}\}).
\]
For \(1\leq i<k\), both \(x_i\) and \(x_{i+1}\) lie in \(S\), so the color
sets differ by moving the hole along \(S\) while keeping the switched vertex
on opposite sides.  The final step into \(q\) is the only step that can repair
the original missing color.

\begin{conjecture}[Single-switch detour shortening]
In a minimal single-switch counterexample, one can choose a root-preserving
trade with no detour arc
\[
  p=x_0\to x_1\to\cdots\to x_k=q
\]
with \(k\geq 2\).  Equivalently, one can choose a trade with \(\tau(p)=q\).
\end{conjecture}

This is a plausible intermediate target between raw deck equality and the
direct card-match conjecture.  The rooted invariant cycle-balance laws apply
on the whole \(p\)-to-\(q\) cycle, while the internal detour vertices have
zero edge-count charge.  If an internal detour exists, one hopes to splice the
cycle or recenter at an internal vertex to obtain either a shorter detour or a
zero-error direct match.

Detouring trades are not themselves pathological.  They occur frequently in
small terminal examples, sometimes even with zero star error.  The point is
selection, not exclusion: the deck may admit many trades, but the conjecture
asks for at least one direct representative.

In the rooted card-match graph, this says that the edge \(p\leadsto q\) should
exist.  If a particular perfect matching uses a detour edge \(p\leadsto s\)
with \(s\in S\), that is not a contradiction; it only means the matching has
chosen a non-direct representative.  The selection problem is to show that
the match graph contains a direct good edge, and preferably that Hall
alternation can modify any detouring perfect matching to one using it.

\begin{lemma}[Detour arcs are degree-flat]
Let
\[
  p=x_0\to x_1\to\cdots\to x_k=q
\]
be the \(p\)-to-\(q\) arc of a root-preserving single-switch trade, with
internal vertices in \(S\).  Then
\[
  \deg_C(x_0)=\deg_C(x_1)=\cdots=\deg_C(x_k).
\]
\end{lemma}

\begin{proof}[Proof sketch]
For any trade edge \(x_i\leadsto x_{i+1}\), equality of the edge counts of
the matched rooted cards gives
\[
  \deg_C(x_i)+\one_{x_i\in A}
  =
  \deg_C(x_{i+1})+\one_{x_{i+1}\in B}.
\]
Along the \(p\)-to-\(q\) arc, \(x_i\in A\) and \(x_{i+1}\in B\) at every
step: \(p\in A\), \(q\in B\), and every internal detour vertex lies in
\(S\subseteq A\cap B\).  The indicator terms therefore cancel at each step.
\end{proof}

More generally, every rooted invariant gives a transport equation along the
detour:
\[
  I_A(x_i)=I_B(x_{i+1}).
\]
For \(I=T\) and \(I=P\), these equations relate the rooted triangle and
two-step profiles of consecutive detour vertices.  A detour is therefore a
string of vertices that are not merely degree-equal; they are successively
indistinguishable by the rooted card profiles determined by the switch.

\begin{lemma}[Natural shadow isolation criterion]
In the single-switch setting, suppose the rooted card
\[
  \Ext(C-p,S)
\]
is not isomorphic to any card of the form
\[
  \Ext(C-s,(S\setminus\{s\})\cup\{q\})
  \qquad (s\in S).
\]
Then
\[
  \Ext(C-p,S)\cong \Ext(C-q,S).
\]
\end{lemma}

\begin{proof}[Proof sketch]
The card \(\Ext(C-p,S)\) appears on the \(A\)-side punctured extension deck.
By equality of punctured decks it must appear on the \(B\)-side.  Root degree
shows that any root-preserving occurrence on the \(B\)-side must delete a
vertex of \(B=S\cup\{q\}\).  By hypothesis it cannot be obtained by deleting
a vertex of \(S\), so it must be the \(q\)-deleted card.
\end{proof}

The difficult case is therefore exactly the non-isolated case: the natural
\(p\)-shadow is profile-equivalent to at least one detour shadow.  The
degree-flat lemma is the first invariant obstruction to such absorption; the
rooted triangle and rooted two-step transport equations provide stronger
filters.

\subsection{Shadow-type flow}

There is a useful multiplicity refinement of the isolation criterion.  Let
\[
  T_p=\Ext(C-p,S)
\]
be the natural \(p\)-deleted rooted card type.  Define
\[
  L(T_p)=\{u\in V(C):\Ext(C-u,A\setminus\{u\})\cong T_p\},
\]
and
\[
  R(T_p)=\{v\in V(C):\Ext(C-v,B\setminus\{v\})\cong T_p\}.
\]
Equality of punctured decks gives
\[
  |L(T_p)|=|R(T_p)|.
\]
The vertex \(p\) always lies in \(L(T_p)\).  The vertex \(q\) lies in
\(R(T_p)\) exactly when the direct edge \(p\leadsto q\) exists.  The vertices
of \(R(T_p)\cap S\) are precisely the \(p\)-absorbers.

Thus if the direct edge is absent, the \(p\)-mass in the deck must be absorbed
inside \(S\):
\[
  q\notin R(T_p)
  \quad\Longrightarrow\quad
  R(T_p)\subseteq S.
\]
Multiplicity balance then forces
\[
  |L(T_p)\setminus\{p\}|=|R(T_p)|.
\]
So every failed direct edge produces not only absorbers on the \(B\)-side, but
also additional suppliers on the \(A\)-side: vertices
\[
  u\neq p
  \quad\text{with}\quad
  \Ext(C-u,A\setminus\{u\})\cong\Ext(C-p,S).
\]

\begin{conjecture}[No closed shadow-type flow]
In the single-switch setting, the natural shadow type \(T_p\) cannot be
supported entirely by a closed supplier--absorber flow avoiding \(q\).  In
other words, multiplicity balance for \(T_p\), together with the absorber
constraints for every vertex of \(R(T_p)\cap S\), forces \(q\in R(T_p)\).
\end{conjecture}

This is a sharper form of detour selection.  Instead of following one chosen
perfect matching, it examines the whole fiber of the natural card type.  A
counterexample would need a balanced population of \(T_p\)-suppliers and
\(T_p\)-absorbers, all satisfying the same rooted profile constraints, with
the natural destination \(q\) missing from the right fiber.

There is also a dual flow.  Let
\[
  T_q=\Ext(C-q,S)
\]
be the natural \(q\)-deleted rooted card type on the \(B\)-side.  If the
direct edge \(p\leadsto q\) is absent, then \(T_p\neq T_q\).  The type \(T_q\)
appears at \(q\) on the \(B\)-side, and equality of punctured decks forces it
to appear on the \(A\)-side.  Root degree again shows that every such
\(A\)-side occurrence must delete a vertex of \(A=S\cup\{p\}\).  Since
deleting \(p\) would give \(T_p\), not \(T_q\), the \(T_q\)-mass must also be
routed through \(S\).

\begin{conjecture}[No coupled closed shadow flows]
In a single-switch same-deck state, it is impossible for both natural shadow
types
\[
  T_p=\Ext(C-p,S),
  \qquad
  T_q=\Ext(C-q,S)
\]
to be routed entirely through the common set \(S\), avoiding the direct
positions \(q\) and \(p\), respectively.
\end{conjecture}

This coupled-flow conjecture is another form of direct-edge selection.  It
says that the common set \(S\) cannot simultaneously absorb all evidence of
the \(p\)-deleted natural card and all evidence of the \(q\)-deleted natural
card.  The two flows would have to pass through vertices with the same total
degree, the same \(S\)-degree constraints, and compatible rooted profiles.
The expectation is that these constraints force the two natural shadow types
to coincide after all, producing the direct edge.

\begin{definition}[Supplier]
In the same setting, a vertex \(u\in S\) is a \(p\)-supplier if
\[
  \Ext(C-u,(S\setminus\{u\})\cup\{p\})
  \cong
  \Ext(C-p,S).
\]
Equivalently, \(u\in L(T_p)\setminus\{p\}\).
\end{definition}

\begin{lemma}[Supplier constraints]
If \(u\in S\) is a \(p\)-supplier, then
\[
  \deg_C(u)=\deg_C(p)
\]
and
\[
  \deg_S(u)=\deg_{S\setminus\{u\}}(p).
\]
\end{lemma}

\begin{proof}[Proof sketch]
The root degree of the \(u\)-deleted \(A\)-side card is
\[
  |A|-\one_{u\in A}.
\]
To match \(T_p=\Ext(C-p,S)\), whose root degree is \(|S|\), we must have
\(u\in A\).  Since \(u\neq p\), this puts \(u\in S\).  Equality of total edge
counts then gives \(\deg_C(u)=\deg_C(p)\).

For the second equality, compare root-using triangle counts.  The natural
\(p\)-shadow has \(|E(C[S])|\) such triangles.  The \(u\)-deleted
\(A\)-shadow has color set \((S\setminus\{u\})\cup\{p\}\), so its root-using
triangle count is
\[
  |E(C[S])|-\deg_S(u)+\deg_{S\setminus\{u\}}(p).
\]
Equality gives the displayed constraint.
\end{proof}

\begin{definition}[Absorber]
In the single-switch setting, a vertex \(s\in S\) is a \(p\)-absorber if
there is a root-preserving isomorphism
\[
  \Ext(C-p,S)
  \cong
  \Ext(C-s,(S\setminus\{s\})\cup\{q\}).
\]
Equivalently, the natural \(p\)-deleted shadow can be matched to the
\(s\)-deleted shadow on the \(B\)-side.
\end{definition}

\begin{lemma}[Absorber constraints]
Suppose \(s\in S\) is a \(p\)-absorber.  Then the following numerical
constraints hold:
\[
  \deg_C(s)=\deg_C(p),
\]
\[
  \deg_S(s)=\deg_{S\setminus\{s\}}(q),
\]
and, assuming the single-switch degree balance
\(\deg_C(p)=\deg_C(q)\),
\[
  \deg_S(p)=\deg_S(q).
\]
\end{lemma}

\begin{proof}[Proof sketch]
The first equality comes from equality of the total edge counts of the
matched rooted cards:
\[
  |E(C-p)|+|S|
  =
  |E(C-s)|+|S|.
\]
For the second equality, compare root-using triangle counts.  The \(p\)-shadow
has
\[
  |E(C[S])|
\]
root-using triangles, while the \(s\)-deleted \(B\)-shadow has
\[
  |E(C[(S\setminus\{s\})\cup\{q\}])|
  =
  |E(C[S])|-\deg_S(s)+\deg_{S\setminus\{s\}}(q).
\]
Equality gives the displayed triangle constraint.

For the last equality, compare the rooted two-step counts \(P\).  The
\(p\)-shadow gives
\[
  \sum_{x\in S}\deg_{C-p}(x)
  =
  \sum_{x\in S}\deg_C(x)-\deg_S(p).
\]
The \(s\)-deleted \(B\)-shadow gives
\[
  \sum_{x\in S\setminus\{s\}}\deg_{C-s}(x)+\deg_{C-s}(q)
  =
  \sum_{x\in S}\deg_C(x)-\deg_C(s)-\deg_S(s)+\deg_C(q)-\one_{sq\in E(C)}.
\]
Using \(\deg_C(s)=\deg_C(p)=\deg_C(q)\), this becomes
\[
  \sum_{x\in S}\deg_C(x)-\deg_S(s)-\one_{sq\in E(C)}.
\]
Together with
\[
  \deg_S(s)=\deg_{S\setminus\{s\}}(q)
  =
  \deg_S(q)-\one_{sq\in E(C)},
\]
the two-step equality gives \(\deg_S(p)=\deg_S(q)\).
\end{proof}

Thus absorption is already quite restrictive.  A \(p\)-absorber must have the
same total degree as \(p\), must see \(S\) exactly as \(q\) sees
\(S\setminus\{s\}\), and is compatible with the global \(S\)-degree balance
forced by triangle reconstruction.  If these absorber constraints fail for
every \(s\in S\), the natural \(p\)-shadow is isolated and the direct
\(p\leadsto q\) edge is forced.

Therefore the hard detour-selection case is not just degree-balanced in
\(C\); it is also balanced relative to the common color:
\[
  \deg_C(p)=\deg_C(q),
  \qquad
  \deg_S(p)=\deg_S(q).
\]
Equivalently, \(p\) and \(q\) have the same number of neighbors in
\(V(C)\setminus(S\cup\{p,q\})\) as well.  The remaining witnesses must be
balanced both inside and outside the common color.

\subsection{Absorber chains}

Absorbers should be organized into paths, not treated as isolated accidents.
Define the absorption digraph for the single-switch state as follows.  Its
vertices are
\[
  S\cup\{q\}.
\]
There is an edge
\[
  p\to y
\]
if the natural \(p\)-deleted shadow can be matched root-preservingly to the
\(y\)-deleted \(B\)-side shadow.  Thus \(y=q\) is the desired direct edge,
and \(y\in S\) is a \(p\)-absorber.  More generally, along a chosen
root-preserving trade cycle
\[
  p=x_0\to x_1\to\cdots\to x_k=q,
\]
each step \(x_i\to x_{i+1}\) is an absorption edge for the current deleted
shadow.

For any such step, equality of rooted card profiles gives transport equations
\[
  I_A(x_i)=I_B(x_{i+1})
\]
for every rooted invariant \(I\).  In particular, the whole chain is
degree-flat:
\[
  \deg_C(x_0)=\deg_C(x_1)=\cdots=\deg_C(x_k),
\]
and it is also flat for every invariant that can be read from the rooted card
after correcting for the moving color set.

\begin{conjecture}[Absorber-chain shortening]
If there is an absorption chain
\[
  p=x_0\to x_1\to\cdots\to x_k=q
\]
with \(k\geq 2\), then either there is a direct edge \(p\to q\), or there is a
shorter absorption chain from \(p\) to \(q\).
\end{conjecture}

This is the Hall-alternation version of detour shortening.  The hope is to
splice a chain at an internal absorber using the fact that all vertices on the
chain have compatible rooted profiles.  The absorber constraints above are
the first low-order consequences of this compatibility; higher rooted
subgraph counts should make long chains increasingly rigid.

\subsection{Choosing the extendable card isomorphism}

Even after detour shortening, there is one remaining issue.  A direct
root-preserving match gives
\[
  \psi:C-p\cong C-q,\qquad \psi(S)=S.
\]
This proves that \(p\) and \(q\) are deletion-similar relative to the color
set \(S\), but it does not by itself prove that \(\psi\) extends over
\(p\mapsto q\).  Extension is exactly the zero-star-error condition
\[
  \psi\bigl(N_C(p)\setminus\{p\}\bigr)=N_C(q)\setminus\{q\}.
\]

Thus the direct single-switch theorem has two layers:
\[
  \text{find a direct card match } C-p\cong C-q,
\]
then
\[
  \text{choose such a match that also carries }N_C(p)\text{ to }N_C(q).
\]
The first layer is detour shortening.  The second layer is neighborhood
selection inside the isomorphism set
\[
  \Iso_S(C-p,C-q).
\]

\begin{conjecture}[Neighborhood selection]
Assume \(\PDeck_C(S\cup\{p\})=\PDeck_C(S\cup\{q\})\), and assume
\(\Iso_S(C-p,C-q)\) is nonempty.  Then there exists
\[
  \psi\in\Iso_S(C-p,C-q)
\]
such that
\[
  \psi(N_C(p)\setminus\{p\})=N_C(q)\setminus\{q\}.
\]
\end{conjecture}

This is a colored version of the familiar obstruction that two vertices may
have isomorphic deleted cards without being related by an automorphism.  The
extra hypothesis here is not merely \(C-p\cong C-q\); it is equality of the
entire punctured extension decks for the two one-vertex colorings.  The
conjecture says that this extra deck information selects an extendable
isomorphism of the common deleted card.

\subsection{Stabilizer-orbit formulation}

Neighborhood selection can be made completely algebraic.  Fix one direct
card isomorphism
\[
  \psi_0:C-p\cong C-q,\qquad \psi_0(S)=S.
\]
Let
\[
  K=C-p,
  \qquad
  U=N_C(p)\setminus\{p\},
  \qquad
  V_0=\psi_0^{-1}(N_C(q)\setminus\{q\}).
\]
Let
\[
  \Aut(K,S)=\{\alpha\in\Aut(K):\alpha(S)=S\}.
\]

\begin{lemma}[Neighborhood selection is an orbit problem]
There exists an \(S\)-preserving isomorphism
\[
  \psi:C-p\cong C-q
\]
that extends over \(p\mapsto q\) if and only if \(U\) and \(V_0\) are in the
same \(\Aut(K,S)\)-orbit.
\end{lemma}

\begin{proof}[Proof sketch]
Every other \(S\)-preserving isomorphism \(C-p\cong C-q\) has the form
\[
  \psi=\psi_0\alpha
\]
with \(\alpha\in\Aut(K,S)\).  The transported \(q\)-neighborhood is then
\[
  \psi^{-1}(N_C(q)\setminus\{q\})
  =
  \alpha^{-1}(V_0).
\]
Thus \(\psi\) has zero star error exactly when
\[
  U=\alpha^{-1}(V_0),
\]
equivalently when \(U\) and \(V_0\) lie in the same orbit.
\end{proof}

So a failure of neighborhood selection is not vague.  It is a pair of subsets
\[
  U,V_0\subseteq V(K)
\]
that are indistinguishable enough to arise from the same full
punctured-extension deck, but still lie in distinct \(\Aut(K,S)\)-orbits.
This is the point where group action can replace some of the deck language.

\begin{lemma}[Finite orbit separator]
If \(U\) and \(V_0\) are not in the same \(\Aut(K,S)\)-orbit, then there is a
finite two-colored rooted pattern \(F\) whose number of induced occurrences
in the colored structure
\[
  (K,S,U)
\]
differs from its number of induced occurrences in
\[
  (K,S,V_0).
\]
\end{lemma}

\begin{proof}[Proof sketch]
The two finite colored structures \((K,S,U)\) and \((K,S,V_0)\) are
non-isomorphic exactly when \(U\) and \(V_0\) are in distinct
\(\Aut(K,S)\)-orbits.  Taking \(F\) to be the whole colored structure
\((K,S,U)\) already separates them: it occurs in itself, with multiplicity
\(|\Aut(K,S,U)|\), and has no surjective color-preserving induced occurrence
in the non-isomorphic structure.  One can then choose a separator of minimal
size or complexity.
\end{proof}

The minimal separator \(F\) is a possible measure of neighborhood-selection
failure.  If \(F\) is small, it should appear in rooted triangle, path, or
other low-order shadow counts.  If \(F\) is large, then all smaller colored
patterns agree, which says \(U\) and \(V_0\) are locally indistinguishable
relative to \(S\).  Either outcome is useful: a small separator feeds cycle
balance, while a large minimal separator gives a strong local-equivalence
normal form.

\begin{lemma}[Anchored Kelly pressure]
Suppose the distinguished root information is available in the deck, so that
we can compare the one-vertex-deleted cards as rooted structures over the
same first root.  Then equality of the anchored decks of
\[
  (K,S,U)
  \quad\text{and}\quad
  (K,S,V)
\]
implies equality of the number of induced occurrences of every proper
two-colored pattern \(F\), meaning every \(F\) with fewer vertices than \(K\).
\end{lemma}

\begin{proof}[Proof sketch]
This is the colored version of Kelly's lemma.  Count pairs
\[
  (W,z)
\]
where \(W\) is an induced occurrence of \(F\) and \(z\) is a vertex outside
that occurrence.  Each occurrence of \(F\) is counted \(|V(K)|-|V(F)|\) times,
and the same pair count can be read from the anchored deck by summing the
number of \(F\)-occurrences in each one-vertex-deleted anchored card.
\end{proof}

Consequently, under anchored deck equality, a finite orbit separator for
\((K,S,U)\) and \((K,S,V)\) cannot be proper.  Any failure of neighborhood
selection must be a full-size colored hypomorphism: all proper colored
patterns agree, but the full colored structures are not isomorphic.  This is
exactly the kind of rigid boundary case where reconstruction difficulties
usually live.

The caveat is important.  The original deck is unrooted.  The anchored Kelly
pressure applies only after the root-moving ambiguity has been controlled, or
inside a strengthened anchored subproblem where the distinguished root is
part of the data.  Still, it explains why low-order invariants either quickly
prove neighborhood selection or force the bad case into a very narrow
full-size hypomorphic corner.

\subsection{Fixed-host subset reconstruction}

The neighborhood-selection problem can be isolated from graph reconstruction
as follows.  Fix a colored host
\[
  (K,S).
\]
For a subset \(U\subseteq V(K)\), define its fixed-host deletion deck
\[
  \Deck_{K,S}(U)
  =
  \{\!\{\, (K-z,\ S\setminus\{z\},\ U\setminus\{z\}) : z\in V(K) \,\}\!\},
\]
where the cards are taken up to color-preserving graph isomorphism.

\begin{conjecture}[Fixed-host subset reconstruction]
If
\[
  \Deck_{K,S}(U)=\Deck_{K,S}(V),
\]
and
\[
  |U|=|V|,
\]
then \(U\) and \(V\) are in the same \(\Aut(K,S)\)-orbit.
\end{conjecture}

This conjecture is exactly the anchored neighborhood-selection problem once
the first root has been recognized.  It is weaker-looking than graph
reconstruction because the uncolored host \(K\) and the first color \(S\) are
held fixed.  The only unknown is the second color class.  The equal-size
hypothesis is essential: for a one-vertex host, the deletion deck cannot
distinguish the empty second color from the singleton second color after the
only vertex is deleted.  In the graph-reconstruction application, equal size
is supplied by root-degree equality.

\begin{remark}
The fixed-host problem is still nontrivial.  The cards are not labelled
subsets of \(K-z\); they are isomorphism types of colored deleted hosts.  A
host with many deletion-similar vertices can hide the location of \(U\), which
is exactly the phenomenon behind absorbers and detours.  But the problem has
a finite group-action target: prove orbit equality under \(\Aut(K,S)\), not
construct an arbitrary graph isomorphism.
\end{remark}

Anchored Kelly pressure says that any counterexample to fixed-host subset
reconstruction must have all proper colored pattern counts equal.  Thus a
minimal counterexample is a fixed-host colored hypomorphism:
\[
  (K,S,U)
  \quad\text{and}\quad
  (K,S,V)
\]
have equal colored deletion decks and equal proper colored substructure
counts, but \(U\) and \(V\) are in different \(\Aut(K,S)\)-orbits.

\begin{remark}[Computational evidence]
A brute-force check found no counterexample to equal-size fixed-host subset
reconstruction for all unlabeled hosts \(K\) on at most five vertices, all
choices of \(S\), and all equal-size pairs \(U,V\).  The same check without
the equal-size hypothesis fails immediately on the one-vertex host, as noted
above.
\end{remark}

\subsection{Fixed-host trades and one-star errors}

The fixed-host obstruction has the same local shape as the centered
obstruction.  Assume
\[
  \Deck_{K,S}(U)=\Deck_{K,S}(V),
  \qquad |U|=|V|.
\]
Choose a matching of the fixed-host decks.  Thus for each \(z\in V(K)\) there
is a vertex \(\tau(z)\in V(K)\) and a color-preserving isomorphism
\[
  \eta_z:
  (K-z,\ S\setminus\{z\},\ U\setminus\{z\})
  \cong
  (K-\tau(z),\ S\setminus\{\tau(z)\},\ V\setminus\{\tau(z)\}).
\]

\begin{lemma}[Deleted colors are preserved]
For every matched pair \(z\leadsto \tau(z)\),
\[
  z\in S \Longleftrightarrow \tau(z)\in S,
  \qquad
  z\in U \Longleftrightarrow \tau(z)\in V.
\]
Consequently
\[
  \tau(S)=S,
  \qquad
  \tau(U)=V.
\]
\end{lemma}

\begin{proof}[Proof sketch]
The first color in the deleted card has size
\[
  |S|-\one_{z\in S}
\]
on the left and
\[
  |S|-\one_{\tau(z)\in S}
\]
on the right.  Since \(\eta_z\) is color-preserving, these sizes are equal,
so \(z\in S\) if and only if \(\tau(z)\in S\).  The same argument for the
second color uses \(|U|=|V|\) and gives
\[
  |U|-\one_{z\in U}
  =
  |V|-\one_{\tau(z)\in V}.
\]
The bijectivity of \(\tau\) then gives \(\tau(S)=S\) and \(\tau(U)=V\).
\end{proof}

Extend \(\eta_z\) to a bijection
\[
  \widehat\eta_z:V(K)\to V(K)
\]
by sending \(z\) to \(\tau(z)\).  The previous lemma says
\[
  \widehat\eta_z(S)=S,
  \qquad
  \widehat\eta_z(U)=V.
\]
Moreover \(\widehat\eta_z\) preserves all graph adjacencies except possibly
those incident to \(z\).  Define its fixed-host star error by
\[
  \Err^{K}_z(\widehat\eta_z)
  =
  \{x\in V(K)\setminus\{z\}:
    (zx\in E(K))
    \not\leftrightarrow
    (\tau(z)\widehat\eta_z(x)\in E(K))
  \}.
\]

\begin{lemma}[Zero fixed-host star error is orbit success]
If
\[
  \Err^{K}_z(\widehat\eta_z)=\varnothing
\]
for some fixed-host deck match, then \(U\) and \(V\) lie in the same
\(\Aut(K,S)\)-orbit.
\end{lemma}

\begin{proof}[Proof sketch]
The map \(\eta_z\) is already an isomorphism away from \(z\).  Empty star
error says it also preserves every adjacency incident to \(z\).  Hence
\(\widehat\eta_z\in\Aut(K)\).  It preserves \(S\) and carries \(U\) to \(V\),
so it is an element of \(\Aut(K,S)\) carrying \(U\) to \(V\).
\end{proof}

Thus a fixed-host counterexample supplies a concrete object: a coherent family
of color-correct one-star maps with no zero-error member.  This is much
sharper than the phrase ``full-size orbit failure'': the failure is witnessed
locally by star errors, one for each chosen deletion-card match.

\begin{lemma}[Fixed-host rehosting]
Let \(z\leadsto w\) be a fixed-host deck match with isomorphism
\[
  \eta:(K-z,S\setminus\{z\},U\setminus\{z\})
  \cong
  (K-w,S\setminus\{w\},V\setminus\{w\}).
\]
Identify both deleted colored hosts with the common colored card
\[
  (L,S_L,U_L).
\]
Then the two full colored hosts are recovered from \(L\) by adding one new
vertex with the same pair of color memberships, and with attachment sets
\[
  N_K(z)\setminus\{z\},
  \qquad
  \eta^{-1}(N_K(w)\setminus\{w\}).
\]
The attachment defect in this rehosted centered presentation is exactly
\[
  \Err^K_z(\widehat\eta).
\]
\end{lemma}

\begin{proof}[Proof sketch]
The deleted-colors lemma gives
\[
  z\in S \Longleftrightarrow w\in S,
  \qquad
  z\in U \Longleftrightarrow w\in V,
\]
so the re-added vertex has the same first-color and second-color status on
both sides.  The deleted colored cards have already been identified by
\(\eta\), so the only remaining information needed to reconstruct the full
colored host is the neighborhood of the re-added vertex in the deleted card.
The difference between these two neighborhoods is precisely the star-error
set.
\end{proof}

This is the key exchange principle for the neighborhood-selection bottleneck.
In the original centered problem, an attachment-color defect can be converted
by a matched card into a one-star graph error.  In the fixed-host problem, a
second-color orbit defect can also be converted by a matched card into a
one-star graph error, but now the deleted card carries both colors as common
data.  A full-size orbit failure would therefore have to be a closed system
where these one-star errors never vanish and never descend.

\begin{conjecture}[Fixed-host one-star descent]
Assume
\[
  \Deck_{K,S}(U)=\Deck_{K,S}(V),
  \qquad |U|=|V|,
\]
and \(U,V\) are not in the same \(\Aut(K,S)\)-orbit.  Then one can choose a
fixed-host trade and a vertex \(z\in U\triangle V\) such that the rehosted
star error satisfies
\[
  |\Err^K_z(\widehat\eta_z)| < |U\triangle V|.
\]
\end{conjecture}

For \(|U\triangle V|=2\), this conjecture asks for zero star error, hence an
orbit-success automorphism.  This is exactly the root-anchored base case
needed after an exact edge-slide: the switch has moved from the original
vertices \(p,q\) to a witness pair inside the direct recentered host, and the
fixed-host descent must show that the new singleton switch cannot be a
nonterminal sink.

\begin{conjecture}[Fixed-host singleton switch]
Let
\[
  U=T\cup\{a\},
  \qquad
  V=T\cup\{b\},
  \qquad
  a,b\notin T.
\]
If
\[
  \Deck_{K,S}(U)=\Deck_{K,S}(V),
\]
then some automorphism of \(K\) preserving \(S\) carries \(U\) to \(V\).
\end{conjecture}

This is the pure fixed-host analogue of single-switch rigidity.  It rules out
nonconjugate singleton moves inside the fixed background.  Proving this
singleton statement would settle the exact two-error endpoint of the
root-anchored descent.

\begin{remark}[Computational evidence for fixed-host singleton]
A brute-force check over all unlabeled hosts \(K\) on at most six vertices,
all choices of \(S\) and \(T\), and all ordered pairs \(a,b\notin T\), found
433,448 fixed-host singleton instances with equal colored deletion decks.  All
433,448 were orbit-successful, and in fact all 433,448 had a direct card edge
\[
  (K-a,S\setminus\{a\},T)
  \cong
  (K-b,S\setminus\{b\},T).
\]
No no-direct closed-flow candidate occurred in this range.  This is stronger
than needed for the inductive singleton package and is good evidence that
detour selection, rather than neighborhood selection, is the core local
obstruction.

A later random stress test on larger hosts sampled 60 random eight-vertex
colored hosts and 25 random nine-vertex colored hosts, grouping all
second-color subsets by their fixed-host deletion deck and extracting the
singleton pairs.  It found 1,024 equal-deck singleton pairs on eight vertices
and 640 on nine vertices, again with no nondirect natural-shadow match.  This
is only evidence, but it supports the stronger direct-shadow target rather
than merely the orbit conclusion.
\end{remark}

\subsection{Attacking the fixed-host singleton switch}

Write
\[
  U=T\cup\{a\},
  \qquad
  V=T\cup\{b\},
  \qquad
  a,b\notin T.
\]
The natural \(a\)-deleted fixed-host card is
\[
  H_a=(K-a,\ S\setminus\{a\},\ T),
\]
and the natural \(b\)-deleted card on the other side is
\[
  H_b=(K-b,\ S\setminus\{b\},\ T).
\]
If \(H_a\) appears on the \(V\)-side by deleting \(b\), we have a direct
fixed-host card edge \(a\leadsto b\).  If it appears by deleting a vertex
\(t\in T\), then the singleton switch has detoured through the common color:
\[
  (K-a,\ S\setminus\{a\},\ T)
  \cong
  (K-t,\ S\setminus\{t\},\ (T\setminus\{t\})\cup\{b\}).
\]

\begin{lemma}[Fixed-host natural shadow isolation]
In the fixed-host singleton switch, suppose
\[
  (K-a,\ S\setminus\{a\},\ T)
\]
is not isomorphic to any detour card
\[
  (K-t,\ S\setminus\{t\},\ (T\setminus\{t\})\cup\{b\})
  \qquad (t\in T).
\]
Then there is a direct fixed-host card edge
\[
  (K-a,\ S\setminus\{a\},\ T)
  \cong
  (K-b,\ S\setminus\{b\},\ T).
\]
\end{lemma}

\begin{proof}[Proof sketch]
The card \(H_a\) occurs in the \(U\)-deck.  By equality of fixed-host decks it
must occur in the \(V\)-deck.  The deleted-colors lemma says that, on the
\(V\)-side, a card matching \(H_a\) must delete a vertex of \(V=T\cup\{b\}\).
By hypothesis it cannot delete a vertex of \(T\), so it must delete \(b\).
\end{proof}

Thus the singleton switch again splits into detour selection and neighborhood
selection.  A direct edge gives an isomorphism
\[
  \eta:(K-a,S\setminus\{a\},T)
  \cong
  (K-b,S\setminus\{b\},T).
\]
Extending \(\eta\) by \(a\mapsto b\) gives a bijection preserving \(S\) and
preserving \(T\cup\{a\}\) as \(T\cup\{b\}\).  If its star error is empty, this
is the desired automorphism.  If its star error has size \(2\), it is again a
single edge-slide at the moved vertex.  So the exact two-error obstruction in
the original graph has reproduced itself as a fixed-host singleton switch.

There is also a profile-pressure formulation.  For a two-colored pattern
\(P\), let
\[
  m_P(x)=\pi_P(T\cup\{x\})-\pi_P(T)
\]
inside the fixed colored host \((K,S)\).  This counts the contribution to
\(P\)-occurrences caused by changing the second-color status of \(x\), with
the common background \((K,S,T)\) subtracted away.  It is best thought of as a
finite-difference invariant: for some patterns it is a direct contribution,
and for others it is a signed net contribution because occurrences containing
\(x\) may change type when \(x\) becomes second-colored.

\begin{lemma}[Moving-vertex profile balance]
Under fixed-host singleton deck equality,
\[
  m_P(a)=m_P(b)
\]
for every two-colored pattern \(P\) with fewer vertices than \(K\).
\end{lemma}

\begin{proof}[Proof sketch]
Proper pattern sums are deck-controlled, so
\[
  \pi_P(T\cup\{a\})=\pi_P(T\cup\{b\})
\]
for every proper pattern \(P\).  The term \(\pi_P(T)\) depends only on the
fixed background \((K,S,T)\), so subtracting it from both sides gives the
displayed equality.
\end{proof}

Write
\[
  \deg_R(x)=|N_K(x)\cap R|
\]
for \(R\subseteq V(K)\).  The first consequences of moving-vertex profile
balance are already strong.

\begin{lemma}[Low-order singleton balances]
Under fixed-host singleton deck equality,
\[
  a\in S \Longleftrightarrow b\in S,
\]
and
\[
  \deg_T(a)=\deg_T(b),
  \qquad
  \deg_{K\setminus T}(a)=\deg_{K\setminus T}(b),
  \qquad
  \deg_K(a)=\deg_K(b).
\]
More generally, for each of the four background classes
\[
  S\cap T,\quad T\setminus S,\quad S\setminus T,\quad
  K\setminus(S\cup T),
\]
the number of neighbors of \(a\) in that class equals the number of neighbors
of \(b\) in that class.
\end{lemma}

\begin{proof}[Proof sketch]
The first-color membership of the moving vertex is detected by the
one-vertex pattern whose unique vertex is second-colored and first-colored:
its finite difference is
\[
  \one_{x\in S}.
\]

For the degree statements, use two-vertex edge patterns and read them as
finite differences.  The pattern in which both vertices are second-colored
has finite difference
\[
  \deg_T(x).
\]
The pattern in which exactly one vertex is second-colored has finite
difference
\[
  \deg_{K\setminus T}(x)-\deg_T(x),
\]
because edges from \(x\) to \(T\) were counted before \(x\) was
second-colored and cease to have exactly one second-colored endpoint after
\(x\) is added.  Since the previous pattern already gives
\(\deg_T(a)=\deg_T(b)\), equality of the exactly-one-second-colored finite
difference gives
\[
  \deg_{K\setminus T}(a)=\deg_{K\setminus T}(b).
\]
Adding the two equalities gives \(\deg_K(a)=\deg_K(b)\).

Refining the same two-vertex edge patterns by first-color status gives a
triangular system whose entries are the neighbor counts in
\[
  S\cap T,\quad T\setminus S,\quad S\setminus T,\quad
  K\setminus(S\cup T).
\]
The entries involving \(T\) are read from the both-second-colored patterns,
and the entries outside \(T\) are then read from the exactly-one-second
patterns after subtracting the known \(T\)-entries.  All of these are proper
pattern counts unless the host has at most two vertices, in which case the
singleton statement is checked directly from the one-vertex cards.
\end{proof}

Now suppose the natural \(a\)-shadow is absorbed by a detour vertex
\(t\in T\):
\[
  (K-a,\ S\setminus\{a\},\ T)
  \cong
  (K-t,\ S\setminus\{t\},\ (T\setminus\{t\})\cup\{b\}).
\]

\begin{lemma}[Fixed-host absorber constraints]
Any such absorber \(t\in T\) satisfies
\[
  t\in S \Longleftrightarrow a\in S \Longleftrightarrow b\in S,
\]
and
\[
  \deg_K(t)=\deg_K(a)=\deg_K(b).
\]
Moreover, the second-color induced edge count gives
\[
  \deg_T(t)=\deg_T(b)-\one_{tb\in E(K)}.
\]
If the common first-color status is \(1\), then the first-color induced edge
count also gives
\[
  \deg_S(t)=\deg_S(a).
\]
\end{lemma}

\begin{proof}[Proof sketch]
The deleted-colors lemma applied to the match \(a\leadsto t\) gives
\[
  a\in S \Longleftrightarrow t\in S,
\]
and comparing the size of the intersection of the two color classes in the
deleted cards gives
\[
  t\in S \Longleftrightarrow b\in S.
\]
The uncolored edge counts of \(K-a\) and \(K-t\) give
\[
  \deg_K(a)=\deg_K(t),
\]
and the low-order singleton balance gives \(\deg_K(a)=\deg_K(b)\).

For the second-color edge equation, the left deleted card has second-color
edge count \(|E(K[T])|\).  The right deleted card has second color
\((T\setminus\{t\})\cup\{b\}\), so its induced edge count is
\[
  |E(K[T])|-\deg_T(t)+\deg_{T\setminus\{t\}}(b).
\]
Equating these and using
\[
  \deg_{T\setminus\{t\}}(b)=\deg_T(b)-\one_{tb\in E(K)}
\]
gives the displayed formula.  Finally, if \(a,t\in S\), the first-color
induced edge counts of \(S\setminus\{a\}\) and \(S\setminus\{t\}\) give
\(\deg_S(a)=\deg_S(t)\).
\end{proof}

Combining the absorber equation with \(\deg_T(a)=\deg_T(b)\), an absorber
must satisfy
\[
  \deg_T(t)=\deg_T(a)-\one_{tb\in E(K)}.
\]
So an absorber either has the same \(T\)-degree as \(a\), if it is not
adjacent to \(b\), or exactly one less, if it is adjacent to \(b\).  This is a
real restriction: detours can only pass through vertices of the same
first-color status and nearly the same background degree profile as the moved
vertex.

\begin{corollary}[Degree-profile isolation]
If no vertex \(t\in T\) satisfies the absorber constraints above, then the
fixed-host singleton switch has a direct card edge \(a\leadsto b\).
\end{corollary}

\begin{proof}[Proof sketch]
If the direct edge were absent, the natural shadow isolation lemma would force
the natural \(a\)-deleted card to be absorbed by some \(t\in T\).  Such a
vertex must satisfy the absorber constraints, contradicting the hypothesis.
\end{proof}

There is a dual set of constraints for absorbers of the natural \(b\)-shadow
on the other side.  Thus a non-direct singleton counterexample would require
two closed flows through \(T\): one absorbing the \(a\)-deleted shadow and one
absorbing the \(b\)-deleted shadow, both restricted to the same first-color
class and to the same narrow degree-profile layer.

\subsection{Fixed-host detour cycles}

The trade-cycle language applies verbatim to the fixed-host singleton switch.
Choose a fixed-host trade, so for each \(z\in V(K)\) there is a matched
deleted-card isomorphism
\[
  (K-z,S\setminus\{z\},U\setminus\{z\})
  \cong
  (K-\tau(z),S\setminus\{\tau(z)\},V\setminus\{\tau(z)\}).
\]
The deleted-colors lemma gives
\[
  \tau(S)=S,
  \qquad
  \tau(U)=V.
\]
In the singleton case \(U=T\cup\{a\}\) and \(V=T\cup\{b\}\), the edge-count
cycle-balance therefore puts \(a\) and \(b\) on the same \(\tau\)-cycle.
More explicitly, along the directed arc
\[
  a=x_0\to x_1\to\cdots\to x_k=b
\]
from \(a\) to \(b\), every internal vertex lies in \(T\), provided the arc is
chosen with no intermediate occurrence of \(a\) or \(b\).

\begin{lemma}[Fixed-host detour arcs are profile-flat]
Let
\[
  a=x_0\to x_1\to\cdots\to x_k=b
\]
be the \(a\)-to-\(b\) arc of a fixed-host singleton trade, with internal
vertices in \(T\).  Then every vertex on the arc has the same first-color
status and the same total degree in \(K\):
\[
  x_i\in S \Longleftrightarrow a\in S,
  \qquad
  \deg_K(x_i)=\deg_K(a)
  \qquad (0\leq i\leq k).
\]
\end{lemma}

\begin{proof}[Proof sketch]
For each trade edge \(x_i\leadsto x_{i+1}\), the deleted first-color sizes
give
\[
  x_i\in S \Longleftrightarrow x_{i+1}\in S.
\]
The uncolored edge counts of the matched deleted hosts give
\[
  |E(K)|-\deg_K(x_i)=|E(K)|-\deg_K(x_{i+1}),
\]
so \(\deg_K(x_i)=\deg_K(x_{i+1})\).  Iterating along the arc gives the claim.
\end{proof}

The absorber constraints sharpen this flatness at the first step: if
\(x_1=t\in T\), then
\[
  \deg_T(t)=\deg_T(a)-\one_{tb\in E(K)}.
\]
The dual constraints sharpen the final step into \(b\).  Thus a closed detour
cycle is a string of vertices in \(T\) that are globally degree-flat and
locally almost \(T\)-degree-flat, while also transporting the full deleted
colored-card type.  This is the object the no-closed-flow conjecture must
rule out.

\subsection{Absorber exchange maps}

A detour absorber has more structure than just its deleted vertex.  Suppose
\[
  \eta:
  (K-a,S\setminus\{a\},T)
  \cong
  (K-t,S\setminus\{t\},(T\setminus\{t\})\cup\{b\})
\]
with \(t\in T\).  Since \(\eta\) maps the second color \(T\) bijectively onto
\((T\setminus\{t\})\cup\{b\}\), there is a unique vertex
\[
  y_\eta\in T
\]
such that
\[
  \eta(y_\eta)=b.
\]
Call \(t\) the entry vertex of the absorber and \(y_\eta\) its exit vertex.

Extending \(\eta\) by \(a\mapsto t\), the second-color transport has the form
\[
  a\mapsto t,\qquad y_\eta\mapsto b,
\]
with the remaining vertices of \(T\setminus\{y_\eta\}\) mapped into
\(T\setminus\{t\}\).  Thus a detour does not simply hide \(a\) in \(T\).  It
injects \(a\) at the entry \(t\), ejects the exit \(y_\eta\) to \(b\), and
permutes the rest of the common second color.

\begin{definition}[Absorber exchange digraph]
For a fixed-host singleton state, define the \(a\)-absorber exchange digraph
on \(T\) by drawing an edge
\[
  t\longrightarrow y
\]
whenever there is an absorber isomorphism with entry \(t\) and exit \(y\).
Define the \(b\)-absorber exchange digraph dually.
\end{definition}

In this language, a no-direct closed flow is not just a set of absorbers; it
is a circulation in these exchange digraphs.  The entry vertex records where
the outside color is inserted into \(T\), and the exit vertex records where a
vertex of \(T\) leaves to become the outside color on the other side.

\begin{conjecture}[No closed absorber circulation]
In a fixed-host singleton state with equal colored deletion decks, the
absorber exchange digraphs cannot support a closed circulation avoiding the
direct edge \(a\leadsto b\).
\end{conjecture}

This is a more combinatorial version of no closed fixed-host singleton flow.
It suggests a possible proof strategy: assign a potential to vertices of
\(T\), built from their rooted profiles relative to \((K,S,T,a,b)\), and show
that every absorber exchange edge strictly changes that potential in one
direction unless a direct edge exists.  The low-order absorber equations give
the first coordinates of such a potential; the witness-charge equations give
the next signed coordinates.

\begin{lemma}[Exit constraints]
Let \(t\to y\) be an absorber exchange edge, arising from
\[
  \eta:
  (K-a,S\setminus\{a\},T)
  \cong
  (K-t,S\setminus\{t\},(T\setminus\{t\})\cup\{b\})
\]
with \(\eta(y)=b\).  Then
\[
  y\in S \Longleftrightarrow b\in S,
\]
\[
  \deg_K(y)-\one_{ay\in E(K)}
  =
  \deg_K(b)-\one_{tb\in E(K)},
\]
and
\[
  \deg_T(y)=\deg_T(b)-\one_{tb\in E(K)}.
\]
\end{lemma}

\begin{proof}[Proof sketch]
The first statement follows because \(\eta\) preserves the first color and
\(\eta(y)=b\).  The second statement compares the uncolored degrees of \(y\)
in the deleted host \(K-a\) and \(b\) in the deleted host \(K-t\).  The third
compares the degrees of \(y\) and \(b\) inside the second-colored induced
subgraphs of the two deleted colored cards.  On the left this degree is
\(\deg_T(y)\); on the right it is
\[
  \deg_{T\setminus\{t\}}(b)=\deg_T(b)-\one_{tb\in E(K)}.
\]
\end{proof}

\begin{corollary}[Closed exchanges avoid witness exits]
Suppose an absorber exchange edge \(t\to y\) lies in a closed circulation, so
that \(y\) is also usable as an absorber entry.  Then
\[
  \one_{ay\in E(K)}=\one_{by\in E(K)}.
\]
In particular, \(y\notin X^-\cup X^+\).
\end{corollary}

\begin{proof}[Proof sketch]
Since \(y\) is also an absorber entry, the entry constraints give
\[
  \deg_K(y)=\deg_K(a)=\deg_K(b)
\]
and
\[
  \deg_T(y)=\deg_T(a)-\one_{yb\in E(K)}.
\]
The exit constraints and the low-order balance
\(\deg_T(a)=\deg_T(b)\) give
\[
  \one_{ay\in E(K)}=\one_{tb\in E(K)}
  \quad\text{and}\quad
  \one_{yb\in E(K)}=\one_{tb\in E(K)}.
\]
Thus \(ay\) and \(by\) have the same adjacency status.
\end{proof}

This is encouraging.  The signed witnesses \(X^-\cup X^+\) are exactly the
vertices that the two-step charge needs in order to balance a non-twin
singleton switch, but a closed absorber circulation cannot use such vertices
as exits.  Therefore a nonterminal witness cycle would have to carry its
witness charges on vertices that are not the exchange exits, while the
exchange exits themselves are \(a/b\)-neutral.  That separation is a promising
place to force a contradiction.

\begin{conjecture}[Fixed-host exchange-chain shortening]
In a minimal no-direct fixed-host singleton state, the absorber exchange
digraph has a closed circulation.  Equivalently, if an absorber has an exit
vertex \(y\in T\) that is not usable as a further absorber entry, then one can
modify the deck matching to obtain either a direct edge \(a\leadsto b\) or a
strictly shorter detour chain.
\end{conjecture}

This is the precise point where the proof could break.  The exit constraints
are strong only once exits close up into entries.  If the deck can keep
throwing exits into vertices that are not themselves absorbers, then the
exchange-potential argument has no cycle to contradict.  Thus the fixed-host
detour campaign has two tasks: first prove exchange-chain shortening, then
rule out closed absorber circulations using the exit and witness-charge
constraints.

\subsection{Natural-shadow fibers in the fixed-host singleton}

There is also a multiplicity formulation, parallel to the earlier
shadow-type flow.  Let
\[
  H_a=(K-a,S\setminus\{a\},T)
\]
be the natural \(a\)-deleted card type.  Define
\[
  L(H_a)=
  \{z\in V(K):
    (K-z,S\setminus\{z\},(T\cup\{a\})\setminus\{z\})\cong H_a
  \}
\]
and
\[
  R(H_a)=
  \{z\in V(K):
    (K-z,S\setminus\{z\},(T\cup\{b\})\setminus\{z\})\cong H_a
  \}.
\]
Deck equality gives
\[
  |L(H_a)|=|R(H_a)|.
\]
The vertex \(a\) always lies in \(L(H_a)\).  The vertex \(b\) lies in
\(R(H_a)\) exactly when the direct natural shadow exists.

Since every card of type \(H_a\) has second-color size \(|T|\), color-size
comparison gives
\[
  L(H_a)\subseteq T\cup\{a\},
  \qquad
  R(H_a)\subseteq T\cup\{b\}.
\]
Thus, if the direct edge is absent,
\[
  R(H_a)\subseteq T
\]
and multiplicity balance becomes
\[
  |R(H_a)\cap T|
  =
  1+|L(H_a)\cap T|.
\]
So a no-direct state has strictly more \(a\)-absorbers than \(a\)-suppliers.
The dual natural shadow \(H_b\) gives the opposite-looking condition for the
\(b\)-flow.  A closed no-direct counterexample must reconcile these two
imbalances inside the same finite set \(T\).

The supplier constraints mirror the absorber constraints.  A vertex
\(u\in T\) is an \(a\)-supplier for \(H_a\) if
\[
  (K-u,S\setminus\{u\},(T\setminus\{u\})\cup\{a\})
  \cong
  (K-a,S\setminus\{a\},T).
\]
Then the same color-size and edge-count comparisons give
\[
  u\in S \Longleftrightarrow a\in S,
  \qquad
  \deg_K(u)=\deg_K(a),
\]
and
\[
  \deg_T(u)=\deg_T(a)-\one_{ua\in E(K)}.
\]
Compare this with the absorber condition
\[
  \deg_T(t)=\deg_T(a)-\one_{tb\in E(K)}.
\]
Thus a vertex \(z\in T\) that is neutral with respect to \(a\) and \(b\),
meaning
\[
  \one_{za\in E(K)}=\one_{zb\in E(K)},
\]
passes the low-order \(T\)-degree test for being a supplier exactly when it
passes the corresponding test for being an absorber.  Only witnesses inside
\(T\) can contribute to a low-order supplier--absorber imbalance.

\begin{lemma}[Low-order imbalance is witness-supported]
Ignoring higher colored-card structure, the difference between the low-order
candidate set for \(a\)-absorbers and the low-order candidate set for
\(a\)-suppliers is supported on
\[
  T\cap (X^-\cup X^+).
\]
\end{lemma}

\begin{proof}[Proof sketch]
The first-color and total-degree tests are the same for suppliers and
absorbers.  The only displayed low-order test that differs is the
\(T\)-degree equation, and that equation differs only by replacing
\(\one_{za\in E(K)}\) with \(\one_{zb\in E(K)}\).  These indicators differ
exactly on \(X^-\cup X^+\).
\end{proof}

This fits the exit constraint above: closed exchanges cannot use witnesses as
exits, but the supplier--absorber imbalance can only be generated by
witnesses.  Therefore a closed no-direct flow would have to create its
cardinality imbalance at witness entries while routing its exits through
neutral vertices.  That asymmetry is the current best candidate for a
contradiction.

\begin{conjecture}[Neutral absorber reversibility]
Let \(t\in T\) be an \(a\)-absorber entry, and suppose \(t\) is neutral with
respect to \(a,b\):
\[
  \one_{ta\in E(K)}=\one_{tb\in E(K)}.
\]
Then either \(t\) is also an \(a\)-supplier for the natural shadow \(H_a\), or
there is a shorter detour chain, or there is a direct edge \(a\leadsto b\).
\end{conjecture}

This is a very small local reversibility statement.  The absorber card
replaces \(t\) by \(b\) in the second color.  The supplier card replaces \(t\)
by \(a\).  If \(t\) sees \(a\) and \(b\) in the same way, then the only reason
absorption need not reverse to supply is a higher-order difference between
the roles of \(a\) and \(b\) in the deleted card \(K-t\).  But such a
difference should be visible in the proper moving-vertex profile unless it
can be pushed further down the detour chain.

\begin{conjecture}[No witness tail into a neutral circulation]
In a minimal no-direct fixed-host singleton state, an absorber exchange path
cannot start at a witness entry
\[
  t\in T\cap(X^-\cup X^+)
\]
and then feed into a closed circulation all of whose vertices are
\(a/b\)-neutral.
\end{conjecture}

This is the refined obstruction.  A closed neutral circulation by itself is
not enough for a contradiction, because an excess absorber could appear as a
tail leading into that circulation.  The multiplicity imbalance has to be
carried by witness entries, while the closed part of the exchange graph is
neutral and reversible.  Therefore the possible surviving shape is a
witness-to-neutral tail.

\begin{proposition}[Exchange reduction for direct shadow]
Assume fixed-host exchange-chain shortening, neutral absorber reversibility,
and no witness tail into a neutral circulation.  Then the fixed-host
direct-shadow theorem holds.
\end{proposition}

\begin{proof}[Proof sketch]
Suppose the direct shadow fails in a minimal counterexample.  By natural
shadow isolation, the \(a\)-natural shadow is absorbed in \(T\).  Repeatedly
follow absorber exchange exits.  Exchange-chain shortening says that, unless
a direct edge or a shorter detour appears, the path can be continued until it
enters a closed circulation.  Minimality excludes directness and shortening.
On the closed part, the exit constraints force every vertex to be
\(a/b\)-neutral; neutral reversibility makes the closed part contribute no
net supplier--absorber imbalance.

The direct-shadow failure requires
\[
  |R(H_a)\cap T|=1+|L(H_a)\cap T|,
\]
so some imbalance must be supplied outside the neutral reversible part.  The
low-order imbalance lemma says such imbalance is supported on witnesses.
Thus the remaining obstruction is exactly a witness entry feeding into a
neutral closed circulation, contradicting the no-witness-tail assumption.
\end{proof}

This proposition isolates the current proof burden into three very local
claims.  Exchange-chain shortening is combinatorial: exits of absorbers must
either close or shorten.  Neutral reversibility is structural: an absorber at
an \(a/b\)-neutral entry should be invertible as a supplier unless some
smaller obstruction appears.  The witness-tail lemma is the bridge between
the two: it must show that a signed witness charge cannot disappear into a
neutral reversible cycle.

\begin{conjecture}[No supplier-absorber imbalance]
In a fixed-host singleton state with equal colored deletion decks, the two
natural-shadow fiber balances cannot both avoid their direct vertices.  More
concretely, it is impossible to have
\[
  b\notin R(H_a)
  \qquad\text{and}\qquad
  a\notin L(H_b).
\]
\end{conjecture}

This is another equivalent-looking form of the direct-shadow theorem.  It may
be more attackable because it isolates a signed cardinality imbalance:
the \(a\)-flow needs one extra absorber in \(T\), while the \(b\)-flow needs
one extra absorber in the opposite direction.  Any proof by exchange
circulation should explain why these two imbalances cannot coexist.

\subsection{The singleton type-flow graph}

There is an even cleaner way to package the same multiplicity bookkeeping.
For each \(t\in T\), define the two card types
\[
  A_t=(K-t,S\setminus\{t\},(T\setminus\{t\})\cup\{a\}),
\]
\[
  B_t=(K-t,S\setminus\{t\},(T\setminus\{t\})\cup\{b\}).
\]
The second-color-size-\(|T|\) part of the deck equality says
\[
  \{\!\{H_a\}\!\}+\{\!\{A_t:t\in T\}\!\}
  =
  \{\!\{H_b\}\!\}+\{\!\{B_t:t\in T\}\!\}.
\]
Build a directed multigraph whose vertices are isomorphism classes of
two-colored cards of this size, with one directed edge
\[
  B_t\longrightarrow A_t
\]
for each \(t\in T\), and one extra artificial edge
\[
  H_b\longrightarrow H_a.
\]

\begin{lemma}[Type-flow decomposition]
The singleton type-flow graph is balanced: every card type has equal indegree
and outdegree.  Consequently every edge lies in a directed cycle.  In
particular, if \(H_a\neq H_b\), removing the artificial edge leaves a directed
path
\[
  H_a\leadsto H_b
\]
using only edges \(B_t\to A_t\) with \(t\in T\).
\end{lemma}

\begin{proof}[Proof sketch]
For a card type \(Y\), its indegree in the graph is
\[
  |\{t:A_t\cong Y\}|+\one_{H_a\cong Y},
\]
and its outdegree is
\[
  |\{t:B_t\cong Y\}|+\one_{H_b\cong Y}.
\]
These are equal exactly by the displayed multiset equality.  A finite
balanced directed multigraph decomposes into directed cycles.  The artificial
edge \(H_b\to H_a\) is in one such cycle; deleting it leaves a directed path
from \(H_a\) to \(H_b\).
\end{proof}

This path is a canonical form of a no-direct detour.  Its first edge is an
absorber of \(H_a\), because \(H_a\cong B_t\) for the first \(t\).  Its final
edge is a supplier into \(H_b\), because \(A_t\cong H_b\) for the last \(t\).
Intermediate edges are switches inside deleted \(T\)-cards:
\[
  (K-t,S\setminus\{t\},(T\setminus\{t\})\cup\{b\})
  \leadsto
  (K-t,S\setminus\{t\},(T\setminus\{t\})\cup\{a\}).
\]
Thus the no-direct problem is equivalent to ruling out nontrivial type-flow
paths from the natural \(a\)-shadow to the natural \(b\)-shadow.

\begin{conjecture}[No nontrivial singleton type-flow path]
In a fixed-host singleton state with equal colored deletion decks, the
type-flow path from \(H_a\) to \(H_b\) must have length zero; equivalently,
\[
  H_a\cong H_b.
\]
\end{conjecture}

This is just the fixed-host direct-shadow theorem in flow language, but it
has two advantages.  First, it is purely a finite directed multigraph
statement attached to the equality of the second-color-size-\(|T|\) deck
slice.  Second, it gives a canonical shortest-path parameter.  A minimal
counterexample can be chosen with the shortest possible nontrivial path
\[
  H_a=Y_0\to Y_1\to\cdots\to Y_r=H_b.
\]
Any repeated card type on this path can be removed, so the intermediate
types are distinct.  The next proof attempt should show that a shortest path
of positive length can be spliced or shortened using the absorber-exchange
constraints above.

\subsection{The complementary type-flow graph}

The deck equality also has a second slice, obtained by deleting vertices that
are not second-colored.  Put
\[
  O=V(K)\setminus(T\cup\{a,b\}).
\]
On the \(a\)-side, the vertices whose deletion leaves second-color size
\(|T|+1\) are \(O\cup\{b\}\); on the \(b\)-side they are \(O\cup\{a\}\).
Thus deck equality also gives
\[
  \{\!\{C_b^a\}\!\}+\{\!\{C_o^a:o\in O\}\!\}
  =
  \{\!\{C_a^b\}\!\}+\{\!\{C_o^b:o\in O\}\!\},
\]
where
\[
  C_z^a=(K-z,S\setminus\{z\},T\cup\{a\}),
  \qquad
  C_z^b=(K-z,S\setminus\{z\},T\cup\{b\}).
\]
As before, this defines a balanced directed type-flow graph, now with edges
\[
  C_o^b\longrightarrow C_o^a
  \qquad(o\in O)
\]
and an artificial edge
\[
  C_a^b\longrightarrow C_b^a.
\]

\begin{lemma}[Complementary type-flow decomposition]
The complementary type-flow graph is balanced.  Hence, if
\[
  C_b^a\not\cong C_a^b,
\]
then there is a directed path from \(C_b^a\) to \(C_a^b\) using only vertices
deleted from \(O\).
\end{lemma}

\begin{proof}[Proof sketch]
The proof is the same indegree-outdegree count as for the singleton
type-flow graph, applied to the second-color-size-\(|T|+1\) deck slice.
\end{proof}

The endpoint cancellation
\[
  C_b^a\cong C_a^b
\]
is also a direct recentering.  It matches the card obtained from
\((K,S,T\cup\{a\})\) by deleting the uncolored vertex \(b\) with the card
obtained from \((K,S,T\cup\{b\})\) by deleting the uncolored vertex \(a\).
The re-added vertex has the same color data on both sides: it has the same
first-color status by low-order singleton balance, and it is not
second-colored on either side.

\begin{lemma}[Complementary direct cancellation is a size drop]
Assume the colored centered extension theorem for centers with fewer than
\(|V(K)|\) vertices.  If
\[
  C_b^a\cong C_a^b,
\]
then the fixed-host singleton switch is solved.
\end{lemma}

\begin{proof}[Proof sketch]
Identify the two complementary endpoint cards with a common colored host
\(L\).  Recenter the two colored full hosts by re-adding \(b\) on the
\(a\)-side and \(a\) on the \(b\)-side.  Since the matched card has one fewer
vertex, the colored centered induction hypothesis applies to the resulting
equal punctured decks over \(L\).  Transporting the resulting isomorphism
back gives a color-preserving isomorphism
\[
  (K,S,T\cup\{a\})\cong (K,S,T\cup\{b\}).
\]
\end{proof}

The low slice routes through \(T\); the high slice routes through \(O\).  A
genuine no-direct singleton obstruction must make both slices balance.  This
is important because the witness set
\[
  X^-\cup X^+
\]
is split between \(T\) and \(O\).  The low-slice supplier--absorber imbalance
can only be generated by witnesses in \(T\), while the complementary slice has
the analogous imbalance supported on witnesses in \(O\).  Thus the two slices
should not be studied separately: a counterexample would need compatible
type-flow paths through both sides of the partition.

\begin{conjecture}[Two-slice incompatibility]
There is no fixed-host singleton counterexample in which both the low
type-flow path through \(T\) and the complementary type-flow path through
\(O\) are nontrivial and shortest.  At least one slice must admit a path
shortening, a direct cancellation, or a size-drop recentering.
\end{conjecture}

\begin{conjecture}[Two-slice direct cancellation]
In a fixed-host singleton state with equal colored deletion decks, at least
one of the two endpoint cancellations holds:
\[
  H_a\cong H_b
  \qquad\text{or}\qquad
  C_b^a\cong C_a^b.
\]
\end{conjecture}

This would be enough for the fixed-host singleton theorem by the two
size-drop lemmas.  The computational evidence actually supports the stronger
first alternative \(H_a\cong H_b\) in all cases checked, but the two-slice
statement may be easier to prove because it allows either color-size slice to
provide the direct recentering.

This is another way the proof might close.  The low slice says excess
absorber mass is supported by witnesses in \(T\); the high slice says the
complementary excess is supported by witnesses in \(O\).  Since the total
witness charge is globally balanced, these two flows should constrain each
other.  A future proof should make this precise by applying the two-step
charge \(Q\) simultaneously to both color-size slices.

\subsection{Two-slice charge balance}

Define the signed witness charge
\[
  \delta(z)=\one_{bz\in E(K)}-\one_{az\in E(K)}
  \qquad (z\notin\{a,b\}).
\]
Thus \(\delta=-1\) on \(X^-\), \(\delta=+1\) on \(X^+\), and
\(\delta=0\) elsewhere.  Degree balance gives the global cancellation
\[
  \sum_{z\notin\{a,b\}}\delta(z)=0.
\]
Equivalently,
\[
  \sum_{t\in T}\delta(t)
  =
  -\sum_{o\in O}\delta(o).
\]

The same charge is visible directly on type-flow paths.  For a low-slice card
type \(Y\) with second color of size \(|T|\), let
\[
  q(Y)=\sum_{x\in\text{second color of }Y}\deg_Y(x).
\]
This is the second-color two-step invariant of that card.

\begin{lemma}[Low path charge formula]
Let
\[
  H_a=Y_0\to Y_1\to\cdots\to Y_r=H_b
\]
be a low-slice type-flow path, where the edge \(Y_{i-1}\to Y_i\) is indexed
by \(t_i\in T\).  Then
\[
  q(Y_i)-q(Y_{i-1})=\delta(t_i),
\]
and therefore
\[
  \sum_{i=1}^r \delta(t_i)
  =
  q(H_b)-q(H_a)
  =
  \sum_{t\in T}\delta(t).
\]
\end{lemma}

\begin{proof}[Proof sketch]
The edge indexed by \(t\) compares
\[
  B_t=(K-t,S\setminus\{t\},(T\setminus\{t\})\cup\{b\})
\]
with
\[
  A_t=(K-t,S\setminus\{t\},(T\setminus\{t\})\cup\{a\}).
\]
The difference of the second-color two-step invariant between these two cards
is exactly
\[
  \deg_{K-t}(a)-\deg_{K-t}(b)
  =
  \one_{bt\in E(K)}-\one_{at\in E(K)}
  =
  \delta(t),
\]
using \(\deg_K(a)=\deg_K(b)\).  Summing along the path telescopes.  Finally,
\[
  q(H_b)-q(H_a)
  =
  \sum_{t\in T}\bigl(\deg_{K-b}(t)-\deg_{K-a}(t)\bigr)
  =
  \sum_{t\in T}\delta(t).
\]
\end{proof}

There is an analogous formula for the complementary path through \(O\), and
its total charge is
\[
  \sum_{o\in O}\delta(o)
  =
  -\sum_{t\in T}\delta(t).
\]
Thus the two type-flow paths carry opposite total two-step charge.

\begin{conjecture}[No opposite shortest charge paths]
In a fixed-host singleton counterexample, one cannot have simultaneously a
shortest positive low-slice path carrying charge
\[
  \sum_{t\in T}\delta(t)
\]
and a shortest positive complementary path carrying the opposite charge.  At
least one of the paths can be shortened or converted into a direct
cancellation.
\end{conjecture}

This conjecture is a charged version of two-slice incompatibility.  If the
charge on one slice is nonzero, then any path in that slice must pass through
signed witnesses.  The other slice must carry exactly the opposite signed
charge.  If both paths were shortest and nontrivial, the two paths would give
two incompatible ways of transporting the same global witness imbalance.  If
the slice charge is zero, then \(q\) alone cannot see the obstruction and one
must repeat the same argument with higher finite-difference pattern
invariants.

\subsection{The rectangular second boundary}

The two-slice picture has a more exact form.  The low slice deletes either
the moving vertex or a vertex of \(T\).  The complementary slice deletes
either the other moving vertex or a vertex of
\[
  O=V(K)\setminus(T\cup\{a,b\}).
\]
If we then delete one more vertex from the opposite side of the partition,
both slices see the same two-hole cards.

For \(t\in T\) and \(o\in O\), define
\[
  D_{t,o}^a
  =
  (K-\{t,o\},S\setminus\{t,o\},(T\setminus\{t\})\cup\{a\}),
\]
\[
  D_{t,o}^b
  =
  (K-\{t,o\},S\setminus\{t,o\},(T\setminus\{t\})\cup\{b\}).
\]
Then the elementary square identities are
\[
  A_t-o=C_o^a-t=D_{t,o}^a,
  \qquad
  B_t-o=C_o^b-t=D_{t,o}^b.
\]

Let \(\partial_R\) denote the fixed-host partial deletion operator that
deletes only vertices in \(R\) from the displayed representatives, and then
passes to colored isomorphism classes.  This operator is not an invariant of
an abstract card type unless \(R\) has been named in the ambient host.  That
is fine here: the following identity is a fixed-host identity before taking
isomorphism classes.

\begin{lemma}[Rectangular second-boundary identity]
In a fixed-host singleton state with equal colored deletion decks,
\[
  \partial_O([H_b]-[H_a])
  =
  \partial_T([C_a^b]-[C_b^a])
\]
and both sides are equal to the rectangular chain
\[
  \sum_{t\in T}\sum_{o\in O}
    \bigl([D_{t,o}^a]-[D_{t,o}^b]\bigr).
\]
\end{lemma}

\begin{proof}[Proof sketch]
The low slice gives
\[
  [H_b]-[H_a]=\sum_{t\in T}([A_t]-[B_t]).
\]
Deleting the vertices of \(O\) from this identity gives
\[
  \partial_O([H_b]-[H_a])
  =
  \sum_{t\in T}\sum_{o\in O}
    \bigl([A_t-o]-[B_t-o]\bigr)
  =
  \sum_{t\in T}\sum_{o\in O}
    \bigl([D_{t,o}^a]-[D_{t,o}^b]\bigr).
\]
The complementary slice gives
\[
  [C_a^b]-[C_b^a]=\sum_{o\in O}([C_o^a]-[C_o^b]).
\]
Deleting the vertices of \(T\) from this identity gives the same double sum,
using \(C_o^a-t=D_{t,o}^a\) and \(C_o^b-t=D_{t,o}^b\).
\end{proof}

This is stronger than the scalar charge balance.  The two-step charge \(q\)
is obtained by applying one numerical functional to the rectangular chain.
Higher finite-difference pattern invariants are obtained by applying other
colored subgraph-count functionals to the same chain.  Thus the right object
is not a single charge but the whole rectangular boundary on \(T\times O\).

\begin{conjecture}[No nonzero rectangular kernel]
Let \(X_a=(K,S,T\cup\{a\})\) and \(X_b=(K,S,T\cup\{b\})\) be a fixed-host
singleton pair with equal colored deletion decks.  If both endpoint
differences
\[
  [H_b]-[H_a],
  \qquad
  [C_a^b]-[C_b^a]
\]
are nonzero, then their rectangular second boundary cannot agree as in the
previous lemma.  Equivalently, the rectangular identity forces at least one
of
\[
  H_a\cong H_b,
  \qquad
  C_b^a\cong C_a^b.
\]
\end{conjecture}

\begin{conjecture}[Double direct cancellation]
In the same fixed-host singleton situation, equal colored deletion decks force
both endpoint cancellations:
\[
  H_a\cong H_b,
  \qquad
  C_b^a\cong C_a^b.
\]
\end{conjecture}

This is stronger than the no-kernel conjecture, and it was supported through
six host vertices.  It is false in general, however; see the cyclic
nine-vertex counterexample below.  The right target is therefore orbit
reconstruction rather than direct endpoint cancellation.

\begin{conjecture}[One-slice orbit reconstruction]
The low color-size slice alone reconstructs the fixed-host orbit.  Namely, if
\[
  \{\!\{H_a\}\!\}+\{\!\{A_t:t\in T\}\!\}
  =
  \{\!\{H_b\}\!\}+\{\!\{B_t:t\in T\}\!\},
\]
then there is an automorphism of the fixed first-colored host \((K,S)\)
carrying
\[
  T\cup\{a\}
  \quad\text{to}\quad
  T\cup\{b\}.
\]
Dually, the complementary color-size slice should imply the same orbit
conclusion.
\end{conjecture}

\begin{remark}[Direct endpoint isolation is false]
The stronger one-slice endpoint statement fails on a connected \(9\)-vertex
fixed host.  Label the vertices by
\[
  (i,j)\in \{0,1,2\}\times\{0,1,2\},
\]
and let the host be invariant under the cyclic shift \(j\mapsto j+1\).  One
example has edge set
\[
\begin{gathered}
  05,06,08,13,16,17,24,27,28,\\
  37,48,56,67,68,78,
\end{gathered}
\]
where \(ij\) denotes the edge between the integer-labelled vertices
\[
  0=(0,0),\ 1=(0,1),\ 2=(0,2),\
  3=(1,0),\ldots,\ 8=(2,2).
\]
Take \(S=\emptyset\), \(T=\{1\}\), \(a=0\), and \(b=2\).  The cyclic
automorphism sends \(T\cup\{a\}=\{0,1\}\) to
\(T\cup\{b\}=\{1,2\}\), so the full fixed-host decks and the low slices agree.
Nevertheless
\[
  (K-a,S-a,T)\not\cong (K-b,S-b,T).
\]
The complementary direct cancellation does hold in this example.  Thus the
six-vertex evidence was seeing a true small-case phenomenon, not the correct
general theorem.  Direct cancellation remains useful as a size-drop move when
available, but it cannot be the main endpoint theorem.

This same example also rules out the tempting shortcut ``low-slice equality
gives either the direct endpoint match or an internal zero-error edge.''  Here
\(T=\{1\}\), and the unique automorphism carrying
\(\{1,2\}=T\cup\{b\}\) to \(\{0,1\}=T\cup\{a\}\) sends \(1\) to \(0\), so no
solving automorphism fixes the internal deleted \(T\)-vertex.  Thus the
zero-error edge lemma is a local win condition, but a proof of one-slice orbit
reconstruction still needs a transport/descent argument that can move the
error along the path rather than demanding a zero-error internal edge at the
start.
\end{remark}

\begin{conjecture}[Active zero-star pair]
Assume the low color-size slice agrees:
\[
  \Deck_{T\cup\{a\}}(K,S,T\cup\{a\})
  =
  \Deck_{T\cup\{b\}}(K,S,T\cup\{b\}).
\]
Then there exist \(x\in T\cup\{a\}\), \(y\in T\cup\{b\}\), and a colored card
isomorphism
\[
  (K-x,S-x,(T\cup\{a\})-\{x\})
  \cong
  (K-y,S-y,(T\cup\{b\})-\{y\})
\]
whose deleted-star error from \(x\) to \(y\) is zero and for which
\(x\in S\iff y\in S\).
\end{conjecture}

This is currently the cleanest local target.  It is stronger than one-slice
orbit reconstruction only in presentation, not in spirit: the formalized
zero-star extension lemma immediately turns such a pair into the desired
automorphism.  Conversely, any automorphism carrying \(T\cup\{a\}\) to
\(T\cup\{b\}\) restricts to such a zero-star card pair, so Lean now proves
that the active zero-star pair condition is equivalent to the fixed-host orbit
conclusion.  The point is not to strengthen the theorem, but to identify the
kind of local witness the one-slice deck must reveal.

It also survives the cyclic \(n=9\) example: the active zero-star pairs are
\((0,1)\) and \((1,2)\), corresponding to the two steps of the cyclic
automorphism \(0\mapsto1\mapsto2\).  Exhaustive \(n=5\) probing found
\(24832\) low-slice equal states and \(24832\) active zero-star pairs.  A
bounded \(n=6\) graph-atlas prefix found \(27922\) low-slice equal states and
\(27922\) active zero-star pairs, and a bounded \(n=7\) graph-atlas prefix
found \(49888\) low-slice equal states and \(49888\) active zero-star pairs.
Sampled random-host probes found \(896/896\) active zero-star pairs at
\(n=8\) and \(263/263\) at \(n=9\).  In those larger probes, most witnesses
were internal active pairs rather than the endpoint pair \(a\to b\), which is
good evidence that the target has the right flexibility.

\begin{definition}[Minimum-error active-card graph]
Form a bipartite graph \(\Gamma\) with left part \(T\cup\{a\}\) and right part
\(T\cup\{b\}\).  Put an edge \(x\to y\) when the colored cards
\[
  (K-x,S-x,(T\cup\{a\})-\{x\}),\qquad
  (K-y,S-y,(T\cup\{b\})-\{y\})
\]
are isomorphic and \(x\in S\iff y\in S\).  Give this edge weight
\[
  \epsilon(x,y)=
  \min_{\phi:x\to y}
  \left|\{z\neq x: K(x,z)\not\leftrightarrow K(y,\phi(z))\}\right|,
\]
where \(\phi\) ranges over all color-preserving card isomorphisms between
these two cards.  The zero-error subgraph \(\Gamma_0\) consists of edges of
weight \(0\).
\end{definition}

Low-slice equality says that \(\Gamma\) has enough card edges to support a
perfect matching of the active deletion deck.  The active zero-star pair
conjecture says the stronger-looking but equivalent statement
\[
  E(\Gamma_0)\neq\emptyset.
\]
Thus a minimal counterexample is exactly a low-slice deck equality whose
minimum-error active-card graph has no zero edge.  This is a clean local
obstruction: every visible active card match exists only after changing at
least one adjacency incident to the deleted vertex.  The remaining proof
burden is to show that such a positive minimum-error obstruction can be
recentered at one of its first changed adjacencies to produce a strictly
smaller active-card obstruction, or else to reduce the total error of the
chosen deck matching.

For finite hosts, the condition \(x\in S\iff y\in S\) is not really an extra
mystery: a color-preserving card isomorphism sees the number of first-colored
vertices left on the card, so it can distinguish whether the deleted vertex
belonged to \(S\).  The genuinely hidden data is the deleted star, i.e. the
adjacency vector from the missing vertex into its card.

\begin{lemma}[Passive status is card-visible]
Let \(K\) be finite.  If
\[
  \phi:(K-x,S-x,U-x)\cong(K-y,S-y,U'-y)
\]
preserves the first color on the deleted cards, then
\[
  x\in S\iff y\in S.
\]
\end{lemma}

\begin{proof}[Proof sketch]
The isomorphism identifies the first-colored vertices surviving in the two
cards.  If \(x\in S\), the left card has one fewer first-colored vertex than
\(S\); if \(x\notin S\), it has the same number.  The same dichotomy holds on
the right.  Equal surviving counts force the same deleted-vertex status.
Lean now formalizes the cardinality lemmas as
\[
\begin{array}{l}
\texttt{FixedHostCardIsoData.deletedFirstColor\_card\_lt\_of\_mem},\\
\texttt{FixedHostCardIsoData.deletedFirstColor\_card\_eq\_of\_not\_mem},
\end{array}
\]
and the visibility statement as
\[
  \texttt{FixedHostCardIsoData.firstColor\_deleted\_status\_iff}.
\]
\end{proof}

\begin{lemma}[First-error localization]
Assume there is no active zero-star pair.  Let \(x\in T\cup\{a\}\),
\(y\in T\cup\{b\}\), and let
\[
  \phi:(K-x,S-x,(T\cup\{a\})-\{x\})
  \cong
  (K-y,S-y,(T\cup\{b\})-\{y\})
\]
be any color-preserving active card isomorphism with \(x\in S\iff y\in S\).
Then there is a surviving vertex \(z\neq x\) such that
\[
  K(x,z)\not\leftrightarrow K(y,\phi(z)).
\]
Moreover, deleting \(z\) as well gives a common two-hole colored card
\[
  (K-\{x,z\},S-\{x,z\},(T\cup\{a\})-\{x,z\})
  \cong
  (K-\{y,\phi(z)\},S-\{y,\phi(z)\},(T\cup\{b\})-\{y,\phi(z)\}).
\]
\end{lemma}

\begin{proof}[Proof sketch]
If no such \(z\) existed, the card isomorphism would preserve the entire
deleted star of \(x\), hence would extend to an automorphism of \(K\) carrying
\(x\) to \(y\) and \(T\cup\{a\}\) to \(T\cup\{b\}\).  That would be an active
zero-star pair.  The two-hole statement is just the restriction of \(\phi\) to
the vertices other than \(x\) and \(z\).  This localization lemma is now
formalized in Lean: a no-pair obstruction produces an explicit star mismatch
and a two-hole colored isomorphism over it.
\end{proof}

Equivalently, Lean proves the normal form
\[
  \text{no active zero-star pair}
  \quad\Longleftrightarrow\quad
  \text{every active card match has a first star mismatch}.
\]
Thus the local obstruction is no longer a vague absence statement.  It is a
family of concrete two-hole witnesses, one above every visible active card
edge.  The next step is to show that a minimal such family cannot be closed
under the active deck matching: following first mismatches either reaches a
zero-star pair or creates a strictly smaller two-hole obstruction.

\begin{proposition}[Formal local obstruction]
Define \(\mathcal L(K,S,T,a,b)\) to be the assertion that every active card
isomorphism
\[
  \phi:(K-x,S-x,(T\cup\{a\})-\{x\})
  \cong
  (K-y,S-y,(T\cup\{b\})-\{y\})
\]
with \(x\in S\iff y\in S\) has a first-error witness \(z\neq x\), together
with the induced two-hole colored isomorphism after also deleting \(z\) and
\(\phi(z)\).  Then
\[
  \mathcal L(K,S,T,a,b)
  \quad\Longleftrightarrow\quad
  \text{there is no active zero-star pair}.
\]
Moreover, each localized first-error witness preserves both colors:
\[
  z\in S\iff \phi(z)\in S,\qquad
  z\in T\cup\{a\}\iff \phi(z)\in T\cup\{b\}.
\]
Thus every first error is either active-active or inactive-inactive; there are
no cross-color local errors.
\end{proposition}

\begin{proof}[Lean status]
This proposition is now formalized by the Lean declarations
\[
\begin{array}{l}
\texttt{LowSliceLocalObstruction},\\
\texttt{LowSliceZeroStarPair.localObstruction\_iff\_no\_pair}.
\end{array}
\]
The individual localized witnesses are packaged as
\texttt{LowSliceCardFirstError}.  The color transport facts are formalized as
\[
\begin{array}{l}
\texttt{LowSliceCardFirstError.sourceFirst\_iff\_targetFirst},\\
\texttt{LowSliceCardFirstError.sourceActive\_iff\_targetActive},\\
\texttt{LowSliceCardFirstError.active\_or\_inactive}.
\end{array}
\]
\end{proof}

\begin{remark}[Chosen low-slice matching]
The low-slice equality is existential, but in the formal campaign we can choose
a concrete matching of the active deletion deck.  Lean packages this as
\[
\begin{array}{l}
\texttt{LowSliceIsoData.ofSameDeck},\\
\texttt{LowSliceIsoData.first\_status},\\
\texttt{LowSliceIsoData.firstError\_of\_no\_pair}.
\end{array}
\]
Thus, over finite hosts, a low-slice deck equality plus the no-pair assumption
produces a first-error witness above every edge of the chosen matching without
requiring any separate passive-color-status hypothesis.
\end{remark}

\begin{corollary}[Active/inactive split]
The formal local obstruction is equivalent to the following split version: above
every visible active card edge, there is either an active-active first error
\[
  z\in T\cup\{a\},\qquad \phi(z)\in T\cup\{b\},
\]
or an inactive-inactive first error
\[
  z\notin T\cup\{a\},\qquad \phi(z)\notin T\cup\{b\}.
\]
There is no mixed case.
\end{corollary}

\begin{proof}[Lean status]
This is formalized as
\[
\begin{array}{l}
\texttt{LowSliceCardFirstError.Active},\\
\texttt{LowSliceCardFirstError.Inactive},\\
\texttt{LowSliceLocalSplitObstruction},\\
\texttt{LowSliceLocalObstruction.iff\_split}.
\end{array}
\]
\end{proof}

\begin{definition}[Formal minimum-error obstruction]
For a chosen low-slice matching \(M\), define its total star error by
\[
  E(M)=\sum_{x\in T\cup\{a\}}
  \left|\{z\neq x:K(x,z)\not\leftrightarrow K(Mx,Mz)\}\right|.
\]
A minimum-error matching is a chosen low-slice matching whose total star error
is minimal among all chosen low-slice matchings.  In Lean this is represented
by
\[
\begin{array}{l}
\texttt{FixedHostCardIsoData.starErrorCount},\\
\texttt{LowSliceIsoData.totalStarErrorCount},\\
\texttt{LowSliceMinimumErrorMatching}.
\end{array}
\]
The formal theorem
\[
  \texttt{LowSliceMinimumErrorMatching.noZeroStarEdge\_iff\_no\_pair}
\]
says that, for a minimum-error matching, having no zero-star edge is equivalent
to the no-active-zero-star-pair obstruction.  Moreover
\[
  \texttt{LowSliceMinimumErrorMatching.firstError\_active\_or\_inactive\_of\_no\_pair}
\]
gives the branch split above each edge of such an obstruction.

In the active branch, the first-error witness also gives a formal recentering
operation.  If the first error above \(x\mapsto y\) occurs at an active
surviving vertex \(z\), then \(z\in T\cup\{a\}\) and \(\phi(z)\in T\cup\{b\}\).
Lean packages this as
\[
\begin{array}{l}
\texttt{LowSliceMinimumErrorMatching.ActiveFirstError.source},\\
\texttt{LowSliceMinimumErrorMatching.ActiveFirstError.target},\\
\texttt{LowSliceMinimumErrorMatching.ActiveFirstError.recenteredTwoHoleIso}.
\end{array}
\]
This is exactly the two-hole transport one would need for a descent or splice
argument.  The missing mathematical step is to upgrade, splice, or compose this
two-hole transport so that it either yields a zero-star pair or lowers total
star error.

The active branch now has a more precise exchange interpretation.  If an active
first error above \(x\mapsto Mx\) points at an active source \(z\) and active
target \(z'\), define the successor of the error to be
\[
  \sigma(x,E)=M^{-1}(z').
\]
There are then two cases.
\[
\begin{array}{ll}
\text{coherent:} & z=\sigma(x,E),\\
\text{noncoherent:} & z\neq \sigma(x,E).
\end{array}
\]
The coherent case says that the first-error target is already matched to its
source by the chosen minimum matching.  This should be a side-cycle or
self-contained two-hole defect, and the intended lemma is that it can be
removed unless it exposes a smaller fixed-host obstruction.  The noncoherent
case is the genuine alternating exchange step: it points from the current
matched edge to a different matched edge.  If every productive local lowering
move fails, these noncoherent steps should assemble into a closed active
correction cycle.

Lean records exactly this successor/coherence structure as
\[
\begin{array}{l}
\texttt{LowSliceMinimumErrorMatching.ActiveFirstError.nextLeft},\\
\texttt{LowSliceMinimumErrorMatching.ActiveFirstError.Coherent},\\
\texttt{LowSliceMinimumErrorMatching.ActiveFirstError.Noncoherent},\\
\texttt{LowSliceMinimumErrorMatching.ActiveFirstError.coherentTwoHoleIso},\\
\texttt{LowSlicePositiveMinimumObstruction.ActiveTransportStep},\\
\texttt{LowSlicePositiveMinimumObstruction.ActiveExchangeRel},\\
\texttt{LowSlicePositiveMinimumObstruction.ActiveNoncoherentRel}.
\end{array}
\]
Inside the \(\texttt{ActiveFirstError}\) namespace, the coherent-target theorem
says that a coherent active target is exactly the matched mate of the active
source, so \(\texttt{coherentTwoHoleIso}\) is a genuine two-hole transport
between matched active pairs.  This is the formal side-cycle shape.
One should not confuse this with immediate removability: the defect is still
seen from the card deleting \(x\), and Lean records the genuine separation
\[
  z\neq x,\qquad z'\neq Mx
\]
as \(\texttt{source\_ne\_base}\) and \(\texttt{target\_ne\_baseTarget}\).
It also proves \(\sigma(x,E)\neq x\).  Thus the coherent case is a matched-pair
two-hole defect witnessed from a third position, not an automatic cancellation;
and the noncoherent case is even stricter, since
\(\texttt{three\_distinct\_of\_noncoherent}\) says that the base \(x\), source
\(z\), and successor \(\sigma(x,E)\) are all distinct.
There are now two directed graphs naturally attached to active errors.  The
relation \(\texttt{ActiveExchangeRel}\) has edge
\[
  z\longrightarrow M^{-1}(z')
\]
for each active first-error witness; this is the correction graph suggested by
the mismatched image of the source.  The relation
\(\texttt{ActiveObserverRel}\) has edge
\[
  x\longrightarrow M^{-1}(z'),
\]
which follows the deleted card from which the defect was observed.  Both
relations are loop-free in their useful parts.  Lean records the elementary
consequence
\[
  \texttt{LowSlicePositiveMinimumObstruction.activeNoncoherentRel\_ne},
\]
so noncoherent correction edges always connect distinct active vertices; it
also proves the analogous observer-ne theorem.

If every edge of a positive obstruction admits an active branch, we may choose
one active first-error witness above each edge.  This gives an
\(\texttt{ActiveObserverSystem}\), i.e. a self-map
\[
  \nu:T\cup\{a\}\to T\cup\{a\},\qquad
  \nu(x)=M^{-1}(z'_x).
\]
The separation theorem gives \(\nu(x)\neq x\) for every \(x\).  Since the active
slice is finite and nonempty, the generic finite-dynamics lemma
\[
  \texttt{exists\_positive\_iterate\_eq\_self\_of\_finite}
\]
implies that \(\nu\) has a nontrivial periodic orbit.  Lean packages this as
\[
\begin{array}{l}
\texttt{LowSlicePositiveMinimumObstruction.ActiveObserverSystem},\\
\texttt{ActiveObserverSystem.exists\_observerCycle},\\
\texttt{exists\_activeObserverCycle\_of\_allActiveBranches},\\
\texttt{LowSlicePositiveMinimumObstruction.ActiveObserverCycle}.
\end{array}
\]
Thus the all-active branch really does reduce to a closed observer cycle.  The
remaining question is whether that observer cycle can be spliced to lower
\(E(M)\), or whether it composes into a zero-star pair.  Lean also proves that
each consecutive pair on the packaged cycle is an actual observer-relation edge
(\(\texttt{ActiveObserverCycle.edge\_at}\)) and that the final edge returns to
the base (\(\texttt{ActiveObserverCycle.last\_edge\_to\_base}\)).
The cycle itself has a second split:
\[
\begin{array}{l}
\texttt{ActiveObserverCycle.allCoherent\_or\_hasNoncoherent},\\
\texttt{ActiveObserverCycle.noncoherentCorrectionEdge\_at}.
\end{array}
\]
So a closed active observer cycle is either entirely coherent, in which case it
has the shape of a removable side cycle, or it contains a noncoherent index,
which supplies a genuine correction edge from the defect source to the next
observer vertex.
Lean names the two active subforks
\[
\begin{array}{l}
\texttt{NoAllCoherentActiveCycleFork},\\
\texttt{NoNoncoherentActiveCycleFork},
\end{array}
\]
and proves
\[
  \texttt{NoActiveObserverCycleFork.of\_no\_active\_cycle\_subforks}.
\]
Thus the active branch has split into two precise no-go statements: no purely
coherent side cycle, and no observer cycle carrying a noncoherent correction
edge.
For the all-coherent subfork, Lean now proves the side-cycle transport
available on each edge:
\[
\begin{array}{l}
\texttt{ActiveObserverSystem.coherentMatchedTwoHoleIsoAt},\\
\texttt{ActiveObserverCycle.coherentMatchedTwoHoleIso\_at}.
\end{array}
\]
Concretely, if \(x_{i+1}=\nu(x_i)\) and the selected witness at \(x_i\) is
coherent, then there is a two-colored isomorphism
\[
  K-\{x_i,x_{i+1}\}
  \cong
  K-\{Mx_i,Mx_{i+1}\}
\]
with the appropriate first and second colors.  Hence an all-coherent active
cycle is literally a cycle of matched two-hole transports.  The next missing
lemma is the side-cycle removal lemma: such a closed coherent transport cycle
should either splice to reduce total star error or expose a smaller fixed-host
obstruction.

Thus the next purely mathematical lemma is no longer vague descent.  It is:
in a positive minimum-error obstruction, coherent active steps are removable,
and noncoherent active steps cannot form an essential closed correction cycle
without either lowering \(E(M)\) or producing an active zero-star pair.

The inactive branch has a parallel split.  An inactive first-error witness
points from a source in \(O\cup\{b\}\) to a target in \(O\cup\{a\}\), where
\(O=V\setminus(T\cup\{a,b\})\).  It is endpoint-exposing if the source is
\(b\) or the target is \(a\), and it is an outer residual if both source and
target lie in \(O\).  Lean records this in the
\(\texttt{InactiveFirstError}\) namespace as
\[
\begin{array}{l}
\texttt{EndpointExposing},\\
\texttt{OuterResidual},\\
\texttt{endpointExposing\_or\_outerResidual}.
\end{array}
\]
The mathematical hope is that endpoint-exposing inactive errors are visible to
restricted Kelly counts involving the distinguished endpoints, while a purely
outer residual is a smaller fixed-host obstruction living away from the switch
endpoints.  The theorem
\(\texttt{not\_endpointExposing\_iff\_outerResidual}\) records that, once
endpoint exposure is ruled out, this outer residual is the only remaining
inactive case.

Similarly, if every edge admits an inactive branch, Lean chooses an
\(\texttt{InactiveObserverSystem}\) and proves that the selected witnesses are
either endpoint-exposing everywhere or include an explicit outer residual:
\[
\begin{array}{l}
\texttt{InactiveObserverSystem.allEndpointExposing\_or\_hasOuterResidual},\\
\texttt{exists\_inactiveObserverSystem\_of\_allInactiveBranches}.
\end{array}
\]
It also names the inactive subforks
\[
\begin{array}{l}
\texttt{NoAllEndpointExposingInactiveFork},\\
\texttt{NoOuterResidualInactiveFork},
\end{array}
\]
with reduction theorem
\[
  \texttt{NoInactiveEndpointOuterFork.of\_no\_inactive\_subforks}.
\]

Combining the active and inactive choices gives the current formal strategic
fork:
\[
  \texttt{LowSlicePositiveMinimumObstruction.strategy\_fork}.
\]
Every positive minimum-error obstruction yields one of:
\[
\begin{array}{ll}
\text{(A)} & \text{a nontrivial active observer cycle},\\
\text{(I)} & \text{an all-inactive endpoint/outer split},\\
\text{(M)} & \text{genuinely mixed active and inactive branches}.
\end{array}
\]
So the problem has been reduced to ruling out these three concrete obstruction
types.  Lean records the exact reduction as
\[
\begin{array}{l}
\texttt{NoActiveObserverCycleFork},\\
\texttt{NoInactiveEndpointOuterFork},\\
\texttt{NoMixedBranchFork},\\
\texttt{NoLowSlicePositiveMinimumObstruction.of\_no\_strategy\_forks}.
\end{array}
\]
Thus a proof of the three branch no-go theorems would prove the
minimum-error-zero route, and hence the zero-star-pair route, by the earlier
equivalence.
After splitting the active and inactive branches further, Lean also proves the
refined five-branch reduction
\[
  \texttt{NoLowSlicePositiveMinimumObstruction.of\_no\_refined\_forks}.
\]
The five no-go targets are:
\[
\begin{array}{ll}
\text{(A1)} & \text{no all-coherent active observer cycle},\\
\text{(A2)} & \text{no active observer cycle with a noncoherent correction edge},\\
\text{(I1)} & \text{no all-inactive endpoint-exposing system},\\
\text{(I2)} & \text{no all-inactive outer-residual system},\\
\text{(M)} & \text{no genuinely mixed active/inactive obstruction}.
\end{array}
\]

The computational minimum-error probe sharpens the picture.  Exhaustively over
the \(n=5\) graph-atlas hosts, all \(24832\) low-slice equal states had minimum
total star error \(0\).  A bounded \(n=6\) graph-atlas run likewise found
\(27922\) low-slice equal states and \(27922\) minimum-total-zero states.  There
were no positive minimum-error states in these probes and hence no finite
no-descent fixture to inspect.  This does not prove the descent theorem, but it
changes the conjectural shape: the hoped-for theorem is not merely that a
positive minimum obstruction descends, but perhaps that low-slice equality
already forces minimum total star error \(0\).

Lean now proves the exact reduction.  Low-slice equality always admits a
minimum-error matching, by applying the well-ordering principle to the set of
natural total-error values:
\[
  \texttt{LowSliceMinimumErrorMatching.exists\_of\_sameDeck}.
\]
For such a matching,
\[
  E(M)=0
  \quad\Longleftrightarrow\quad
  \text{there is an active zero-star pair},
\]
formalized as
\[
  \texttt{LowSliceMinimumErrorMatching.totalStarErrorCount\_eq\_zero\_iff\_pair}.
\]
Thus the current best target is the minimum-error-zero conjecture:
\[
  \text{low-slice equality}\quad\Longrightarrow\quad
  \min_M E(M)=0.
\]
Lean packages this route as
\[
\begin{array}{l}
\texttt{LowSliceMinimumErrorZeroConjecture},\\
\texttt{LowSliceZeroStarPairConjecture.of\_minimumErrorZero}.
\end{array}
\]
It also packages the exact failure object:
\[
\begin{array}{l}
\texttt{LowSlicePositiveMinimumObstruction},\\
\texttt{LowSlicePositiveMinimumObstruction.nonempty\_iff\_not\_minimumErrorZero},\\
\texttt{NoLowSlicePositiveMinimumObstruction},\\
\texttt{NoLowSlicePositiveMinimumObstruction.iff\_minimumErrorZero}.
\end{array}
\]
Such an obstruction automatically has no zero-star pair, no zero-star edge in
its minimum matching, and an active/inactive first-error split above every
matched edge.
\end{definition}

The one-slice computational probe still supports the orbit version through the
exhaustive six-vertex atlas scan: low-slice equality, complementary-slice
equality, full deck equality, and fixed-host orbit success have exactly the
same state count.  A further seven-vertex atlas scan with \(S=\emptyset\)
checked \(1403136\) singleton states.  It found \(143488\) exact low-slice
equalities and no one-slice orbit failures.  Random NetworkX probes on
eight- and nine-vertex hosts also found no failures in \(400000\) additional
states.

The same six-vertex scan found that this common count is also exactly the
number of singleton pairs in the same orbit under \(\operatorname{Aut}(K,S)\).
This suggests the surviving factorization of the proof:
\[
  \text{one-slice equality}
  \Longrightarrow
  T\cup\{a\}\sim_{\operatorname{Aut}(K,S)}T\cup\{b\}.
\]
This is a reconstruction statement for a marked deletion deck.  It is sharper
than the original rectangular no-kernel claim, but it does not force the
specific direct endpoint cards.

\begin{lemma}[Restricted Kelly principle for a marked deletion deck]
Let \(X=(K,S,U)\) be a two-colored fixed-host graph, and let
\(\Deck_U(X)\) be the multiset of cards obtained by deleting only vertices of
the second color \(U\).  For any two-colored pattern \(P\) using
\(s<|U|\) second-colored vertices, the number of induced colored copies of
\(P\) in \(X\) is determined by \(\Deck_U(X)\).
\end{lemma}

\begin{proof}[Proof sketch]
Double-count pairs \((u,\iota)\), where \(u\in U\) and \(\iota\) is an induced
colored copy of \(P\) in \(X-u\).  A fixed copy of \(P\) in \(X\) avoids exactly
\(|U|-s\) vertices of \(U\).  Hence
\[
  \sum_{u\in U} N_P(X-u)=(|U|-s)N_P(X).
\]
Since \(|U|-s\neq 0\), the marked deletion deck determines \(N_P(X)\).
\end{proof}

Apply this with \(U=T\cup\{a\}\).  The endpoint card \(H_a\) uses only
\(|U|-1\) second-colored vertices, so one-slice equality forces the type
\(H_a\) to occur on the \(b\)-side.  If it does not occur as \(H_b\), it must
occur as some \(B_t\), producing the first edge of a type-flow path.  Iterating
gives the shortest path
\[
  H_a=Y_0\to Y_1\to\cdots\to Y_r=H_b.
\]
Thus the only information a nontrivial path can hide from the one-slice deck is
the extension profile of the single deleted second-colored vertex.  A proof of
one-slice orbit reconstruction should show that transporting this hidden
extension profile around a shortest path either produces a fixed-host
automorphism carrying \(T\cup\{a\}\) to \(T\cup\{b\}\), or creates a smaller
centered-extension obstruction.

\begin{definition}[Low-slice transport path]
A low-slice transport path is a shortest type-flow path
\[
  H_a=Y_0\to Y_1\to\cdots\to Y_r=H_b
\]
together with choices of colored card isomorphisms realizing each equality
along the path.  Equivalently, it is a sequence of deleted vertices
\[
  x_0=a,\quad x_1,\ldots,x_r\in T,\quad x_{r+1}=b
\]
and card isomorphisms between the successive punctured fixed hosts obtained by
deleting \(x_i\) and \(x_{i+1}\), with the active color transported from
\(T\cup\{a\}\) toward \(T\cup\{b\}\).
\end{definition}

\begin{lemma}[A zero-error transport edge solves the orbit problem]
Suppose some edge on a low-slice transport path is realized by a colored card
isomorphism
\[
  (K-x,S-x,(T\cup\{a\})-\{x\})
  \cong
  (K-y,S-y,(T\cup\{b\})-\{y\})
\]
where \(x\in T\cup\{a\}\), \(y\in T\cup\{b\}\), the deleted vertices have the
same \(S\)-status, and the star-extension error from \(x\) to \(y\) is zero.
Then \(T\cup\{a\}\) and \(T\cup\{b\}\) lie in the same
\(\operatorname{Aut}(K,S)\)-orbit.
\end{lemma}

\begin{proof}[Proof sketch]
The zero-error condition says exactly that the card isomorphism extends across
the deleted vertex by sending \(x\) to \(y\).  The extension is an automorphism
of \(K\) preserving \(S\).  The card isomorphism preserves the displayed second
colors away from the deleted vertices, and the hypotheses \(x\in T\cup\{a\}\),
\(y\in T\cup\{b\}\) handle the reinserted vertex.  Hence the extension carries
\(T\cup\{a\}\) to \(T\cup\{b\}\).

This lemma has now been formalized in Lean as the general card-extension
result for chosen fixed-host card isomorphisms, with the previous same-deleted
vertex transport edge as a special case.  The formal statement makes the
hidden condition completely explicit: a chosen punctured-card isomorphism must
transport the deleted vertex's neighbor set to the neighbor set of the target
deleted vertex.  Lean then extends that bijection by sending the deleted source
vertex to the deleted target vertex and checks both the graph automorphism
condition and the two color-preservation conditions.
\end{proof}

Lean also proves the matching-level version: a full automorphism induces a
perfect zero-star matching of the whole active deletion deck, and any such
matching gives a zero-star pair by choosing the card over \(a\).  Hence the
orbit conclusion, the active zero-star pair, and the existence of a perfect
zero-star matching are three equivalent formulations of the same target.  The
advantage of the matching formulation is tactical: low-slice equality already
provides a perfect matching of card types, so the remaining question is whether
one can lower its total star error to zero.

\begin{lemma}[Every transport edge is erroneous in a minimal orbit failure]
In a minimal counterexample to one-slice orbit reconstruction, no edge on any
low-slice transport path admits a zero-error realization.
\end{lemma}

\begin{proof}[Proof sketch]
If a transport edge had a zero-error realization, the previous lemma would
give an automorphism of \((K,S)\) carrying \(T\cup\{a\}\) to \(T\cup\{b\}\),
contradicting orbit failure.
\end{proof}

\begin{lemma}[First-error descent for transport]
There is no minimal one-slice orbit failure in which every transport edge has
unavoidable nonzero star-extension error.
\end{lemma}

\begin{proof}[Proof sketch]
Choose a shortest transport path, then among those minimize the total
star-error support over the chosen card isomorphisms.  Pick an edge whose
nonzero star error is minimal.  Recenter at its common punctured card and
regard the failed extension across the deleted vertex as a smaller colored
centered-extension problem.  If the recentered problem has a zero-error
representative, replace the offending transport edge and lower the error
potential.  If it has no zero-error representative, it is a smaller
counterexample to the centered-extension theorem.  Minimality forbids both
outcomes.
\end{proof}

\begin{proposition}[Transport integration proves one-slice orbit reconstruction]
Assume the first-error descent lemma for all smaller centered-extension
states.  Then low-slice equality implies fixed-host orbit success.
\end{proposition}

\begin{proof}[Proof sketch]
Restricted Kelly gives at least one low-slice transport path from \(H_a\) to
\(H_b\).  In a minimal counterexample, choose it shortest and then
minimum-error.  If some edge is zero-error, the zero-error edge lemma gives
orbit success.  If no edge is zero-error, first-error descent gives a smaller
centered-extension obstruction.  Both contradict minimality.
\end{proof}

One should not try to prove the stronger statement that the whole low
type-flow graph is trivial.  It is false in small solved states.  For example,
the \(n=5\) exhaustive probe found \(1152\) equal-deck singleton states with
some internal \(T\)-edge satisfying \(A_t\not\cong B_t\), even though the
endpoint cancellation \(H_a\cong H_b\) always held.  Moreover, the endpoint
card type itself can coincide with internal detour cards in solved states; in
the same \(n=5\) probe this happened in \(6472\) states.  Thus the proof has to
distinguish orbit success from endpoint cancellation and from uniqueness of the
endpoint type.  Harmless internal circulation, even circulation incident to the
endpoint card type, may still exist.

The natural proof attempt is a paired minimality argument.  Suppose, for
example, that the low endpoint defect \([H_b]-[H_a]\) is nonzero.  Its named
\(O\)-boundary is the rectangular chain.  If the complementary endpoint defect
vanishes, then the rectangular identity says this boundary is zero, so the low
defect is a named deletion cycle.  If the complementary endpoint defect is
nonzero, then the two nonzero endpoint defects share the same rectangular
boundary.  In both cases, choose a smallest colored separator for one endpoint
defect and run the support-escalation argument.  The computation suggests
that the terminal case is impossible even without assuming the other endpoint
defect is nonzero.

The no-kernel conjecture remains the minimal statement needed for the
size-drop induction.  A counterexample would be a nonzero rectangular kernel
element: two different endpoint defects whose named second deletions across
the opposite side of the partition are indistinguishable.  The double direct
cancellation conjecture says that even this weaker kind of endpoint defect
cannot occur.  Either obstruction is much more specific than an arbitrary
homology class of the deletion complex.

The important nuance is that the rectangular chain is the common image of two
endpoint defects, not the defect itself.  A moment of the \(D_{t,o}^\ast\)
cards only studies this common image.  To kill the kernel, one has to show
that a nonzero endpoint defect cannot be invisible to its named partial
boundary.

\begin{conjecture}[Full-support rectangular squeeze]
Work in a minimal fixed-host singleton counterexample, minimizing first
\(|V(K)|\) and then the total length of the low and complementary type-flow
paths.  Assume both endpoint cancellations fail.  Let \(P\) be a smallest
colored pattern that distinguishes the two low endpoints \(H_a\) and
\(H_b\).  If a copy of \(P\) responsible for the difference misses some
vertex of \(O\), then the matching term in
\(\partial_T([C_a^b]-[C_b^a])\) produces either a smaller fixed-host
singleton obstruction, a path shortening, or a direct cancellation.  Hence,
in a minimal unresolved counterexample, every such low separator must use all
vertices of \(O\).  Dually, every complementary separator must use all
vertices of \(T\).  If the support can then be escalated across the other
side, the only nonproductive exposed terms left are terminal two-hole card
imbalances.
\end{conjecture}

This is now the best candidate for the ``make or break'' lemma.  If it is
true, the fixed-host singleton theorem follows by the two size-drop lemmas.
If it is false, the strategy has not merely failed to find the right charge:
there is a genuine finite rectangular kernel, and the centered-extension
campaign needs a new ingredient that is not just another Kelly-style moment.

The intuition is simple.  The partial boundary \(\partial_O([H_b]-[H_a])\)
detects every low-endpoint discrepancy that survives after deleting at least
one vertex of \(O\).  Such a visible discrepancy is not automatically a
contradiction, because the rectangular identity says that it is matched by
the complementary partial boundary.  But that match is a two-hole card
isomorphism, and therefore should be exactly the kind of smaller detour from
which we can extract a path shortening or a size drop.  If no such extraction
is possible, the first visible difference must be supported on all of
\(O\).  The dual argument forces the complementary first difference to be
supported on all of \(T\).  This is not yet a contradiction by itself; it is a
support-escalation step.  A genuine kernel would have to push every
nonterminal visible difference into a whole two-hole card imbalance.  That is
a very small corner of the search space, and it is the corner to attack next.

\begin{remark}[Candidate proof of the squeeze]
Assume a minimal counterexample to the fixed-host singleton theorem in which
both endpoint cancellations fail.

\begin{enumerate}[leftmargin=2em]
  \item By the smaller colored centered theorem, the nonisomorphic low
  endpoints \(H_a,H_b\) cannot have identical colored decks.  Choose a
  deleted low-endpoint card type \(Z\) whose multiplicity differs, minimizing
  the size of a colored pattern detecting the difference.

  \item If \(Z\) arises by deleting a vertex of \(O\), then \(Z\) appears in
  \(\partial_O([H_b]-[H_a])\).  The rectangular identity matches this
  occurrence inside \(\partial_T([C_a^b]-[C_b^a])\).  Unwinding the matched
  two-hole card gives a smaller fixed-host singleton state over
  \(K-\{t,o\}\), or else an isomorphism that splices one of the two type-flow
  paths.  Minimality rules this out unless it is already a direct
  cancellation.

  \item Therefore every first low deck discrepancy must come from deleting a
  vertex of \(T\).  Symmetrically, every first complementary deck discrepancy
  must come from deleting a vertex of \(O\).

  \item Repeat the argument one level down.  A low discrepancy obtained by
  deleting \(t\in T\) is a singleton difference inside the smaller host
  \(K-t\).  If any detecting copy misses a vertex of \(O\), the same
  rectangular matching produces a smaller obstruction or a path shortening.
  Hence the detecting copy must use all of \(O\).

  \item The dual argument forces the complementary detecting copy to use all
  of \(T\).  If either detector still omits a vertex on the other side, the
  previous step applies again and gives a splice.  Thus an unresolved exposed
  term must be terminal: it is an imbalance of entire two-hole card types, not
  a smaller pattern-count separator.
\end{enumerate}

The fragile step is the second one: a matched two-hole card must really
produce a smaller singleton obstruction or a valid splice.  This is the exact
place where the approach can break.  The good news is that it is now a
finite, local statement about two deleted vertices, not a global statement
about arbitrary graph reconstruction.
\end{remark}

\subsection{The two-hole splice lemma}

We now isolate the local statement hidden in the previous paragraph.  A
rectangular cancellation pairs a positive two-hole card with a negative
two-hole card.  Thus its basic datum is
\[
  t,t'\in T,\qquad o,o'\in O,\qquad
  \theta:D_{t,o}^a\cong D_{t',o'}^b.
\]
Unpacked, \(\theta\) is a color-preserving isomorphism
\[
  (K-\{t,o\},S\setminus\{t,o\},(T\setminus\{t\})\cup\{a\})
  \cong
  (K-\{t',o'\},S\setminus\{t',o'\},(T\setminus\{t'\})\cup\{b\}).
\]

\begin{definition}[Two-hole ports]
For such a two-hole match \(\theta\), define its low ports by looking at the
second color:
\[
  \theta(a)\in (T\setminus\{t'\})\cup\{b\},
  \qquad
  \theta^{-1}(b)\in (T\setminus\{t\})\cup\{a\}.
\]
If \(\theta(a)\neq b\), call \(\theta(a)\) the low exit.  If
\(\theta^{-1}(b)\neq a\), call \(\theta^{-1}(b)\) the low entry.

Define its complementary ports by looking at the vertices that are not
second-colored:
\[
  \theta(b)\in (O\setminus\{o'\})\cup\{a\},
  \qquad
  \theta^{-1}(a)\in (O\setminus\{o\})\cup\{b\}.
\]
If \(\theta(b)\neq a\), call \(\theta(b)\) the complementary exit.  If
\(\theta^{-1}(a)\neq b\), call \(\theta^{-1}(a)\) the complementary entry.
\end{definition}

\begin{lemma}[Two-hole color bookkeeping]
Every two-hole match has exactly the four alternatives described above:
the image of \(a\) is either the moved second-colored vertex \(b\) or a
vertex of \(T\setminus\{t'\}\); the preimage of \(b\) is either \(a\) or a
vertex of \(T\setminus\{t\}\); the image of \(b\) is either \(a\) or a
vertex of \(O\setminus\{o'\}\); and the preimage of \(a\) is either \(b\) or
a vertex of \(O\setminus\{o\}\).
\end{lemma}

\begin{proof}[Proof sketch]
The first two statements are exactly preservation of the second color.  The
last two use preservation of the complement of the second color, together
with the fact that \(a,b\notin T\) and \(o,o'\in O\).
\end{proof}

The productive cases are the ones where a two-hole match identifies an
endpoint with its counterpart, or gives a chord in one of the type-flow
paths.  For example, if \(\theta(a)=b\), then the match respects the moved
second-colored vertex.  If also \(\theta(b)=a\), then it respects the moved
uncolored vertex as well; the only remaining ambiguity is how to reinsert
the two deleted holes \(t,o\) and \(t',o'\).  This is a two-star error, hence
a smaller centered problem over the common two-hole card.  If one of the
low ports lies in \(T\), then the match says that the low type-flow path can
enter or leave the same card type through an internal \(T\)-vertex; in a
shortest path, this should create a chord.  The complementary ports do the
same for the high path through \(O\).

\begin{conjecture}[Two-hole splice lemma]
Work in a minimal fixed-host singleton counterexample, minimizing
\(|V(K)|\) and then the total length of the chosen low and complementary
type-flow paths.  Let
\[
  \theta:D_{t,o}^a\cong D_{t',o'}^b
\]
be a two-hole match appearing in a balanced residual rectangular matching.
Then at least one of the following holds.

\begin{enumerate}[leftmargin=2em]
  \item The match is endpoint-compatible:
  \[
    \theta(a)=b,\qquad \theta(b)=a.
  \]
  Re-inserting the two deleted holes either gives a direct endpoint
  cancellation or a smaller colored centered state over the two-hole card.

  \item The low ports define a chord in the low type-flow path, so the path
  from \(H_a\) to \(H_b\) can be shortened.

  \item The complementary ports define a chord in the complementary type-flow
  path, so the path from \(C_b^a\) to \(C_a^b\) can be shortened.

  \item The match exposes a smaller fixed-host singleton obstruction inside
  one of the deleted-hole cards, contradicting minimality of \(|V(K)|\).
\end{enumerate}
\end{conjecture}

Equivalently, a minimal unresolved counterexample would have to make the
remaining residual matches \emph{essentially locked}: neither moved vertex
maps to its counterpart, all four ports are internal, and the induced low and
complementary port relations give no chord in either shortest path.

\begin{conjecture}[No essential locked two-hole match]
There is no essential locked two-hole match in a minimal fixed-host singleton
counterexample, after all direct cancellations, smaller-obstruction moves,
path-splicing chords, and exposed unbalanced rectangular terms have been
removed.
\end{conjecture}

This is the most concrete local obstruction currently visible.  A locked
match would have data
\[
  \theta(a)=x_T\in T\setminus\{t'\},
  \qquad
  \theta^{-1}(b)=y_T\in T\setminus\{t\},
\]
\[
  \theta(b)=x_O\in O\setminus\{o'\},
  \qquad
  \theta^{-1}(a)=y_O\in O\setminus\{o\},
\]
with no shortening relation among the corresponding low and complementary
flow vertices.  The color bookkeeping already forces
\[
  a\in S \Longleftrightarrow x_T\in S
  \Longleftrightarrow y_T\in S
  \Longleftrightarrow b\in S
\]
and
\[
  b\in S \Longleftrightarrow x_O\in S
  \Longleftrightarrow y_O\in S
  \Longleftrightarrow a\in S.
\]
Degree and two-step balance should add the same equal-degree and witness
neutrality constraints as in the absorber analysis.  Thus a locked match is
not an arbitrary two-hole isomorphism; it is a highly constrained
four-port near-automorphism of the host with two stars removed.

\begin{remark}[Locked maps are not forbidden absolutely]
The adjective ``essential'' matters.  Locked two-hole isomorphisms can occur
in already solved states.  For example, in the empty six-vertex graph with
\(S=\emptyset\), \(T=\{0,1\}\), \(a=2\), \(b=3\), and \(t=t'=0\),
\(o=o'=4\), the two-hole cards are just colored empty graphs.  A
color-preserving bijection can send \(a\) to the remaining \(T\)-vertex and
\(b\) to an \(O\)-vertex, making all ports internal.  But the singleton
switch is trivially solved by a direct automorphism.  Therefore the target is
not to prohibit locked local maps; it is to show that no locked map can be
the last surviving term in a minimal unresolved rectangular kernel.

The same warning applies even after extension errors are included.  A random
seven-vertex solved singleton instance produced a locked two-hole map with
nonzero low and high extension errors, but with the low ports equal and the
complementary ports equal.  Such a map is not essential: the port loops are
exactly the local splice signal.  This reinforces the intended target.  The
forbidden object is a locked map with nonzero unresolved errors and no
port-loop, no path chord, no direct cancellation, and no smaller obstruction.
\end{remark}

\begin{lemma}[First-order equations for a locked match]
Let \(\theta:D_{t,o}^a\cong D_{t',o'}^b\) be locked, with ports
\[
  \theta(a)=x_T,\qquad \theta(y_T)=b,
  \qquad
  \theta(b)=x_O,\qquad \theta(y_O)=a.
\]
Then the following total-degree equations hold:
\[
  \deg_K(a)-\one_{at}-\one_{ao}
  =
  \deg_K(x_T)-\one_{x_Tt'}-\one_{x_To'},
\]
\[
  \deg_K(y_T)-\one_{y_Tt}-\one_{y_To}
  =
  \deg_K(b)-\one_{bt'}-\one_{bo'},
\]
\[
  \deg_K(b)-\one_{bt}-\one_{bo}
  =
  \deg_K(x_O)-\one_{x_Ot'}-\one_{x_Oo'},
\]
\[
  \deg_K(y_O)-\one_{y_Ot}-\one_{y_Oo}
  =
  \deg_K(a)-\one_{at'}-\one_{ao'}.
\]
The second-color degree equations are
\[
  \deg_T(a)-\one_{at}
  =
  \deg_T(x_T)-\one_{x_Tt'}+\one_{x_Tb},
\]
\[
  \deg_T(y_T)-\one_{y_Tt}+\one_{y_Ta}
  =
  \deg_T(b)-\one_{bt'},
\]
\[
  \deg_T(b)-\one_{bt}+\one_{ba}
  =
  \deg_T(x_O)-\one_{x_Ot'}+\one_{x_Ob},
\]
and
\[
  \deg_T(y_O)-\one_{y_Ot}+\one_{y_Oa}
  =
  \deg_T(a)-\one_{at'}+\one_{ab}.
\]
\end{lemma}

\begin{proof}[Proof sketch]
The total-degree equations are obtained by comparing degrees in the deleted
uncolored hosts \(K-\{t,o\}\) and \(K-\{t',o'\}\).  The second-color
equations compare the number of neighbors of the given vertex inside the
second-color class.  On the domain side the second color is
\((T\setminus\{t\})\cup\{a\}\); on the codomain side it is
\((T\setminus\{t'\})\cup\{b\}\).  Substituting the four port images gives the
four displayed formulas.
\end{proof}

These equations are deliberately elementary.  They are the first Lean-facing
normal form for the locked case: all terms are degrees in \(K\) plus
adjacency indicators among the ten named vertices
\[
  a,b,t,t',o,o',x_T,y_T,x_O,y_O,
\]
with possible coincidences among the ports.  The low-order singleton balance
can then replace \(\deg_K(a)\) by \(\deg_K(b)\) and
\(\deg_T(a)\) by \(\deg_T(b)\).  What remains is a finite system of
\(\{0,1\}\)-corrections.  A locked match is essential only if this system is
compatible while still producing no port equality, no chord, and no smaller
singleton state.

\begin{definition}[Two-hole extension errors]
Let \(\theta:D_{t,o}^a\cong D_{t',o'}^b\).  The low extension error is
\[
  \Err_O(\theta)
  =
  \theta\bigl(N_K(o)\setminus\{t,o\}\bigr)
  \triangle
  \bigl(N_K(o')\setminus\{t',o'\}\bigr),
\]
viewed as a subset of \(V(K)\setminus\{t',o'\}\).  The high extension error
is
\[
  \Err_T(\theta)
  =
  \theta\bigl(N_K(t)\setminus\{t,o\}\bigr)
  \triangle
  \bigl(N_K(t')\setminus\{t',o'\}\bigr).
\]
\end{definition}

The notation reflects which deleted hole is being reinserted.  The low error
measures whether \(\theta\) extends after re-inserting the two \(O\)-holes
\(o,o'\), hence whether it gives an isomorphism between the low one-hole
cards \(A_t\) and \(B_{t'}\).  The high error measures whether \(\theta\)
extends after re-inserting the two \(T\)-holes \(t,t'\), hence whether it
gives an isomorphism between the complementary one-hole cards \(C_o^a\) and
\(C_{o'}^b\).

It is useful to keep signs, not just supports.  For
\(z\in V(K)\setminus\{t,o\}\), define
\[
  \sigma_O^\theta(z)
  =
  \one_{oz\in E(K)}-\one_{o'\theta(z)\in E(K)},
\]
and
\[
  \sigma_T^\theta(z)
  =
  \one_{tz\in E(K)}-\one_{t'\theta(z)\in E(K)}.
\]
Then
\[
  |\Err_O(\theta)|=\sum_z |\sigma_O^\theta(z)|,
  \qquad
  |\Err_T(\theta)|=\sum_z |\sigma_T^\theta(z)|.
\]
The sign records whether the domain reinserted hole has an extra adjacency
\((+1)\) or the codomain reinserted hole has an extra adjacency \((-1)\).

\begin{lemma}[Zero extension error gives a chord]
Let \(\theta:D_{t,o}^a\cong D_{t',o'}^b\).

\begin{enumerate}[leftmargin=2em]
  \item If \(o\in S\Longleftrightarrow o'\in S\) and
  \(\Err_O(\theta)=\emptyset\), then \(\theta\) extends by \(o\mapsto o'\)
  to a colored isomorphism
  \[
    A_t\cong B_{t'}.
  \]
  Thus, if \(A_t\) and \(B_{t'}\) occur as nonconsecutive vertices of the
  chosen low type-flow path, that path has a chord and can be shortened.

  \item If \(t\in S\Longleftrightarrow t'\in S\) and
  \(\Err_T(\theta)=\emptyset\), then \(\theta\) extends by \(t\mapsto t'\)
  to a colored isomorphism
  \[
    C_o^a\cong C_{o'}^b.
  \]
  Thus, if \(C_o^a\) and \(C_{o'}^b\) occur as nonconsecutive vertices of the
  chosen complementary type-flow path, that path has a chord and can be
  shortened.
\end{enumerate}
\end{lemma}

\begin{proof}[Proof sketch]
For the first statement, the only graph adjacencies not already checked by
\(\theta\) are the adjacencies incident to the reinserted \(O\)-hole.  These
are exactly compared by \(\Err_O(\theta)\).  The condition
\(o\in S\Longleftrightarrow o'\in S\) says that the reinserted vertices have
the same first-color status; both are outside the second color.  The second
statement is identical, with \(t,t'\) reinserted instead.  In that case both
reinserted vertices are second-colored, so only first-color compatibility is
additional.
\end{proof}

Therefore an essential locked match must satisfy four negative conditions:
the two endpoint maps \(a\mapsto b\) and \(b\mapsto a\) fail, the low
extension either has nonzero error or wrong first-color status at \(o,o'\),
and the high extension either has nonzero error or wrong first-color status
at \(t,t'\).  The locked case is not just a four-port condition; it is a
simultaneous pair of nonzero one-star extension errors over the same
two-hole card.

\begin{lemma}[Extension-error size tests]
For a two-hole match \(\theta:D_{t,o}^a\cong D_{t',o'}^b\), put
\[
  A_O=N_K(o)\setminus\{t,o\},
  \qquad
  B_O=N_K(o')\setminus\{t',o'\},
\]
and
\[
  A_T=N_K(t)\setminus\{t,o\},
  \qquad
  B_T=N_K(t')\setminus\{t',o'\}.
\]
Then
\[
  |\theta(A_O)|-|B_O|
  =
  \deg_K(o)-\one_{ot\in E(K)}
  -
  \deg_K(o')+\one_{o't'\in E(K)},
\]
and
\[
  |\theta(A_T)|-|B_T|
  =
  \deg_K(t)-\one_{to\in E(K)}
  -
  \deg_K(t')+\one_{t'o'\in E(K)}.
\]
In particular, if the corresponding reinserted one-hole cards have the same
edge count, then the relevant extension error has even cardinality.
\end{lemma}

\begin{proof}[Proof sketch]
The first two displayed identities just count the attachment sets of the
reinserted vertices into the common two-hole card.  The difference between
the edge counts of the low one-hole cards \(A_t\) and \(B_{t'}\) is exactly
the first displayed quantity; the difference between the edge counts of the
complementary one-hole cards \(C_o^a\) and \(C_{o'}^b\) is exactly the
second.  If the two attachment sets have the same cardinality, then
\[
  |A\triangle B|=2|A\setminus B|
\]
after transporting one set by \(\theta\), so the extension error is even.
\end{proof}

Thus a one-edge locked error is immediately visible as an edge-count
imbalance between the corresponding one-hole cards.  In a shortest type-flow
obstruction, such an imbalance must be compensated elsewhere along the same
path.  This suggests a second potential: sum the sizes of the low and high
extension errors over a rectangular matching.  Any productive splice reduces
this potential to zero on one rectangle; an essential locked matching would
have to keep the potential positive while maintaining all edge-count
telescoping identities.

\begin{definition}[Balanced residual rectangular potential]
The rectangular chain
\[
  R_{T,O}
  =
  \sum_{t\in T}\sum_{o\in O}
    \bigl([D_{t,o}^a]-[D_{t,o}^b]\bigr)
\]
need not be the zero chain.  Its nonzero coefficients are exposed
rectangular boundary terms.  After removing all exposed terms that can be
used for a support-escalation squeeze, a terminal-exposure obstruction, or a
smaller obstruction, a remaining balanced residual subchain has, for each
isomorphism class of two-hole cards, equally many positive and negative
occurrences.

On such a balanced residual subchain, choose a bijective matching between
the positive occurrences \(D_{t,o}^a\) and the negative occurrences
\(D_{t',o'}^b\), together with a specific isomorphism \(\theta\) for each
matched pair.  Define
\[
  \Phi
  =
  \sum_{\theta}
    \bigl(|\Err_O(\theta)|+|\Err_T(\theta)|\bigr).
\]
Among all such residual rectangular matchings, choose one minimizing
\(\Phi\).  Subject to that, minimize the number of essential locked matches.
\end{definition}

\begin{remark}[Rigidified vector form]
Fix an isomorphism class \(\mathcal D\) of two-hole cards, and choose a
representative colored card \(R\).  For every positive occurrence
\[
  D_{t,o}^a\cong R,
\]
choose one identification \(\iota:D_{t,o}^a\cong R\).  This occurrence then
determines two attachment vectors on \(V(R)\):
\[
  u_O(z)=\one_{o\,\iota^{-1}(z)\in E(K)},
  \qquad
  u_T(z)=\one_{t\,\iota^{-1}(z)\in E(K)}.
\]
Similarly every negative occurrence \(D_{t',o'}^b\cong R\) determines two
attachment vectors
\[
  v_O(z)=\one_{o'\,\iota^{-1}(z)\in E(K)},
  \qquad
  v_T(z)=\one_{t'\,\iota^{-1}(z)\in E(K)}.
\]
After this rigidification, choosing an isomorphism
\(\theta:D_{t,o}^a\cong D_{t',o'}^b\) is the same as choosing compatible
identifications with \(R\), and its contribution to \(\Phi\) is the Hamming
transport cost
\[
  d(u_O,v_O)+d(u_T,v_T).
\]
Zero cost in the \(O\)-coordinate means the low one-hole cards match.  Zero
cost in the \(T\)-coordinate means the complementary one-hole cards match.
\end{remark}

This vector form clarifies the nature of the remaining problem.  The
rectangular boundary does not merely say that two multisets of card types are
equal.  For each two-hole card type \(R\), it gives two finite multisets of
pairs of binary attachment vectors
\[
  (u_O,u_T)
  \qquad\text{and}\qquad
  (v_O,v_T).
\]
The direct-cancellation campaign first has to deal with exposed unbalanced
two-hole terms.  Only after those terms are removed or converted into smaller
obstructions does one get a balanced residual transport problem.  In that
residual problem, an essential locked rectangular kernel is an optimal
transport plan in which every matched pair differs in both coordinates and
none of those differences gives a path chord.

\begin{remark}[The transport statement is not abstractly true]
The last sentence is not a theorem about arbitrary multisets of binary-vector
pairs.  Abstractly, one can have a positive multiset containing a single pair
\((u_O,u_T)\) and a negative multiset containing a single pair
\((v_O,v_T)\) with \(u_O\neq v_O\) and \(u_T\neq v_T\); then every matching
has positive cost in both coordinates.  Therefore the missing argument must
use the extra structure: these vectors arise as the two deleted-hole
attachment vectors inside one fixed graph, they sit in the low and
complementary type-flow paths, and their aggregate comes from equality of an
actual one-vertex deletion deck.  This is a useful guardrail for the Lean
formalization: the minimum-error lemma cannot be stated as a standalone
transport theorem, and it only applies after the exposed rectangular boundary
terms have been handled.
\end{remark}

\begin{definition}[Active rectangle of two shortest paths]
Fix a shortest positive low type-flow path
\[
  H_a=Y_0\to Y_1\to\cdots\to Y_r=H_b,
\]
where the edge \(Y_{i-1}\to Y_i\) is indexed by \(t_i\in T\).  Also fix a
shortest positive complementary type-flow path
\[
  C_b^a=Z_0\to Z_1\to\cdots\to Z_s=C_a^b,
\]
where the edge \(Z_{j-1}\to Z_j\) is indexed by \(o_j\in O\).  The active
rectangle is the multiset of pairs
\[
  \{1,\ldots,r\}\times\{1,\ldots,s\},
\]
whose \((i,j)\)-entry is the two-hole pair
\[
  D_{t_i,o_j}^a,\qquad D_{t_i,o_j}^b.
\]
\end{definition}

\begin{lemma}[Active rectangular identity]
For the two chosen paths,
\[
  \sum_{i=1}^r\sum_{j=1}^s
    \bigl([D_{t_i,o_j}^a]-[D_{t_i,o_j}^b]\bigr)
\]
is obtained both by deleting the complementary path holes \(o_j\) from the
low path telescoping identity and by deleting the low path holes \(t_i\) from
the complementary path telescoping identity.
\end{lemma}

\begin{proof}[Proof sketch]
The low path gives the chain identity
\[
  [H_b]-[H_a]=\sum_{i=1}^r([A_{t_i}]-[B_{t_i}]).
\]
Deleting each \(o_j\) and summing over \(j\) gives the displayed double sum.
The complementary path gives
\[
  [C_a^b]-[C_b^a]=\sum_{j=1}^s([C_{o_j}^a]-[C_{o_j}^b]).
\]
Deleting each \(t_i\) and summing over \(i\) gives the same double sum, since
\[
  A_{t_i}-o_j=C_{o_j}^a-t_i=D_{t_i,o_j}^a,
  \qquad
  B_{t_i}-o_j=C_{o_j}^b-t_i=D_{t_i,o_j}^b.
\]
\end{proof}

This active rectangle is the right object for the splice lemma, but it has
the same caveat as the full rectangle: it need not be balanced as a
positive-negative chain.  If it has an exposed unbalanced two-hole type, that
type is already a candidate separator.  If an active balanced residual
remains, then a zero \(O\)-extension error in a matched entry gives an
isomorphism
\[
  A_{t_i}\cong B_{t_{i'}}
\]
for some active \(i,i'\), hence a chord or endpoint cancellation in the low
path.  A zero \(T\)-extension error similarly gives a chord or endpoint
cancellation in the complementary path.  Rectangular matches outside the
active rectangle may be harmless side cycles; they should not be part of the
minimum-error obstruction.

\begin{definition}[Exposed and residual active terms]
For a colored two-hole isomorphism class \(Y\), define its active rectangular
coefficient by
\[
  c_{\rm act}(Y)
  =
  \#\{(i,j):D_{t_i,o_j}^a\cong Y\}
  -
  \#\{(i,j):D_{t_i,o_j}^b\cong Y\}.
\]
If \(c_{\rm act}(Y)\neq 0\), call \(Y\) an exposed active type.  The balanced
active residual is what remains after all exposed active types have either
been converted into a size drop, a path shortening, an off-path side
absorption, or a terminal-exposure obstruction.
\end{definition}

\begin{definition}[Active isolation]
The full rectangular coefficient of a two-hole type \(Y\) is
\[
  c_{\rm full}(Y)
  =
  \#\{(t,o)\in T\times O:D_{t,o}^a\cong Y\}
  -
  \#\{(t,o)\in T\times O:D_{t,o}^b\cong Y\}.
\]
An exposed active type is actively isolated if its active imbalance cannot be
cancelled by occurrences outside the chosen shortest low and complementary
paths without creating a shorter path, a direct endpoint cancellation, or a
smaller fixed-host obstruction.
\end{definition}

\begin{conjecture}[Active isolation lemma]
In a minimal fixed-host singleton counterexample with shortest chosen paths,
every exposed active type is either productive, in the sense above, or actively
isolated.  Equivalently, an inactive occurrence that cancels an active
two-hole imbalance is not harmless: matching the active occurrence to that
inactive occurrence gives a path replacement, a removable side cycle, a direct
endpoint cancellation, or a smaller obstruction.
\end{conjecture}

\begin{remark}[Why active isolation is a real extra issue]
The active rectangle is chosen from two shortest paths; it is not the whole
rectangle \(T\times O\).  Therefore an active coefficient \(c_{\rm act}(Y)\)
can be nonzero even when the full coefficient \(c_{\rm full}(Y)\) is zero,
because the compensating terms live off the two paths.  The proof cannot
simply declare an exposed active type to be a deck-count contradiction.  It
must first show that off-path compensation either shortens the chosen paths or
is a side cycle that can be removed from the obstruction.  If this active
isolation lemma is false, the minimal bad certificate must be enlarged by a
neutral side absorber.
\end{remark}

\begin{definition}[Support saturation in a two-hole card]
For an occurrence inside \(D_{t,o}^a\) or \(D_{t,o}^b\), call the support
\(O\)-saturated if it contains every vertex of \(O\setminus\{o\}\), and
\(T\)-saturated if it contains every vertex of \(T\setminus\{t\}\).  Call it
terminal if it is both \(O\)- and \(T\)-saturated and also contains the moved
vertices \(a\) and \(b\), so that it is the whole two-hole card.
\end{definition}

\begin{conjecture}[Support escalation]
Let \(Y\) be an actively isolated exposed active type, and let \(P\) be a
lowest-complexity separator witnessing its exposure.  If an occurrence of \(P\)
is not \(O\)-saturated, or is not \(T\)-saturated, then a further deletion in
the missing side produces a productive two-hole splice: a direct cancellation,
a path chord, or a smaller fixed-host obstruction.  Hence a nonproductive
actively isolated exposure must be terminal.
\end{conjecture}

\begin{remark}[The terminal exposure fork]
This is the sharpest newly visible fork.  Nonterminal exposure should be
productive because a missing vertex gives another deletion direction in which
the same discrepancy is still visible.  Terminal exposure is different: it is
not a low-order pattern-count failure but an imbalance of entire two-hole card
types along the active rectangle.  If terminal actively isolated exposure can
occur without a path shortening, then the rectangular campaign needs a new
global ingredient.  If terminal exposure cannot occur, the only remaining
case is the balanced residual matching problem.
\end{remark}

\begin{remark}[The tempting false shortcut]
There is an attractive but invalid one-line proof at this point.  The low
slice equality says
\[
  [H_a]+\sum_{t\in T}[A_t]
  =
  [H_b]+\sum_{t\in T}[B_t].
\]
If one could apply ``delete every uncolored vertex of the displayed fixed
host'' as a well-defined operator on abstract colored card types, then applying
this operation to the low slice would seem to force the rectangular chain to
vanish.  That would give direct cancellation immediately.

The flaw is exactly the obstruction we are studying.  The operation
``delete the vertices of \(O\)'' is not intrinsic to the colored card type.
After a card isomorphism, the image of the fixed-host set \(O\) may be a
different uncolored subset.  Thus named partial deletion is not a functor on
the abstract deck; it is only available on the chosen fixed-host
representatives.  The active-isolation and support-escalation lemmas are
attempts to recover enough functoriality from shortest-path minimality.  This
is why the rectangular side street is promising rather than cosmetic: it
pinpoints the precise missing naturality statement.
\end{remark}

\begin{remark}[Strategic assessment]
Subjectively, this is a good side street, but not yet the main road.  It is
good because it has reduced the vague question ``why should the deck choose
the right card isomorphisms?'' to a finite naturality problem for named
partial deletions.  It is not yet the main road because active isolation could
fail: off-path terms might absorb every exposed imbalance without shortening
the chosen paths.  If that happens, the current campaign needs a new global
minimality or flow argument.

The best-case outcome is that shortest-path minimality restores named
partial-deletion functoriality: any non-natural card isomorphism creates a
chord, side-cycle removal, or smaller obstruction.  Then terminal exposure
disappears, the residual matching is balanced, and the locked-match correction
cycle becomes the only local lemma left.  The worst-case useful outcome is a
neutral side absorber, which would be a concrete finite obstruction to this
proof strategy and a clear computational target.
\end{remark}

\begin{conjecture}[Shortest-path naturality for named deletion]
Let a two-hole card isomorphism
\[
  \theta:D_{t_i,o_j}^a\cong D_{t',o'}^b
\]
match an active occurrence to another occurrence of the same two-hole type.
If \(\theta\) fails to transport the named \(T/O\)-sides coherently, then one
of the port or extension-error mechanisms produces a shorter low path, a
shorter complementary path, a direct endpoint cancellation, or a smaller
fixed-host singleton obstruction.  If \(\theta\) does transport the named
sides coherently, then the two occurrences belong to the same removable side
cycle and do not contribute to a minimal active obstruction.
\end{conjecture}

\begin{definition}[Name-defect type]
Let
\[
  \theta:D_{t,o}^a\cong D_{t',o'}^b
\]
be a two-hole match.  Say \(\theta\) has a low name defect if
\[
  \theta(a)\neq b,
\]
equivalently if \(\theta^{-1}(b)\neq a\).  Say \(\theta\) has a complementary
name defect if
\[
  \theta(b)\neq a,
\]
equivalently if \(\theta^{-1}(a)\neq b\).  Thus \(\theta\) is:
\begin{enumerate}[leftmargin=2em]
  \item name-coherent if it has no name defect;
  \item purely low-defective if it has a low name defect but no complementary
  name defect;
  \item purely complementary-defective if it has a complementary name defect
  but no low name defect;
  \item mixed-defective if it has both defects.
\end{enumerate}
\end{definition}

\begin{lemma}[Name coherence means side coherence]
If \(\theta:D_{t,o}^a\cong D_{t',o'}^b\) is name-coherent, then
\[
  \theta(T\setminus\{t\})=T\setminus\{t'\},
  \qquad
  \theta(O\setminus\{o\})=O\setminus\{o'\},
\]
and \(\theta(a)=b\), \(\theta(b)=a\).  Hence \(\theta\) transports every named
side visible in the two-hole card correctly.
\end{lemma}

\begin{proof}[Proof sketch]
The second color in the domain is \((T\setminus\{t\})\cup\{a\}\), and the
second color in the codomain is \((T\setminus\{t'\})\cup\{b\}\).  If
\(\theta(a)=b\), preservation of the second color forces the remaining
second-colored vertices to map exactly to \(T\setminus\{t'\}\).  The complement
of the second color is \((O\setminus\{o\})\cup\{b\}\) in the domain and
\((O\setminus\{o'\})\cup\{a\}\) in the codomain.  If \(\theta(b)=a\), the same
argument gives the \(O\)-side equality.
\end{proof}

\begin{conjecture}[Name-defect case split]
In a minimal fixed-host singleton counterexample with shortest active low and
complementary paths, a two-hole match used to compensate an exposed active
term falls into the following cases.

\begin{enumerate}[leftmargin=2em]
  \item Name-coherent matches are removable side-cycle matches unless
  reinserting \(t,o\) and \(t',o'\) exposes a smaller fixed-host obstruction.

  \item Purely low-defective matches produce a chord or endpoint cancellation
  in the low type-flow path, unless their low extension error is part of a
  smaller fixed-host obstruction.

  \item Purely complementary-defective matches produce a chord or endpoint
  cancellation in the complementary type-flow path, unless their complementary
  extension error is part of a smaller fixed-host obstruction.

  \item Mixed-defective matches are precisely the candidates for essential
  locked residual matches.  They must be handled by the extension-error
  potential \(\Phi\) and the locked correction-cycle argument.
\end{enumerate}
\end{conjecture}

\begin{remark}[Why this split is encouraging]
The naturality problem now has only one genuinely hard branch.  Zero defect is
not an obstruction to naturality; it is naturality.  One-sided defect points in
one type-flow direction and should shorten that path.  The two-sided defect is
the locked case already isolated by the port and extension-error analysis.
Thus the side street is feeding back into the main proof rather than spawning
an unrelated new problem.
\end{remark}

\begin{remark}[One-sided defects have an extension fork]
The previous case split hides one important condition.  A purely low-defective
match gives an actual low path chord only when the \(O\)-extension error
vanishes, and a purely complementary-defective match gives an actual
complementary chord only when the \(T\)-extension error vanishes.  If the
relevant extension error is nonzero, the match does not yet contradict
shortest-path minimality.  Instead it creates a smaller one-star attachment
problem over the common two-hole card.

For a purely low-defective match
\[
  \theta:D_{t,o}^a\cong D_{t',o'}^b
\]
with \(\theta(b)=a\), the two possible low extensions compare the attachment
sets of \(o\) and \(o'\) inside the common two-hole card.  If
\(\Err_O(\theta)=\emptyset\), then \(A_t\cong B_{t'}\) and the low path has a
chord or endpoint cancellation.  If \(\Err_O(\theta)\neq\emptyset\), then the
remaining data are exactly a one-star defect:
\[
  \theta(N_K(o)\setminus\{t,o\})
  \neq
  N_K(o')\setminus\{t',o'\}.
\]
This is smaller in vertex count, but it is not automatically a smaller
counterexample unless the surrounding rectangular identities imply the
appropriate punctured-deck equality for this one-star problem.
\end{remark}

\begin{conjecture}[One-sided defect inheritance]
In a minimal fixed-host singleton counterexample, let a purely one-sided
two-hole match appear in the active rectangle.  If its corresponding extension
error is nonzero, then the common two-hole card together with the two
reinserted-hole attachment sets inherits a smaller centered-extension
counterexample.  Equivalently, the nonzero extension error cannot be an
isolated local mismatch; it carries enough deck equality from the rectangular
telescoping identities to be handled by the outer size induction.
\end{conjecture}

\begin{remark}[Inheritance is not local]
This conjecture is deliberately not a statement about an arbitrary pair of
one-vertex extensions of an arbitrary two-hole card.  Locally, it is easy to
take a colored card \(R\), attach two new vertices \(o,o'\) to different
subsets of \(V(R)\), and arrange a nonzero extension error.  There is no
reason for those two one-star extensions to have equal punctured decks.

The missing information must come from the fact that the two one-star
extensions sit on a shortest type-flow path inside a global fixed-host
singleton state.  Thus the inheritance problem is really a path-localization
problem: can the segment of the low or complementary path between the two
one-hole cards be pushed down through the common two-hole deletion?
\end{remark}

\begin{conjecture}[One-sided path localization]
Suppose a pure low-defect match
\[
  \theta:D_{t,o}^a\cong D_{t',o'}^b
\]
has \(\Err_O(\theta)\neq\emptyset\), and suppose the corresponding low cards
\[
  A_t,\qquad B_{t'}
\]
occur on the chosen shortest low type-flow path.  Then the low path segment
between \(A_t\) and \(B_{t'}\), after deleting the \(O\)-holes seen by that
segment, induces equality of the punctured extension decks of the two
one-star extensions of the common two-hole card.  Consequently the nonzero
\(O\)-extension error is a smaller centered-extension counterexample.

The dual statement holds for pure complementary defects with nonzero
\(T\)-extension error.
\end{conjecture}

\begin{remark}[Path localization only gives a partial deck]
The previous conjecture is still slightly too optimistic.  A low path segment
between \(A_t\) and \(B_{t'}\) does not by itself compare every card of the
two one-star extensions over the common two-hole card
\[
  R \cong D_{t,o}^a\cong D_{t',o'}^b.
\]
It compares only those deletions that are visible along the segment.  The
active rectangle supplies additional deletions by the complementary holes
appearing in the chosen complementary path.  Together, a low segment and a
complementary segment form a rectangular strip.  That strip may still miss
some vertices of \(R\).

Thus the honest inheritance target has two parts.  First, the path strip gives
equality of a \emph{partial} punctured deck for the two one-star extensions of
\(R\).  Second, minimality or support escalation must show that the missing
deleted vertices are irrelevant, productive, or terminal.  Only after this
cover/completion step do we get a smaller centered-extension counterexample.
\end{remark}

\begin{definition}[Localization strip]
For a pure low-defect match at an active rectangle entry, choose the shortest
subpath of the low type-flow path joining the one-hole cards \(A_t\) and
\(B_{t'}\).  Pair this subpath with the complementary indices used in the
active rectangle.  The resulting set of two-hole deletions is the localization
strip.  Its visible vertices are the vertices of the common two-hole card
whose deletion appears in one of these strip terms.
\end{definition}

\begin{conjecture}[Partial localization]
The localization strip of a pure one-sided match with nonzero extension error
induces equality of the partial punctured decks of the two one-star extensions
over the common two-hole card, restricted to the visible vertices of the
strip.
\end{conjecture}

\begin{conjecture}[Localization cover/completion]
In a minimal fixed-host singleton counterexample, every invisible vertex of a
localization strip is productive or terminal.  More explicitly, if deleting an
invisible vertex from the two one-star extensions changes the partial deck
comparison, then the resulting exposed term gives a path shortening, direct
cancellation, or smaller obstruction.  If no invisible deletion is productive,
then the one-sided defect is supported on all invisible vertices and becomes a
terminal-exposure obstruction.

Consequently, a nonterminal pure one-sided nonzero extension error inherits a
full smaller centered-extension counterexample.
\end{conjecture}

\begin{remark}[Sharper success or failure]
The one-sided obstruction has now split into two testable claims.  Partial
localization is mostly bookkeeping: it should follow by restricting the
telescoping identity to the chosen strip.  Cover/completion is the real
mathematical point.  If it fails, the failure is a \emph{coverage gap}: a
shortest path strip sees a nonzero one-star extension error but misses exactly
the deletions needed to certify punctured-deck equality downstairs.  This is a
more concrete obstruction than the earlier vague ``inheritance failure''.
\end{remark}

\begin{remark}[Concrete success or failure]
This is now the most concrete near-term test of the side street.  If
one-sided path localization is true, then pure one-sided defects are removed
by the same size induction that removes direct recentering errors.  If it is
false, the obstruction is very explicit: a shortest path segment whose
endpoints share a two-hole deletion but whose induced one-star attachments do
not inherit deck equality.  That would not kill the reconstruction conjecture,
but it would kill the current ``shortest-path naturality'' campaign unless a
new global flow invariant is added.
\end{remark}

\begin{remark}[A concrete failure mode]
If the one-sided defect inheritance conjecture is false, the campaign fails in
a very specific way.  There would be a pure low or pure complementary
name-defect match with nonzero extension error, but the induced one-star
attachment problem over the two-hole card would not have equal punctured decks.
Such a match would not be a locked residual match, because one pair of moved
vertices is correctly aligned; it would also not be a chord, because the
extension error is nonzero.  This is a fourth computational falsification
target, and it is arguably easier to search for than a full mixed locked
certificate.
\end{remark}

\begin{remark}[Why this would settle exposed terms]
This is the active-isolation lemma in its most conceptual form.  The whole
problem is that named deletion does not descend to abstract card types.  The
conjecture says that in a minimal shortest-path counterexample, every failure
of descent is productive, and every nonproductive descent failure is actually
a removable side cycle.  If this holds, then exposed active coefficients can
be read as genuine fixed-host imbalances rather than artifacts of a bad choice
of path representatives.
\end{remark}

\begin{conjecture}[Exposed term dichotomy]
In a minimal fixed-host singleton counterexample with shortest low and
complementary type-flow paths, every exposed active type \(Y\) has one of the
following two outcomes.

\begin{enumerate}[leftmargin=2em]
  \item It is productive: an occurrence of \(Y\) gives a direct endpoint
  cancellation, a smaller fixed-host singleton obstruction, or a chord in one
  of the two chosen type-flow paths.

  \item It is terminal: after active isolation and support escalation, every
  lowest-complexity separator realizing the coefficient \(c_{\rm act}(Y)\) is
  the whole two-hole card on the remaining vertices.
\end{enumerate}

The second outcome is now the only exposed-term failure mode.  It says that
all lower-support attempts to see the imbalance have become productive, and
the remaining imbalance is a pure two-hole card-type imbalance along the
active rectangle.
\end{conjecture}

\begin{remark}[How to try to prove the exposed dichotomy]
Assume, for instance, that \(c_{\rm act}(Y)>0\).  Choose an occurrence
\[
  D_{t_i,o_j}^a\cong Y
\]
and a smallest colored pattern \(P\) whose count contributes to the positive
exposure.  If \(P\) misses some vertex \(o\in O\), then the same pattern is
visible after deleting \(o\) from the low path telescoping identity.  The
active rectangular identity says that this visible excess is also visible
after deleting a \(T\)-vertex from the complementary path.  Unwinding the
corresponding two-hole isomorphism gives exactly the port configuration of the
two-hole splice lemma.

Thus the attempt has only two possible exits.  Either the port configuration
is productive, contradicting minimality of the paths or of \(|V(K)|\), or the
pattern could not have missed any remaining vertex on either side of the
two-hole rectangle.  Repeating the same argument from both directions forces
terminal support.  At that point the obstruction is no longer a pattern-count
separator; it is an imbalance of entire two-hole card types in the active
rectangle.  That terminal case must either be killed by active isolation
against the full rectangle, or carried forward as part of the minimal bad
certificate.
\end{remark}

\begin{conjecture}[Minimum-error residual matching has no essential lock]
In a minimal fixed-host singleton counterexample, after all exposed
rectangular boundary terms have been converted into productive moves,
terminal-exposure obstructions, or smaller obstructions, a minimum-error
matching of the balanced residual rectangular chain contains no essential
locked two-hole match.
\end{conjecture}

This is a more operational version of the two-hole splice lemma.  If a
locked match has nonzero low or high extension error, then some adjacency
incident to one of the deleted holes is mismatched.  In the balanced residual
chain, the same mismatched adjacency should be corrected by another occurrence
of the same two-hole card type.  Swapping the partners along the resulting
alternating correction cycle should either reduce \(\Phi\), create a
zero-error extension in one direction, or expose a smaller singleton
obstruction.  In a \(\Phi\)-minimal residual matching, none of these can
happen, so an essential locked match should be impossible.

\begin{conjecture}[Locked correction cycle]
Let a minimum-error matching of the balanced residual rectangular chain
contain an essential locked match \(\theta\).  Choose a signed low or high
error, say
\(\sigma_O^\theta(z)=+1\).  Then the rectangular balance supplies a finite
alternating sequence of matched two-hole cards carrying the same named
adjacency discrepancy with signs alternating \(+1,-1,+1,\ldots,-1\).  Swapping
partners around this correction cycle strictly decreases
\[
  \sum_\theta\sum_z
    \bigl(|\sigma_O^\theta(z)|+|\sigma_T^\theta(z)|\bigr),
\]
unless one of the swaps creates a zero-extension chord or a smaller
singleton obstruction.
\end{conjecture}

This conjecture is the combinatorial heart of the minimum-error approach.
It is the rectangular analogue of cancelling opposite signs in an optimal
matching problem.  The nontrivial point is that the signs are attached to
named adjacencies involving the deleted holes, while the two-hole cards are
unlabelled up to isomorphism.  The two-hole isomorphisms must transport the
names coherently enough for the alternating cycle to be defined.

\begin{definition}[Minimal bad certificate]
A minimal bad certificate for the rectangular campaign consists of the
following finite data.

\begin{enumerate}[leftmargin=2em]
  \item A fixed-host singleton state
  \[
    (K,S,T,a,b)
  \]
  with equal colored deletion decks, no automorphism of \((K,S)\) carrying
  \(T\cup\{a\}\) to \(T\cup\{b\}\), and no direct low or complementary endpoint
  cancellation.

  \item Shortest nontrivial low and complementary type-flow paths
  \[
    H_a=Y_0\to\cdots\to Y_r=H_b,
    \qquad
    C_b^a=Z_0\to\cdots\to Z_s=C_a^b.
  \]

  \item An active rectangular chain
  \[
    R_{\rm act}
    =
    \sum_{i=1}^r\sum_{j=1}^s
      \bigl([D_{t_i,o_j}^a]-[D_{t_i,o_j}^b]\bigr)
  \]
  whose exposed active types are all either actively isolated and terminal in
  the sense of the exposed term dichotomy, or are cancelled by neutral
  off-path side absorbers that do not shorten either path.

  \item A balanced residual subchain of \(R_{\rm act}\), together with a
  typewise matching and specific two-hole isomorphisms minimizing
  \[
    (\Phi,\#\text{essential locked matches}).
  \]

  \item At least one essential locked match in this minimum, and every
  correction cycle based at such a lock is neutral: partner swaps neither
  lower \(\Phi\), nor create a zero-extension chord, nor expose a smaller
  fixed-host obstruction.
\end{enumerate}
\end{definition}

\begin{remark}[Why this certificate is the right failure target]
This is now a finite obstruction, not a philosophical gap.  Proving that no
minimal bad certificate exists proves the fixed-host singleton theorem.  On
the other hand, constructing such a certificate would give a precise reason
the rectangular campaign fails: either active exposure really can evade
path-shortening and remain as terminal two-hole exposure, or the balanced
residual really can support an essential locked minimum-error transport.  Both
failure modes are concrete enough to search for computationally and to
formalize in Lean.
\end{remark}

\begin{remark}[Computational falsification targets]
A useful search should no longer ask only whether equal fixed-host singleton
decks are orbit-successful.  The small cases already suggest that they are.
The sharper search is for one of four internal witnesses, with every
two-hole match classified by name-defect type.

\begin{enumerate}[leftmargin=2em]
  \item A neutral side absorber: an active exposed two-hole type whose
  cancelling occurrences lie off the chosen shortest paths, but where the
  induced off-path match gives no shorter path and no removable side cycle.

  \item Terminal exposure: an actively isolated exposed type for which every
  nonterminal support witness is productive, but the whole two-hole card type
  remains imbalanced along the active rectangle.

  \item A one-sided inheritance failure: a pure low or pure complementary
  name-defect match with nonzero corresponding extension error, where the
  induced one-star attachment problem over the common two-hole card does not
  inherit equal punctured decks.  The sharpened version is a coverage gap:
  the localization strip gives only a partial deck and the invisible deletions
  are neither productive nor terminal.

  \item An essential locked residual match: after deleting exposed terms and
  choosing a minimum-error residual matching, a locked two-hole isomorphism
  survives with nonzero low and high extension errors and no correction cycle
  reducing \((\Phi,\#\text{locks})\).
\end{enumerate}

Finding any one of these would not necessarily disprove reconstruction, but it
would disprove the current proof campaign as stated.  Failing to find any of
them in exhaustive fixed-host singleton searches is evidence for exactly the
lemmas we need to formalize next.
\end{remark}

\begin{remark}[First fixed-host singleton probe]
The script
\[
  \texttt{research/computational/data/fixed\_host\_singleton\_search.py}
\]
exhaustively checks the fixed-host singleton states for \(n\leq 5\), using a
tiny brute-force colored canonical form.  The \(n=5\) run examined \(174080\)
singleton states over the \(34\) unlabeled host graphs.  It found \(24832\)
equal-deck singleton states, no orbit failures, and in every one of these
states both direct low cancellation and direct complementary cancellation held.

The same run is more informative locally.  Among the two-hole matches inside
equal-deck states, it found
\[
  7392
\]
pure low-defective matches and \(7392\) pure complementary-defective matches.
Of these, \(1920\) in each direction had nonzero corresponding extension
error.  Thus pure one-sided nonzero extension errors are real local artifacts
even in solved states.  They cannot be forbidden locally.  This supports the
new formulation: a one-sided nonzero extension error must either be irrelevant
because the state is already solved, or inherit a smaller centered-extension
counterexample from the surrounding path/deck structure.

No mixed-defective matches occurred in this exhaustive \(n=5\) equal-deck
probe.  A bounded labelled \(n=6\) smoke test over \(100000\) singleton states
again found no orbit failures; every equal-deck state in the sample had both
direct endpoint cancellations.  It did find \(21398\) mixed-defective local
matches, \(624\) of them with nonzero low and high extension errors.  Thus
mixed locked-looking maps also cannot be forbidden locally.  They must be
excluded only after exposure
handling, residual balancing, and \(\Phi\)-minimal partner choice.

The cached Python canonicalizer is adequate for small bounded probes, but
exhaustive \(n=6\) and beyond should use nauty or a stronger
colored-canonical-label backend.
\end{remark}

\begin{remark}[Direct-cancellation probe through six vertices]
Using NetworkX's graph atlas to supply the \(156\) unlabeled six-vertex host
graphs, the same script can skip local two-hole classification and test the
upstream direct-cancellation prediction exhaustively for \(n=6\).  The run
\[
  \texttt{fixed\_host\_singleton\_search.py 6 --atlas --skip-local}
\]
examined \(4792320\) fixed-host singleton states and found \(406528\)
equal-deck singleton states.  It found no orbit failures and, more
importantly for the rectangular campaign, every equal-deck state had both
direct low cancellation and direct complementary cancellation:
\[
  H_a\cong H_b
  \qquad\text{and}\qquad
  C_b^a\cong C_a^b.
\]

This is the strongest current computational evidence for double direct
cancellation in small hosts, but the cyclic \(n=9\) example above shows that
double direct cancellation is not a general theorem.  The useful surviving
lesson is that the small equal-deck states were all orbit-successful.

A separate one-slice probe tests the orbit prediction directly, without
assuming equality of the whole fixed-host deck.  Through \(n=5\), four counts
agree at \(24832\): full deck equality, low-slice equality, complementary-slice
equality, and fixed-host orbit success.  The exhaustive six-vertex atlas run
gives the same common count \(406528\).  So far there is no computational
evidence for a one-slice orbit failure.
\end{remark}

\begin{remark}[Why this is Lean-shaped]
The potential \(\Phi\) is finite and purely combinatorial.  A Lean proof would
not need to reason about arbitrary choices forever: choose a minimizing
residual rectangular matching by finite well-ordering, then prove that any locked
match admits a local partner-swap decreasing the lexicographic potential
\[
  (\Phi,\ \#\text{essential locks}).
\]
This is exactly the same architecture as the earlier minimal-defect strategy,
but now the moving part is a finite matching of two-hole cards rather than an
unrooted deck matching.
\end{remark}

\begin{remark}[Next local attack]
The immediate task is now to prove one-slice orbit reconstruction.  The cleanest
route is:
\[
  \text{low-slice equality}
  \quad\Longrightarrow\quad
  T\cup\{a\}\sim_{\operatorname{Aut}(K,S)}T\cup\{b\}.
\]
The first pass should use the restricted Kelly principle for marked-vertex
decks to show that any positive type-flow path transports only the hidden
extension profile of the deleted marked vertex.  Then one tries to prove that a
shortest transport path integrates to an automorphism of the fixed host, or
recovers a smaller centered-extension obstruction.

If this orbit-reconstruction route stalls, the backup is the rectangular locked
residual: use first-color status, total degree, degree into \(T\), degree into
\(O\), and the signed witness charge \(\delta\) to force a loop or chord in the
two-hole port relation, then run the \(\Phi\)-minimal partner-swap descent.
\end{remark}

\begin{remark}[Lean readiness of the fixed-host block]
The following parts of the fixed-host singleton campaign are now
formalization-ready definitions or straightforward lemmas.

\begin{enumerate}[leftmargin=2em]
  \item The fixed-host colored deletion deck \(\Deck_{K,S}(U)\).
  \item Low and complementary color-size slice decks.
  \item The restricted Kelly counting lemma for deleting only marked vertices.
  \item The low and complementary type-flow decompositions.
  \item The rectangular second-boundary identity.
  \item Two-hole ports and the color-bookkeeping lemma.
  \item The first-order locked-match equations.
  \item The extension-error sets \(\Err_O,\Err_T\), their signed versions
  \(\sigma_O,\sigma_T\), and the zero-extension-error chord lemma.
  \item The rigidified vector form for each two-hole card type, and the
  finite potential \(\Phi\) for a rectangular matching.
\end{enumerate}

The part that is not yet Lean-ready as a proof is orbit reconstruction: from a
shortest positive type-flow path, construct either an automorphism carrying
\(T\cup\{a\}\) to \(T\cup\{b\}\), a removable cycle, or a smaller
centered-extension obstruction.  If that fails, the older correction-cycle
descent for \(\Phi\)-minimal rectangular matchings is the fallback local proof
burden.
\end{remark}

\begin{proposition}[Rectangular kernel exclusion solves singleton]
Assume the colored centered extension theorem for colored centers with fewer
than \(|V(K)|\) vertices.  Assume also the no nonzero rectangular kernel
statement for the fixed-host singleton state \((K,S,T,a,b)\).  Then the
fixed-host singleton switch is solved.
\end{proposition}

\begin{proof}[Proof sketch]
The rectangular second-boundary identity holds for every equal-deck
singleton state.  By no nonzero rectangular kernel, at least one endpoint
cancellation holds:
\[
  H_a\cong H_b
  \qquad\text{or}\qquad
  C_b^a\cong C_a^b.
\]
In the first case, the direct-edge size-drop lemma applies.  In the second
case, the complementary direct-cancellation size-drop lemma applies.  Either
way, the problem recenters over a host with one fewer vertex, and the colored
centered induction hypothesis gives the required color-preserving
automorphism of \(K\) carrying \(T\cup\{a\}\) to \(T\cup\{b\}\).
\end{proof}

So the fixed-host singleton problem can now be attacked without mentioning
type-flow paths at all.  Type flows remain useful for finding a minimal
obstruction, but the rectangular kernel is the algebraic object that must be
proved empty.  This is a meaningful simplification: instead of showing that
no routing through \(T\) and no routing through \(O\) can coexist, it is
enough to show that the common \(T\times O\) two-hole boundary separates one
of the two endpoint differences.

\subsection{The second-order boundary issue}

There is a tempting but invalid shortcut.  Each type-flow edge
\[
  B_t\to A_t
\]
compares two singleton colorings of the smaller fixed host \(K-t\):
\[
  B_t=(K-t,S\setminus\{t\},(T\setminus\{t\})\cup\{b\}),
\]
\[
  A_t=(K-t,S\setminus\{t\},(T\setminus\{t\})\cup\{a\}).
\]
If these two smaller colorings had equal fixed-host decks, the induction
hypothesis would solve the edge.  But original deck equality does not give
that equality edge-by-edge.  It gives only the aggregate balance of all such
edges, plus the two endpoint shadows \(H_a,H_b\).

Thus a proof cannot simply apply induction to every type-flow edge.  The
failure of a single edge to have equal smaller decks is a second-order
boundary.  When one deletes another vertex \(w\neq t\), the discrepancy
between \(A_t-w\) and \(B_t-w\) must be compensated elsewhere in the global
second-deletion data.

This suggests a hierarchy.  Let \(\mathcal G_m\) be the free abelian group on
isomorphism classes of finite colored graphs with \(m\) vertices, and define
the deletion boundary
\[
  \partial[X]=\sum_{v\in V(X)}[X-v].
\]
The fixed-host singleton equality says that, for the two colored full hosts
\[
  X_a=(K,S,T\cup\{a\}),
  \qquad
  X_b=(K,S,T\cup\{b\}),
\]
we have
\[
  \partial[X_a]=\partial[X_b].
\]
Equivalently,
\[
  \partial([X_a]-[X_b])=0.
\]
The direct-shadow theorem asks for one specific cancellation inside this
boundary: the term obtained by deleting \(a\) on the left should match the
term obtained by deleting \(b\) on the right.

\begin{remark}[Why this is not a free lunch]
For arbitrary colored graphs, injectivity of \(\partial\) is just the colored
reconstruction problem.  So the boundary language does not by itself prove
anything.  Its value here is that \([X_a]-[X_b]\) is a very special
one-vertex color difference inside a fixed host.  The hope is to prove
acyclicity for this special class of chains, not for all graph chains at
once.
\end{remark}

\begin{conjecture}[Singleton boundary acyclicity]
Let \(X_a=(K,S,T\cup\{a\})\) and \(X_b=(K,S,T\cup\{b\})\) be a fixed-host
singleton pair.  If
\[
  \partial([X_a]-[X_b])=0,
\]
then the natural boundary terms cancel:
\[
  [K-a,S\setminus\{a\},T]
  =
  [K-b,S\setminus\{b\},T].
\]
\end{conjecture}

This is exactly the fixed-host direct-shadow theorem, rewritten as a special
acyclicity statement.  A shortest positive type-flow path is a candidate
nontrivial cycle in this singleton boundary complex.  The absorber exchange,
witness-charge, and neutral-reversibility constraints are attempts to prove
that such a special cycle cannot exist.

\begin{conjecture}[Higher-order path shortening]
If a shortest positive singleton type-flow path exists, then examining its
second-order boundary either produces a shorter positive type-flow path or a
direct natural-shadow match \(H_a\cong H_b\).
\end{conjecture}

This is the broader version of exchange-chain shortening.  It says that
second-deletion discrepancies cannot merely circulate at the next level; they
must expose a shorter obstruction.  If this higher-order shortening fails,
then the approach has likely run into genuine homology of the colored
deletion complex, which would explain why the current strategy cannot close
the theorem without a new idea.

\begin{conjecture}[No closed fixed-host singleton flow]
In the fixed-host singleton switch, the two natural shadow types
\[
  H_a=(K-a,S\setminus\{a\},T),
  \qquad
  H_b=(K-b,S\setminus\{b\},T)
\]
cannot both be routed entirely through detours in \(T\), avoiding the direct
positions \(b\) and \(a\), respectively.
\end{conjecture}

The computational evidence suggests a still stronger and cleaner statement.

\begin{conjecture}[Fixed-host direct-shadow theorem]
In every fixed-host singleton switch with equal colored deletion decks,
\[
  (K-a,S\setminus\{a\},T)
  \cong
  (K-b,S\setminus\{b\},T).
\]
Equivalently, the natural shadow of the moved vertex on one side always
matches the natural shadow of the moved vertex on the other side.
\end{conjecture}

The direct-shadow theorem immediately implies no closed singleton flow.  It
is stronger than what the induction needs, but it may be the right local
lemma: the deck appears unable to route the natural singleton shadow through
the common second color at all.

This is now the most local detour-selection target.  A counterexample would
not merely need deletion-similar vertices in \(T\); it would need a balanced
population of absorbers in the same color and degree-profile layer, with all
proper moving-vertex profiles still agreeing.  The next invariant to try is a
two-step count relative to \(T\): vertices adjacent to exactly one of \(a,b\)
should create signed shifts in these detour cards, just as the witness sets
\(W^-\) and \(W^+\) did in the original single-switch analysis.

\subsection{Witness charges in the fixed-host singleton switch}

Define the second-color two-step shadow
\[
  Q_z(U)=\sum_{x\in U\setminus\{z\}}\deg_{K-z}(x).
\]
Equivalently, \(Q_z(U)\) counts oriented length-two walks in the deleted
colored host \(K-z\) whose first vertex lies in the second color
\(U\setminus\{z\}\).  This is a numerical invariant of the deleted
two-colored card.

\begin{lemma}[Fixed-host two-step defect]
In the singleton setting \(U=T\cup\{a\}\), \(V=T\cup\{b\}\), we have
\[
  Q_a(U)-Q_a(V)=-\deg_{K-a}(b),
\]
\[
  Q_b(U)-Q_b(V)=\deg_{K-b}(a),
\]
and for every \(z\notin\{a,b\}\),
\[
  Q_z(U)-Q_z(V)=\one_{bz\in E(K)}-\one_{az\in E(K)}.
\]
\end{lemma}

\begin{proof}[Proof sketch]
For \(z\notin\{a,b\}\), the common contribution of the vertices in
\(T\setminus\{z\}\) cancels.  The remaining difference is
\[
  \deg_{K-z}(a)-\deg_{K-z}(b)
  =
  \deg_K(a)-\deg_K(b)-\one_{az\in E(K)}+\one_{bz\in E(K)}.
\]
The low-order singleton balance gives \(\deg_K(a)=\deg_K(b)\), leaving the
displayed formula.  The cases \(z=a\) and \(z=b\) are the same computation
with one of the moving vertices absent from one side's second color.
\end{proof}

Define the fixed-host witness sets
\[
  X^-=\{z\notin\{a,b\}:az\in E(K),\; bz\notin E(K)\},
\]
\[
  X^+=\{z\notin\{a,b\}:bz\in E(K),\; az\notin E(K)\}.
\]
Then degree balance gives
\[
  |X^-|=|X^+|,
\]
and the two-step defect formula says
\[
  Q_z(U)-Q_z(V)=
  \begin{cases}
    -1, & z\in X^-,\\
    +1, & z\in X^+,\\
    0, & \text{otherwise}
  \end{cases}
\]
for \(z\notin\{a,b\}\).

Thus the fixed-host singleton switch has the same signed witness geometry as
the original single-switch problem.  If \(a\) and \(b\) are twins relative to
the whole host, then the transposition \(a\leftrightarrow b\) preserves \(K\);
because \(a\in S\) if and only if \(b\in S\), and because both are outside
\(T\), it also preserves \(S\) and carries \(T\cup\{a\}\) to
\(T\cup\{b\}\).  So any counterexample must have nonempty balanced witness
sets \(X^-\) and \(X^+\).

After choosing a fixed-host trade, the same telescoping argument as cycle
balance applies to \(Q\).  On every cycle of the trade permutation, the
aligned \(Q\)-defects must sum to zero.  Hence witnesses cannot be hidden
independently: every cycle that carries a positive witness charge must also
carry a compensating negative charge, with the special vertices \(a\) and
\(b\) contributing the larger endpoint charges
\[
  -\deg_{K-a}(b),
  \qquad
  +\deg_{K-b}(a).
\]

\begin{conjecture}[No nonterminal fixed-host witness cycle]
In a fixed-host singleton counterexample, there is no trade cycle that
balances the endpoint charges and the witness charges while routing both
natural shadows through \(T\) and avoiding a zero-star-error direct edge.
\end{conjecture}

This is the fixed-host version of the edge-slide obstruction.  The smallest
non-twin case has
\[
  |X^-|=|X^+|=1.
\]
Then any nonterminal trade cycle has essentially one lost neighbor and one
gained non-neighbor to account for.  The expected contradiction is that
deleting those witnesses exposes the slide unless the direct \(a\leadsto b\)
edge already has a zero-error representative.

\subsection{Direct edges as a size drop}

The direct-edge case has an important extra feature: even if the direct
isomorphism does not have zero star error, it lowers the size of the centered
problem.

Suppose we have a direct fixed-host card edge
\[
  \eta:(K-a,S\setminus\{a\},T)
  \cong
  (K-b,S\setminus\{b\},T).
\]
Identify these two colored cards with a common colored host
\[
  (L,S_L,T_L).
\]
The two full colored hosts
\[
  (K,S,T\cup\{a\}),
  \qquad
  (K,S,T\cup\{b\})
\]
are now two one-vertex extensions of \(L\).  The re-added vertex has the same
first-color status on both sides and is second-colored on both sides.  Its
two attachment sets are
\[
  A_\eta=N_K(a)\setminus\{a\},
  \qquad
  B_\eta=\eta^{-1}(N_K(b)\setminus\{b\}),
\]
viewed as subsets of \(V(L)\).  Their symmetric difference is exactly the
direct star error of \(\eta\).

\begin{lemma}[Direct recentering drops the host size]
Under fixed-host singleton deck equality, a direct card edge produces a
colored centered state over \(L\), with
\[
  |V(L)|=|V(K)|-1,
\]
and
\[
  \PDeck_L(A_\eta)=\PDeck_L(B_\eta)
\]
in the colored sense, where the common colors \(S_L,T_L\) are part of the
card data.
\end{lemma}

\begin{proof}[Proof sketch]
The two colored full hosts have equal vertex decks by hypothesis.  The direct
card edge identifies one card on each side with the same colored host
\((L,S_L,T_L)\).  Cancelling this matched card from the two colored decks
leaves equality of the punctured colored extension decks over \(L\).  The
center has one fewer vertex because \(L\cong K-a\cong K-b\).
\end{proof}

\begin{conjecture}[Colored centered extension theorem]
Fix a finite graph \(L\) together with any finite list of vertex colors
\(\mathcal C\).  For attachment sets \(A,B\subseteq V(L)\), form the
one-vertex extensions in which the new vertex has the same prescribed color
data on both sides.  If the punctured decks of these colored extensions agree
as colored cards, then the two colored full extensions are isomorphic.
\end{conjecture}

This is not meant as a fundamentally new conjecture.  A finite list of colors
can be encoded by adjoining a bounded amount of rooted gadget data, or treated
directly in a Lean formalization as extra predicates on vertices.  The point
of naming it is bookkeeping: direct recentering from a fixed-host singleton
switch lands in the same centered problem, but with passive colors carried
along.

\begin{remark}[Simultaneous colored induction]
The proof should not appeal to the colored centered theorem as an external
black box.  Instead, run the whole centered-extension induction
simultaneously for finite colored centers.  The induction statement at size
\(n\) is:

for every finite graph \(C\) with any fixed finite list of vertex colors, and
for attachment sets \(A,B\subseteq V(C)\), if the colored punctured extension
decks agree and the added root has the same color data on both sides, then
the two colored extensions are isomorphic.

The uncolored centered theorem is the special case with no colors.  All
moving-hole, star-error, cycle-balance, and detour arguments are unchanged:
card isomorphisms are simply required to preserve the passive colors.  The
direct-edge size-drop step is then a legitimate appeal to the induction
hypothesis at size \(n-1\), not a circular appeal to a stronger theorem at
the same size.
\end{remark}

\begin{proposition}[Direct edge plus smaller colored centered theorem]
Assume the colored centered extension theorem is known for colored centers
with fewer than \(|V(K)|\) vertices.  Then any fixed-host singleton switch
with a direct card edge is solved: there is an automorphism of \(K\) preserving
\(S\) and carrying \(T\cup\{a\}\) to \(T\cup\{b\}\).
\end{proposition}

\begin{proof}[Proof sketch]
Choose a direct card edge and form the smaller colored centered state over
\((L,S_L,T_L)\).  If the direct star error is empty, the extended direct
isomorphism is already the desired automorphism.  Otherwise the smaller
colored centered theorem applies to the equal punctured decks over \(L\), so
the two one-vertex colored extensions over \(L\) are isomorphic.  Transporting
this isomorphism back through the direct recentering gives a color-preserving
isomorphism
\[
  (K,S,T\cup\{a\})\cong (K,S,T\cup\{b\}),
\]
which is exactly an element of \(\Aut(K,S)\) carrying \(T\cup\{a\}\) to
\(T\cup\{b\}\).
\end{proof}

Thus the fixed-host singleton problem splits cleanly.  Direct edges are
inductive size-drop moves.  The same-size obstruction is detour selection:
show that the natural shadows cannot both be hidden inside \(T\) forever.  In
this sense, the closed-flow conjecture is the real fixed-host singleton
heart; neighborhood selection after a direct edge can be delegated to the
smaller centered theorem.

\begin{proposition}[Inductive singleton package]
Assume the following two statements.

\begin{enumerate}[leftmargin=2em]
  \item The colored centered extension theorem holds for colored centers with
  fewer than \(|V(K)|\) vertices.

  \item The no closed fixed-host singleton flow conjecture holds for the
  fixed background \((K,S,T)\).
\end{enumerate}

Then the fixed-host singleton switch over \((K,S,T)\) is solved.
\end{proposition}

\begin{proof}[Proof sketch]
If there is a direct card edge \(a\leadsto b\), the direct-edge size-drop
proposition applies and gives the desired color-preserving automorphism.
Suppose no direct card edge exists.  By natural shadow isolation, the natural
\(a\)-shadow must be routed through some detour vertex of \(T\).  Applying the
same argument with \(a\) and \(b\) interchanged, the natural \(b\)-shadow must
also be routed through \(T\).  This is exactly the closed singleton flow
forbidden by the second assumption.  Hence the no-direct case is impossible,
and the direct case has already been handled.
\end{proof}

This is a useful strategic compression.  We no longer need to prove
fixed-host singleton rigidity in one breath.  It is enough to prove the
same-size detour-selection statement; all direct failures are passed down to a
smaller centered problem.  This is the first place where a genuine induction
on the number of vertices looks structurally plausible.

Consequently, if no automorphism of the background \((K,S,T)\) sends \(a\) to
\(b\), then any rooted separator between \(a\) and \(b\) must use the whole
host.  The singleton obstruction is therefore a full-size pseudo-similarity:
all proper moving-vertex profiles agree, and the only possible distinction is
the full colored orbit.  This is precisely the boundary where the detour-flow
and star-error arguments have to enter.

\begin{proposition}[Fixed-host reconstruction implies neighborhood selection]
The equal-size fixed-host subset reconstruction conjecture implies the
neighborhood selection conjecture.
\end{proposition}

\begin{proof}[Proof sketch]
After choosing a direct card isomorphism and transporting to
\[
  K=C-p,
\]
neighborhood selection asks whether
\[
  U=N_C(p)\setminus\{p\}
  \quad\text{and}\quad
  V=\psi^{-1}(N_C(q)\setminus\{q\})
\]
are in the same \(\Aut(K,S)\)-orbit.  Direct recentering gives an anchored
state with equal punctured decks.  Once the distinguished first root is kept
as part of the anchored data, the subdeck obtained by deleting ordinary
vertices of \(K\) is precisely the fixed-host deletion deck, so we get
\[
  \Deck_{K,S}(U)=\Deck_{K,S}(V).
\]
The two subsets have equal size because \(\deg_C(p)=\deg_C(q)\).  Fixed-host
subset reconstruction therefore gives the desired orbit equality.
\end{proof}

\subsection{Pattern-sum invariants}

Let \(P\) be a finite graph whose vertices carry two colors: the first color
corresponds to membership in \(S\), and the second color corresponds to
membership in \(U\).  Define the pattern-sum invariant
\[
  \pi_P(U)
\]
to be the number of induced subsets \(R\subseteq V(K)\) such that the
two-colored induced structure
\[
  (K[R],S\cap R,U\cap R)
\]
is isomorphic to \(P\).

The full list of pattern-sum invariants determines the
\(\Aut(K,S)\)-orbit of \(U\): in particular, the full-size pattern
\[
  (K,S,U)
\]
separates \(U\) from any subset not in its orbit.  The fixed-host deck and
anchored Kelly pressure explain which of these invariants are already forced.

\begin{lemma}[Proper pattern sums are deck-controlled]
Under the hypotheses of equal-size fixed-host deck equality,
\[
  \Deck_{K,S}(U)=\Deck_{K,S}(V),
  \qquad
  |U|=|V|,
\]
we have
\[
  \pi_P(U)=\pi_P(V)
\]
for every two-colored pattern \(P\) with fewer vertices than \(K\).
\end{lemma}

\begin{proof}[Proof sketch]
The invariant \(\pi_P(U)\) is exactly a proper colored induced-pattern count
in the fixed host.  Anchored Kelly pressure reconstructs this count from the
fixed-host deletion deck.  Equal decks therefore give equal proper pattern
sums.
\end{proof}

Thus a fixed-host counterexample is exactly a pair of equal-size subsets
\(U,V\) that agree on every proper pattern-sum invariant but disagree on their
full orbit.  This is an invariant-theoretic restatement of the full-size
colored hypomorphism obstruction.

\begin{conjecture}[No full-degree-only subset obstruction]
For every finite colored host \((K,S)\), equal-size subsets \(U,V\) that agree
on all proper pattern-sum invariants are in the same \(\Aut(K,S)\)-orbit.
\end{conjecture}

This conjecture is equivalent in spirit to fixed-host subset reconstruction,
but it removes the card language.  It says that equal-size subsets of a fixed
colored host have no orbit distinction visible only at full degree.  If
false, a counterexample would be a very explicit finite pattern-counting
phenomenon: two nonconjugate subsets with identical proper colored induced
substructure statistics.

\subsection{Recursive view of neighborhood selection}

Fix a direct card isomorphism
\[
  \psi:C-p\cong C-q,\qquad \psi(S)=S,
\]
and transport the \(q\)-side back to
\[
  K=C-p.
\]
Inside \(K\), define two neighborhood attachment sets
\[
  U=N_C(p)\setminus\{p\},
  \qquad
  V=\psi^{-1}(N_C(q)\setminus\{q\}).
\]
The direct card isomorphism extends over \(p\mapsto q\) exactly when
\[
  U=V.
\]
If \(U\neq V\), then the direct recentering produces a new centered state
whose defect is
\[
  U\triangle V.
\]
In fact the new center is
\[
  D=\Ext(K,S),
\]
and the two re-added vertices are attached to
\[
  U\cup\{\ast\}
  \qquad\text{and}\qquad
  V\cup\{\ast\},
\]
where \(\ast\) is the root of \(\Ext(K,S)\).  The root belongs to both
attachment sets, so the new defect is still \(U\triangle V\).

Thus neighborhood selection is exactly the assertion that the direct
recentered state is terminal.  If it is not terminal, the star error
\[
  |U\triangle V|
\]
is the new defect.  In the single-switch base case, parity forces this new
defect to be either \(0\) or at least \(2\).  Exact two-error is precisely the
case in which \(U\triangle V=\{a,b\}\), giving the slide from the switched
pair \((p,q)\) to the witness pair \((a,b)\).

This recursive view explains the slide graph: it is the graph of repeated
direct recenterings in the special situation where the defect remains \(2\).
The zero-companion lemma says that this graph has no nonterminal vertices
reachable from a genuine same-deck single switch.

\subsection{Root-anchored centered states}

The direct recentering produces a special kind of centered state.  The new
center is
\[
  D=\Ext(K,S),
\]
with distinguished root \(\ast\), and both new attachment sets contain
\(\ast\):
\[
  A'=U\cup\{\ast\},
  \qquad
  B'=V\cup\{\ast\}.
\]
Thus the defect avoids the distinguished root.  Moreover, automorphisms of
\(D\) fixing \(\ast\) are exactly automorphisms of \(K\) preserving \(S\):
\[
  \Aut(D,\ast)\cong \Aut(K,S).
\]
So neighborhood selection is equivalent to saying that the root-anchored
state \((D,A',B')\) is terminal by an automorphism fixing the distinguished
root.

\begin{lemma}[Direct recentering preserves punctured deck equality]
In the single-switch setting, suppose a direct root-preserving card
isomorphism identifies
\[
  \Ext(C-p,S)
  \cong
  \Ext(C-q,S)
\]
with a common center \(D=\Ext(K,S)\).  Let
\[
  A'=U\cup\{\ast\},
  \qquad
  B'=V\cup\{\ast\}
\]
be the two recentered attachment sets described above.  Then
\[
  \PDeck_D(A')=\PDeck_D(B').
\]
\end{lemma}

\begin{proof}[Proof sketch]
The recentered presentations describe the same two full graphs as before:
\[
  \Ext(C,S\cup\{p\})\cong \Ext(D,A'),
  \qquad
  \Ext(C,S\cup\{q\})\cong \Ext(D,B').
\]
The original graphs have equal full decks, and the direct card match
identifies one occurrence of \(D\) on each side.  Cancelling that common card
from the full decks gives equality of the punctured extension decks over
\(D\).
\end{proof}

\begin{conjecture}[Root-anchored centered theorem]
Let \(D\) be a centered graph with a distinguished vertex \(\ast\).  Suppose
\[
  A'=\{\ast\}\cup U,
  \qquad
  B'=\{\ast\}\cup V,
\]
and
\[
  \PDeck_D(A')=\PDeck_D(B').
\]
Then either \(U\) and \(V\) are in the same orbit under
\(\Aut(D,\ast)\), or there is a recentering move that keeps the distinguished
root in both attachment sets and strictly decreases \(|U\triangle V|\).
\end{conjecture}

If this anchored theorem holds, then neighborhood selection follows in the
single-switch situation: the direct recentering gives a root-anchored state
with equal punctured decks.  If it is terminal, the stabilizer-orbit lemma
selects an extendable direct card isomorphism.  If it descends, then in the
defect-\(2\) base case the only possible strict descent is defect \(0\), again
giving the desired zero-error companion.

This anchored version may be easier than the full centered theorem because
the distinguished root is common to both attachment sets.  It removes one
source of ambiguity: the defect lives entirely among old vertices of \(K\),
while \(\ast\) records the original color set \(S\).

\begin{proposition}[Detour selection plus anchored descent proves the base case]
Assume the following two statements.

\begin{enumerate}[leftmargin=2em]
  \item In every single-switch state with equal punctured extension decks, the
  rooted card-match graph contains the direct edge \(p\leadsto q\).

  \item The root-anchored centered theorem holds for the direct recentered
  state.
\end{enumerate}

Then the single-switch rigidity conjecture holds.
\end{proposition}

\begin{proof}[Proof sketch]
By the first assumption, choose a direct root-preserving card isomorphism
\[
  \psi:C-p\cong C-q,\qquad \psi(S)=S.
\]
Direct recentering gives the root-anchored state
\[
  (D,A',B')
  =
  \bigl(\Ext(K,S),\, U\cup\{\ast\},\, V\cup\{\ast\}\bigr)
\]
with
\[
  \PDeck_D(A')=\PDeck_D(B').
\]
By the anchored theorem, either \(U\) and \(V\) are in the same
\(\Aut(D,\ast)\)-orbit, or there is a strict anchored descent.  Iterating the
anchored theorem must terminate because the anchored defect is finite.  At
termination the two anchored full extensions are isomorphic, so the original
single-switch full extensions are isomorphic as well.  If termination happens
already at the first anchored state, the stabilizer-orbit formulation gives
an extendable direct card isomorphism, hence an automorphism of \(C\) carrying
\(S\cup\{p\}\) to \(S\cup\{q\}\).
\end{proof}

\subsection{Two-root interpretation}

The root-anchored state can also be read as a graph with two adjacent special
vertices.  Start with \(K=C-p\).  Add a first root \(r\) adjacent to \(S\).
Then add a second root \(x\) adjacent to \(r\) and to \(U\).  The resulting
graph is
\[
  \Ext(\Ext(K,S),U\cup\{r\}).
\]
The comparison state replaces \(U\) by \(V\):
\[
  \Ext(\Ext(K,S),V\cup\{r\}).
\]
Thus the neighborhood-selection problem is a two-color reconstruction
problem on \(K\): the first color \(S\) is encoded by the neighbors of \(r\),
and the second color \(U\) or \(V\) is encoded by the neighbors of \(x\).  The
two roots are adjacent in both graphs.

In this language, an automorphism fixing \(r\) and carrying \(U\) to \(V\) is
exactly an isomorphism of the two-root structures that fixes the first root
and sends the second root to the second root.  A failure of neighborhood
selection would therefore be a pair of two-colorings \((S,U)\) and \((S,V)\)
of the same graph \(K\) whose associated two-root graphs have the same deck,
but whose second colors are not in the same \(\Aut(K,S)\)-orbit.

This is a useful warning as well as a simplification.  General colored graph
reconstruction is not obviously easier than graph reconstruction.  The hope
is that the present two-root situation is special: the first root is common,
the second root is adjacent to it on both sides, and the state arose from a
single-switch comparison, so the second-color defect begins in the tightly
controlled neighborhood set \(U\triangle V\).

\subsection{Minimal bad form for the single-switch theorem}

Putting the pieces together, a counterexample to the single-switch theorem
would have to satisfy a rigid list of conditions.

\begin{enumerate}[leftmargin=2em]
  \item \(A=S\cup\{p\}\), \(B=S\cup\{q\}\), and
  \(\PDeck_C(A)=\PDeck_C(B)\).

  \item \(p\) and \(q\) are degree-balanced:
  \[
    \deg_C(p)=\deg_C(q).
  \]

  \item \(p\) and \(q\) are not twins outside \(\{p,q\}\), since otherwise
  twin collapse gives the automorphism.

  \item The witness sets
  \[
    W^-=\{w:pw\in E(C),\;qw\notin E(C)\},
    \qquad
    W^+=\{w:qw\in E(C),\;pw\notin E(C)\}
  \]
  have equal nonzero size.

  \item The witnesses are balanced relative to the common color:
  \[
    \deg_S(p)=\deg_S(q).
  \]

  \item More generally, \(p\) and \(q\) have identical proper root-link
  profiles: every proper induced pattern using the added root occurs equally
  often in \(\Ext(C,S\cup\{p\})\) and \(\Ext(C,S\cup\{q\})\).

  \item The rooted card-match graph admits perfect matchings, but no direct
  edge \(p\leadsto q\) with zero star error.

  \item If the natural shadow \(\Ext(C-p,S)\) is not isolated, every detour
  absorbing it must pass through vertices of \(S\) with the same degree and
  compatible rooted triangle/two-step profiles.

  \item If a direct card edge \(p\leadsto q\) exists, every corresponding
  \(S\)-preserving isomorphism \(C-p\cong C-q\) must fail neighborhood
  selection, producing a nonempty even neighborhood defect.
\end{enumerate}

The computation through six vertices says this normal form is empty.  The
proof campaign is now to show it is empty in general.  The two plausible
routes are:
\[
  \text{detour selection/isolation}
  \quad\text{or}\quad
  \text{neighborhood selection}.
\]
Either one would prove the single-switch base case.

This is a smaller target than the full single-switch rigidity conjecture.  It
does not ask us to understand the whole deck at once; it only asks us to show
that the deck equality forces one of the two defect-deleted shadows to match
with no old-card neighborhood error.

\subsection{Shape of an exact two-error}

The first bad case to exclude is therefore very specific.  Suppose a
root-preserving match based at \(p\) gives a bijection \(f=f_p\) with
\[
  f(A)=B
\]
and
\[
  |\Err_p(f)|=2.
\]
Let \(v=f(p)\).  Since the two compared neighborhoods have the same size,
there are unique vertices \(a,b\in V(C)\setminus\{p\}\) such that, after
transport by \(f\),
\[
  a\in N_C(p),
  \qquad
  f(a)\notin N_C(v),
\]
and
\[
  b\notin N_C(p),
  \qquad
  f(b)\in N_C(v).
\]
Every other adjacency from \(p\) is preserved by \(f\), and every adjacency
not involving \(p\) is preserved by \(f\).  Thus \(f\) is an isomorphism from
\(C\) to the graph obtained from \(C\) by deleting the edge \(pa\) and adding
the edge \(pb\), after relabelling by \(f\).  In other words, the obstruction
is a single edge-slide at the deleted vertex.

This suggests a more concrete base-case attack.  A single-switch
counterexample with no zero-error match would have to make every
defect-deleted card look correct while hiding an edge-slide at the missing
vertex.  But the two witnesses \(a\) and \(b\) reappear as actual deleted
vertices in their own shadows.  Deleting \(a\) removes the lost-edge witness;
deleting \(b\) removes the gained-edge witness.  The deck should have to
match these two shadows in a way that is compatible with the same color
transport, and that compatibility is where the contradiction ought to live.

\begin{conjecture}[Edge-slide exposure]
In the single-switch setting, an exact two-error match at \(p\) forces a
distinguishable imbalance between the shadows obtained by deleting the lost
neighbor \(a\) and the gained non-neighbor \(b\), unless a zero-error match
already exists.
\end{conjecture}

\begin{conjecture}[Zero companion for edge-slides]
In the single-switch setting, suppose a defect-rooted root-preserving card
match has exact star error \(2\).  Then there is another defect-rooted
root-preserving card match with star error \(0\).
\end{conjecture}

This is the sharpened base-case target suggested by computation.  Exact
two-error matches are not themselves contradictory; they appear even in small
examples where the switched colorings are globally isomorphic.  The claim is
that an exact edge-slide cannot be isolated.  If the deck permits the hidden
edge \(pa\) to slide to \(pb\), then the other shadows must also contain a
companion alignment that repairs the slide completely.

\begin{remark}[Fallback: edge-slide descent]
The zero-error target may be stronger than necessary.  If a defect-\(2\)
single-switch state admits an exact two-error recentering, then recentering
does not decrease the raw defect; the new defect again has size \(2\).  But it
changes the switched pair.  In the notation above, the new switched vertices
are precisely the lost/gained witnesses \(a\) and \(b\) inside the punctured
card.  Thus an exact two-error move is not noise: it moves the single switch
from \((p,q)\) to the local edge-slide witness pair \((a,b)\).

This suggests a secondary base-case potential.  For a switched pair \(p,q\),
define
\[
  \omega_C(p,q)
  =
  \frac{1}{2}
  \left|
    \{w\notin\{p,q\} :
      (pw\in E(C))\not\leftrightarrow(qw\in E(C))
    \}
  \right|,
\]
so the comparison ignores the possible edge \(pq\).  The twin case is
\(\omega_C(p,q)=0\).  The smallest non-twin case is
\(\omega_C(p,q)=1\), exactly one lost witness and one gained witness.  A
weaker but perhaps more realistic base-case theorem would be: every exact
two-error recentering either gives an isomorphism or moves to a new
single-switch state with smaller value of some refinement of \(\omega\), for
example the multiset of rooted profiles of the witness pair.  Then defect
\(2\) would terminate by this secondary potential even without an immediate
zero-error match.
\end{remark}

\subsection{The slide graph at defect two}

The fallback can be organized as a finite directed graph.  A vertex of the
slide graph is a defect-\(2\) centered state
\[
  (C,S,p,q),
  \qquad
  A=S\cup\{p\},\quad B=S\cup\{q\},
\]
with \(\PDeck_C(A)=\PDeck_C(B)\).  A terminal vertex is one admitting a
zero-error defect-rooted match.  A directed edge is an exact two-error
recenter:
\[
  (C,S,p,q)\longrightarrow (D,S',a,b),
\]
where \(a\) is the lost neighbor and \(b\) is the gained non-neighbor in the
edge-slide description above.  Here \(D\) is the matched punctured card used
as the new center, and the new attachment sets have the form
\[
  A'=S'\cup\{a\},
  \qquad
  B'=S'\cup\{b\}.
\]

The single-switch base case follows from either of the following statements.

\begin{conjecture}[No nonterminal slide cycle]
Every vertex in the defect-\(2\) slide graph reaches a terminal vertex.
\end{conjecture}

\begin{conjecture}[Slide potential]
There is a finite-ranking potential \(\Omega(C,S,p,q)\), defined on
defect-\(2\) states, such that every nonterminal exact two-error recentering
strictly decreases \(\Omega\) after choosing the recentering appropriately.
\end{conjecture}

The witness sets give the first candidate for \(\Omega\).  For a switched
pair \(p,q\), let
\[
  W^-=\{w: pw\in E(C),\;qw\notin E(C)\},
  \qquad
  W^+=\{w: qw\in E(C),\;pw\notin E(C)\},
\]
again outside \(\{p,q\}\).  The quantity
\[
  \omega_C(p,q)=|W^-|=|W^+|
\]
is the coarsest measure of non-twinness.  If \(\omega=0\), we are terminal by
twin collapse.  If \(\omega=1\), every exact two-error slide uses the unique
pair of witnesses.  If \(\omega>1\), one can hope to choose a witness pair
whose recentering lowers \(\omega\), or at least lowers the signed multiset of
rooted witness profiles.

This reframes the base case as a small dynamical problem: a counterexample is
not just a pair of switched vertices.  It is a closed slide system in which
the deck keeps moving the switch from one witness pair to another without ever
reaching twins or a zero-error card match.

The most promising invariants for exposing the edge-slide are rooted
two-step and triangle counts.  In a shadow \(\Ext(C-w,A\setminus\{w\})\),
the number of triangles using the root is exactly
\[
  |E(C[A\setminus\{w\}])|.
\]
Similarly, the number of rooted length-two paths from the root through the
card records
\[
  \sum_{x\in A\setminus\{w\}} \deg_{C-w}(x).
\]
For \(A=S\cup\{p\}\) and \(B=S\cup\{q\}\), these quantities change in
controlled ways when \(w\) is one of \(p,q,a,b\).  The base-case proof should
try to show that an edge-slide cannot preserve all of these rooted local
profiles across the full punctured deck.

Here are the aligned formulas.  Write
\[
  \deg_R(x)=|N_C(x)\cap R|
\]
for the number of neighbors of \(x\) in a set \(R\subseteq V(C)\).  Also write
\[
  T_w(A)=|E(C[A\setminus\{w\}])|
\]
for the number of root-using triangles in the \(w\)-deleted shadow, and
\[
  P_w(A)=\sum_{x\in A\setminus\{w\}}\deg_{C-w}(x)
\]
for the number of oriented root-using two-step walks whose first step leaves
the root.  In the single-switch case \(A=S\cup\{p\}\) and
\(B=S\cup\{q\}\), the aligned triangle differences are
\[
  T_p(A)-T_p(B)=-\deg_S(q),
\]
\[
  T_q(A)-T_q(B)=\deg_S(p),
\]
and, for \(w\notin\{p,q\}\),
\[
  T_w(A)-T_w(B)
  =
  \deg_{S\setminus\{w\}}(p)
  -
  \deg_{S\setminus\{w\}}(q).
\]
The aligned two-step differences are
\[
  P_p(A)-P_p(B)=-\deg_{C-p}(q),
\]
\[
  P_q(A)-P_q(B)=\deg_{C-q}(p),
\]
and, for \(w\notin\{p,q\}\),
\[
  P_w(A)-P_w(B)=\deg_{C-w}(p)-\deg_{C-w}(q).
\]
Using degree balance, this last formula becomes
\[
  P_w(A)-P_w(B)
  =
  \one_{qw\in E(C)}-\one_{pw\in E(C)}.
\]
Thus every vertex \(w\) that distinguishes \(p\) from \(q\) produces a
visible \(\pm 1\) shift in the aligned rooted two-step count.

Define the witness sets
\[
  W^- = \{w\notin\{p,q\}: pw\in E(C),\; qw\notin E(C)\},
  \qquad
  W^+ = \{w\notin\{p,q\}: qw\in E(C),\; pw\notin E(C)\}.
\]
Equivalently,
\[
  W^- = (N_C(p)\setminus N_C(q))\setminus\{q\},
  \qquad
  W^+ = (N_C(q)\setminus N_C(p))\setminus\{p\}.
\]
Degree balance gives
\[
  |W^-|=|W^+|.
\]
For \(w\notin\{p,q\}\), the two-step formula says
\[
  P_w(A)-P_w(B)=
  \begin{cases}
    -1, & w\in W^-,\\
    +1, & w\in W^+,\\
    0, & \text{otherwise.}
  \end{cases}
\]
So a non-twin single-switch counterexample must hide a balanced collection of
negative and positive witness shadows.  The smallest possible obstruction has
\[
  |W^-|=|W^+|=1,
\]
which is exactly the edge-slide picture above.

After choosing a centered trade, the cycle-balance lemma below makes this
stronger.  For the rooted two-step invariant \(P\), every cycle of the trade
permutation must contain balanced \(P\)-defect.  Ignoring the special vertices
\(p\) and \(q\), that means a cycle cannot contain a \(W^+\)-witness without
also containing compensating \(W^-\)-defect, and conversely.  Thus the deck
cannot hide witnesses independently; it must organize them into balanced
cycles of the matching permutation.

More explicitly, let
\[
  r=\deg_{C-p}(q)=\deg_{C-q}(p).
\]
Then the aligned \(P\)-defect is the signed charge distribution
\[
  \delta_P(p)=-r,\qquad \delta_P(q)=+r,
\]
\[
  \delta_P(w)=-1\quad(w\in W^-),
  \qquad
  \delta_P(w)=+1\quad(w\in W^+),
\]
and \(\delta_P(w)=0\) for all remaining vertices.  Every cycle of any
centered trade must have total charge zero.  Thus a minimal single-switch
counterexample is not just balanced globally; its witness charges must be
partitioned into zero-sum cycles compatible with the full rooted card types.
For \(\omega_C(p,q)=1\), this leaves very few combinatorial possibilities.

Deck equality says that the multiset of the full rooted profiles
\[
  \{\!\{\, (T_w(A),P_w(A),\ldots) : w\in V(C) \,\}\!\}
\]
equals the corresponding \(B\)-side multiset.  The aligned formulas show what
has to be hidden by the matching.  In particular, if \(p\) and \(q\) have
different adjacency profiles into \(S\), the rooted triangle shadows already
force a nontrivial rematching pattern.  If all these first moments agree, then
\(p\) and \(q\) are already very similar relative to \(S\), which is exactly
the situation where a zero-error automorphism becomes more plausible.

\begin{lemma}[Twin collapse]
If \(p\) and \(q\) have the same neighbors outside \(\{p,q\}\), then the
transposition swapping \(p\) and \(q\) and fixing all other vertices is an
automorphism of \(C\).  Consequently it carries
\[
  S\cup\{p\}
  \quad\text{to}\quad
  S\cup\{q\},
\]
so the single-switch case is solved.
\end{lemma}

\begin{proof}[Proof sketch]
All edges not incident to \(p\) or \(q\) are fixed.  For any
\(x\notin\{p,q\}\), the edge \(px\) exists if and only if \(qx\) exists, so
the transposition preserves the edges from the pair \(\{p,q\}\) to the rest
of the graph.  The edge \(pq\), if present, is fixed as an unordered edge.
\end{proof}

Thus a single-switch counterexample must be non-twin: there is some witness
vertex \(x\) adjacent to exactly one of \(p,q\).  The low rooted moments above
measure such witnesses.  The remaining difficulty is that deck equality may
permute the shadows, so the proof has to show that the witness cannot be
hidden globally by a rematching of all cards.

So the one-step descent conjecture has a concrete form: if the identity
comparison between \(A\) and \(B\) has color error
\[
  A\triangle B,
\]
then some card match based at a defect vertex should give a color-preserving
bijection with \(\widehat\psi(A)=B\) whose one-star graph error is smaller
than that color error.  If the one-star error is empty, \(\widehat\psi\) is
already the desired automorphism of \(C\).

\begin{conjecture}[One-step defect descent]
Suppose
\[
  \PDeck_C(A)=\PDeck_C(B)
\]
and \(A\) and \(B\) are not in the same \(\Aut(C)\)-orbit.  Then there exist
\[
  u\in A\triangle B,\qquad v\in V(C),
\]
and a root-preserving isomorphism
\[
  \theta:
  \Ext(C-u,A\setminus\{u\})
  \cong
  \Ext(C-v,B\setminus\{v\})
\]
such that the recentered defect satisfies
\[
  |\Delta_\theta(u,v)| < |A\triangle B|.
\]
\end{conjecture}

\begin{proposition}[Descent proves the centered theorem]
The one-step defect descent conjecture implies the centered extension theorem,
and hence the reconstruction conjecture.
\end{proposition}

\begin{proof}[Proof sketch]
Assume the centered extension theorem fails, and choose a counterexample
\((C,A,B)\) with \(|A\triangle B|\) minimal, so
\[
  \PDeck_C(A)=\PDeck_C(B)
  \quad\text{but}\quad
  \Ext(C,A)\not\cong\Ext(C,B).
\]
The defect is nonzero.  In particular \(A\) and \(B\) are not in the same
\(\Aut(C)\)-orbit, so one-step descent applies.  Recenter through a
root-preserving matched card to obtain a new centered presentation
\[
  (D,A',B')
\]
of the same two full extension graphs with
\[
  |A'\triangle B'| < |A\triangle B|.
\]
Because the full extension decks agree and the recentering identifies the
chosen root-deleted cards, cancellation gives
\[
  \PDeck_D(A')=\PDeck_D(B').
\]
By minimality, \(A'\) and \(B'\) lie in the same \(\Aut(D)\)-orbit, so the two
full extensions are isomorphic.  Transporting back through the recentering
shows \(\Ext(C,A)\cong\Ext(C,B)\), contradicting that \((C,A,B)\) was a
counterexample.  The centered reduction then gives reconstruction.
\end{proof}

\begin{remark}
This proposition proves the centered extension theorem, not automatically the
stronger root-preserving centered theorem.  After recentering, the original
root is just a vertex of the new center \(D\); a terminal automorphism of
\((D,A',B')\) may move it.  To prove the stronger theorem, one would need to
carry an additional marked-root invariant through the state graph.  For the
reconstruction conjecture itself, no such extra marking is needed.
\end{remark}

\subsection{The star-error campaign}

The proof campaign is now sharply focused.  Suppose a positive-defect
counterexample exists with
\[
  d=|A\triangle B|
\]
minimal.  A root-preserving centered trade supplies, for each deleted vertex
that participates in the trade, a bijection
\[
  f_u:V(C)\to V(C)
\]
with three properties:
\[
  f_u(A)=B,
\]
\[
  f_u \text{ preserves all adjacencies away from the star of } u,
\]
and
\[
  |\Err_u(f_u)|=|\Delta(u,f_u(u))|.
\]
Minimality says that no available \(f_u\), with \(u\in A\triangle B\), has
\[
  0<|\Err_u(f_u)|<d.
\]
Thus every useful one-star error is either zero, which proves the theorem, or
has size at least \(d\).  This is an extremely rigid normal form.

The next idea is to compare two such near-automorphisms.  If \(f_u\) and
\(f_w\) are two root-preserving trade maps, set
\[
  \gamma_{u,w}=f_w^{-1}f_u.
\]
Then
\[
  \gamma_{u,w}(A)=A.
\]
Moreover, \(\gamma_{u,w}\) can fail to preserve the graph structure only near
two stars:
\[
  \text{star}(u)
  \quad\text{and}\quad
  \text{star}\bigl(f_u^{-1}(f_w(w))\bigr).
\]
Indeed, \(f_u\) only creates errors at \(u\), and \(f_w^{-1}\) only creates
errors at the preimage of the \(f_w\)-hole.

\begin{conjecture}[Two-hole rigidity]
In a minimal positive-defect centered trade, the transition maps
\(\gamma_{u,w}\) cannot all be nontrivial two-star errors.  More precisely,
some comparison of two trade maps either produces an automorphism preserving
\(A\), or produces a one-star recentering move whose error is strictly smaller
than \(d=|A\triangle B|\).
\end{conjecture}

This cannot be a statement about arbitrary two-star errors.  Any transposition
of two same-colored vertices preserves the color set \(A\) and can only change
edges incident to those two vertices.  Such maps are everywhere.  The useful
rigidity has to use the whole deck-indexed family of one-star maps, together
with the fact that each \(f_u\) comes from an actual rooted card isomorphism.
So the object to rule out is not an isolated two-star near-automorphism, but a
closed coherent system of them arising from a centered trade.

This is the most plausible route to proving one-step descent.  The deck gives
many one-star near-automorphisms from \((C,A)\) to \((C,B)\).  A genuine
counterexample would force all of them to have large star error, but their
pairwise transition maps are color-preserving near-automorphisms of
\((C,A)\) with error supported on at most two stars.  The hope is that such a
large family of low-support graph errors cannot close up unless it comes from
actual automorphisms.

\subsection{An averaging target}

There is a slightly softer version of one-step descent that may be easier to
prove.  For a chosen root-preserving trade, write \(f_u\) for the extended
one-star map at \(u\), and define the defect-star energy
\[
  \mathcal E_D(f)
  =
  \sum_{u\in A\triangle B} |\Err_u(f_u)|.
\]
If one can choose the trade so that
\[
  \mathcal E_D(f) < |A\triangle B|^2,
\]
then some defect vertex \(u\) satisfies
\[
  |\Err_u(f_u)| < |A\triangle B|,
\]
which is exactly one-step descent.

Thus a minimal counterexample must satisfy the opposite inequality for every
available root-preserving trade:
\[
  \mathcal E_D(f) \geq |A\triangle B|^2.
\]
This is a useful target because \(\mathcal E_D\) is an honest sum of local
adjacency errors.  It may be possible to compare it with the color transport
cost of the same trade.  Since each \(f_u\) carries \(A\) to \(B\), the
identity color defect \(A\triangle B\) is being paid for by moving colored
vertices; the conjectural inequality says that the deck cannot force every
such payment to reappear as at least the same amount of graph error at every
defect vertex.

\subsection{Case 2: the root does not map to the root}

If \(\theta\) sends the root of the left extension to a non-root vertex on the
right, then the shadow already exhibits root ambiguity.  This should be
treated as a potential direct success: it may allow the root of
\(\Ext(C,A)\) to be identified with an old vertex of \(\Ext(C,B)\), producing
an isomorphism of the full extensions after one more extension step.

\begin{conjecture}[Root-moving lemma]
If some punctured-card isomorphism between
\(\Ext(C-u,A\setminus\{u\})\) and
\(\Ext(C-v,B\setminus\{v\})\) moves root to non-root, then either
\(\Ext(C,A) \cong \Ext(C,B)\), or there is a centered presentation with
strictly smaller defect.
\end{conjecture}

\begin{conjecture}[Root-preserving extraction]
In a minimal centered counterexample, after applying all possible root-moving
success or ambiguity-descent moves, there is a centered trade whose card
isomorphisms are root-preserving.
\end{conjecture}

This conjecture is the gateway into the charge-cycle machinery.  The intended
proof structure is:
\[
  \text{root-moving match}
  \Rightarrow
  \text{success or ambiguity descent},
\]
while
\[
  \text{no useful root-moving match}
  \Rightarrow
  \text{root-preserving trade}.
\]
Once the latter trade is chosen, the color-boundary and cycle-balance lemmas
apply.

\section{Root recognition and ambiguity}

The root-moving case deserves its own invariant.  In a centered extension
\(\Ext(C,A)\), the root is certainly a vertex whose deletion gives \(C\).
But it need not be the only such vertex.  Define the root-ambiguous set
\[
  R_C(A)
  =
  \{u\in V(C) : \Ext(C-u,A\setminus\{u\}) \cong C\}.
\]
These are the old vertices whose deletion produces a card isomorphic to the
root-deleted card.  Equivalently, they are the old vertices that can masquerade
as the root from the point of view of the deck.

A first necessary condition is obtained from edge counts.  If
\(u\in R_C(A)\), then
\[
  |E(C)|
  =
  |E(C-u)| + |A|-\one_{u\in A},
\]
so
\[
  \deg_C(u)+\one_{u\in A}=|A|.
\]
Thus a root-ambiguous old vertex has the same degree in the full extension as
the root.  Higher card invariants impose stronger versions of this condition.

If \(R_C(A)=\varnothing\), then the root-deleted card \(C\) is unique in
\(\Deck(\Ext(C,A))\).  In that case, every deck isomorphism must match root
card to root card, and all remaining work is root-preserving.

If \(R_C(A)\neq\varnothing\), then choosing \(u\in R_C(A)\) and recentering at
the card \(\Ext(C,A)-u\cong C\) produces a new centered presentation of the
same full graph in which an old vertex has become the root.  This is not a
failure mode; it is a symmetry/ambiguity of the extension.  It suggests a
second ranking parameter:
\[
  \rho(C,A,B)=|R_C(A)|+|R_C(B)|.
\]

\begin{conjecture}[Root-ambiguity descent]
If a root-moving punctured-card isomorphism does not already extend to an
isomorphism \(\Ext(C,A)\cong \Ext(C,B)\), then recentering through it decreases
either the attachment defect or the root-ambiguity parameter \(\rho\).
\end{conjecture}

This would give a lexicographic potential of the form
\[
  \bigl(|A\triangle B|,\rho(C,A,B),\text{higher card-profile data}\bigr).
\]
The root-preserving case changes attachment defect into old-card neighborhood
defect; the root-moving case should consume root ambiguity.

\section{State graph and potential functions}

It may be too optimistic to expect the raw defect size \(|A\triangle B|\) to
decrease after every useful recentering.  The root-preserving formula shows why:
after recentering at \(u\) and \(v\), the new defect is not directly the old
attachment defect; it is the mismatch between the old-card neighborhoods of
\(u\) and \(v\):
\[
  \bigl(N_C(u)\setminus\{u\}\bigr)
  \triangle
  \psi^{-1}\bigl(N_C(v)\setminus\{v\}\bigr).
\]
So a defect in the attachment set can be traded for a defect in the card
structure.

This suggests replacing the one-shot descent lemma by a finite state graph.
A state is a centered presentation
\[
  (C,A,B)
\]
of the same two full graphs \(G\) and \(H\).  A move chooses an old vertex
\(u\in V(C)\), matches the punctured card
\(\Ext(C-u,A\setminus\{u\})\) to some punctured card
\(\Ext(C-v,B\setminus\{v\})\), and recenters through a chosen isomorphism.

The centered theorem is equivalent to saying that every state with
\(\PDeck_C(A)=\PDeck_C(B)\) can reach a terminal state:
\[
  A=B
  \quad\text{or, more generally,}\quad
  A \text{ and } B \text{ are in the same } \Aut(C)\text{-orbit}.
\]

\begin{conjecture}[No positive-defect sink]
In the recentering state graph, every nonterminal state has at least one
outgoing move to a state with smaller value for some finite-ranking potential
\(\Phi(C,A,B)\).
\end{conjecture}

Candidate potentials, in increasing order of sophistication, are:
\[
  \Phi_0(C,A,B)=|A\triangle B|,
\]
\[
  \Phi_1(C,A,B)=
  \left(
    |A\triangle B|,
    \{\!\{\, \deg_C(w) : w\in A\triangle B \,\}\!\}
  \right),
\]
with lexicographic order, and a more structural potential
\[
  \Phi_2(C,A,B)=
  \{\!\{\, (C-w, A\setminus\{w\}, B\setminus\{w\}) :
      w\in A\triangle B \,\}\!\},
\]
ordered by the multiset extension of a recursively defined card-complexity
order.

The advantage of the state-graph view is that a proof need not show raw defect
always decreases.  It is enough to rule out a directed cycle of genuinely
positive-defect states in which every recentering merely moves the defect
around.  Such a positive cycle would be a ``balanced trade'' of attachment
sets and old-card neighborhoods.  The proof should aim to show that balanced
trades cannot exist.

\section{Normal form for a hypothetical counterexample}

Assume the centered extension theorem is false, and choose a counterexample
\((C,A,B)\) minimizing the eventual potential \(\Phi\).  Thus
\(\PDeck_C(A)=\PDeck_C(B)\), but
\(\Ext(C,A)\not\cong\Ext(C,B)\).  Even before \(\Phi\) is finalized, a
minimal counterexample should satisfy the following constraints.

\begin{enumerate}[leftmargin=2em]
  \item \(A\) and \(B\) are not in the same \(\Aut(C)\)-orbit.

  \item The defect \(A\triangle B\) is nonempty, and \(|A|=|B|\).

  \item For every \(u\in V(C)\), every matched punctured-card isomorphism
  \[
    \Ext(C-u,A\setminus\{u\})
    \cong
    \Ext(C-v,B\setminus\{v\})
  \]
  either fails to produce a smaller state, or it already preserves the
  putative minimal potential.

  \item In particular, if such an isomorphism is root-preserving, with
  restriction \(\psi:C-u\cong C-v\), then the recentered defect
  \[
    \bigl(N_C(u)\setminus\{u\}\bigr)
    \triangle
    \psi^{-1}\bigl(N_C(v)\setminus\{v\}\bigr)
  \]
  is nonempty; otherwise the empty-defect extension lemma gives a full
  isomorphism.

  \item If a root-moving isomorphism exists, then root ambiguity is not
  consumed by recentering.  Equivalently, root ambiguity must persist around a
  directed cycle of centered presentations.
\end{enumerate}

Thus a counterexample cannot merely have an attachment defect.  It must have a
self-sustaining trade in which attachment defects, old-card neighborhood
defects, and root ambiguities continually transform into one another without
ever vanishing.  The next proof task is to show that such a self-sustaining
trade is impossible.

\section{Trades}

It is useful to name the object that a counterexample would have to supply.

\begin{definition}[Rooted card-match graph]
Given \(C,A,B\), define the rooted card-match graph
\[
  \mathcal M(C,A,B)
\]
to be the bipartite graph with left and right vertex sets both \(V(C)\), where
there is an edge \(u\leadsto v\) if there exists a root-preserving
isomorphism
\[
  \Ext(C-u,A\setminus\{u\})
  \cong
  \Ext(C-v,B\setminus\{v\}).
\]
Each edge \(u\leadsto v\) carries a set of possible star-error values, one
for each such isomorphism.
\end{definition}

A root-preserving centered trade is exactly a perfect matching in
\(\mathcal M(C,A,B)\), together with a choice of one isomorphism for every
matched edge.  In this language, one-step descent asks for a defect-left
vertex \(u\in A\triangle B\) and an incident edge \(u\leadsto v\) carrying a
star error smaller than \(|A\triangle B|\).  The cycle-balance lemmas apply
after choosing a perfect matching, but the existence of a good recentering is
an edge-existence statement.

\begin{definition}[Centered trade]
A centered trade from \((C,A)\) to \((C,B)\) consists of a bijection
\(\tau:V(C)\to V(C)\) and, for every \(u\in V(C)\), an isomorphism
\[
  \theta_u:
  \Ext(C-u,A\setminus\{u\})
  \cong
  \Ext(C-\tau(u),B\setminus\{\tau(u)\}).
\]
It is root-preserving if every \(\theta_u\) sends root to root.
\end{definition}

Equality \(\PDeck_C(A)=\PDeck_C(B)\) is exactly the assertion that at least one
centered trade exists.  The stronger root-preserving centered theorem would
say that the existence of such a trade already forces a global automorphism of
\(C\), after possibly choosing a better trade:
\[
  \exists \alpha\in\Aut(C),\quad \alpha(A)=B.
\]

For a root-preserving trade, each \(\theta_u\) restricts to a colored-card
isomorphism
\[
  (C-u,A\setminus\{u\})
  \cong
  (C-\tau(u),B\setminus\{\tau(u)\}).
\]
Thus a root-preserving trade is precisely a hypomorphism between the two
one-colored structures \((C,A)\) and \((C,B)\).  The remaining question is
whether every such hypomorphism is induced by an actual colored isomorphism.

\begin{lemma}[Root-preserving trades preserve the deleted color]
If \((\tau,\theta_u)\) is a root-preserving centered trade, then
\[
  \tau(A)=B.
\]
\end{lemma}

\begin{proof}[Proof sketch]
The card isomorphism at \(u\) identifies the rooted degrees, so
\[
  |A|-\one_{u\in A}
  =
  |B|-\one_{\tau(u)\in B}.
\]
The full extensions have equal edge count, hence \(|A|=|B|\).  Therefore
\[
  u\in A
  \quad\Longleftrightarrow\quad
  \tau(u)\in B.
\]
Since \(\tau\) is a bijection, it carries \(A\) onto \(B\).
\end{proof}

\begin{lemma}[Cycle balance for rooted invariants]
Let \(I\) be any numerical invariant of rooted cards, and write
\[
  I_A(u)=I\bigl(\Ext(C-u,A\setminus\{u\})\bigr),
  \qquad
  I_B(u)=I\bigl(\Ext(C-u,B\setminus\{u\})\bigr).
\]
Suppose \((\tau,\theta_u)\) is a centered trade.  Define the aligned defect
\[
  \delta_I(u)=I_A(u)-I_B(u).
\]
Then on every cycle \(\Gamma\) of the permutation \(\tau\),
\[
  \sum_{u\in\Gamma}\delta_I(u)=0.
\]
\end{lemma}

\begin{proof}[Proof sketch]
The trade gives
\[
  I_A(u)=I_B(\tau(u)).
\]
Thus
\[
  \delta_I(\tau(u))
  =
  I_A(\tau(u))-I_B(\tau(u))
  =
  I_A(\tau(u))-I_A(u).
\]
Summing this identity around a \(\tau\)-cycle telescopes.
\end{proof}

This lemma is small but important.  Equality of decks gives only a multiset
equality until we choose a trade.  Once a trade is chosen, every scalar rooted
card invariant becomes a discrete potential whose aligned defects must
balance on each cycle of the trade permutation.  In the single-switch case,
using rooted two-step counts, this forces the \(+1\) and \(-1\) witness
signals to be paired cycle-by-cycle, not merely in aggregate.

\begin{corollary}[Defect vertices lie in balanced trade cycles]
For the edge-count invariant,
\[
  \delta_E(u)=\one_{u\in B}-\one_{u\in A}.
\]
Hence every cycle of a centered trade contains equally many vertices of
\(A\setminus B\) and \(B\setminus A\).  In particular, in the single-switch
case \(A=S\cup\{p\}\), \(B=S\cup\{q\}\), the vertices \(p\) and \(q\) lie in
the same cycle of the trade permutation.
\end{corollary}

\begin{proof}[Proof sketch]
The formula for \(\delta_E\) is exactly the aligned edge-count defect lemma.
Apply cycle balance to the edge-count invariant.  In the single-switch case
there is only one negative charge and one positive charge, so they must occur
in the same cycle.
\end{proof}

\begin{lemma}[Color-boundary form of a root-preserving trade cycle]
Assume the centered trade is root-preserving.  Let
\[
  x_0,x_1,\ldots,x_{m-1}
\]
be a cycle of \(\tau\), indexed so that \(\tau(x_i)=x_{i+1}\).  Put
\[
  b_i=\one_{x_i\in B}.
\]
Then
\[
  \one_{x_i\in A}=b_{i+1},
\]
and hence
\[
  x_i\in B\setminus A
  \quad\Longleftrightarrow\quad
  b_i=1,\ b_{i+1}=0,
\]
while
\[
  x_i\in A\setminus B
  \quad\Longleftrightarrow\quad
  b_i=0,\ b_{i+1}=1.
\]
Thus the defect vertices on a root-preserving trade cycle are exactly the
boundaries between \(B\)-colored and \(B\)-uncolored arcs, and the two defect
signs alternate around the cycle.
\end{lemma}

\begin{proof}[Proof sketch]
The root-preserving trade gives \(\tau(A)=B\), equivalently
\[
  x_i\in A
  \quad\Longleftrightarrow\quad
  \tau(x_i)=x_{i+1}\in B.
\]
The displayed boundary descriptions follow immediately.
\end{proof}

\subsection{Defect cycles}

The preceding lemma suggests a more economical normal form for a minimal
counterexample.  Once a root-preserving trade is chosen, each trade cycle is a
cyclic word in two colors: vertices in \(B\) and vertices outside \(B\).  The
defect vertices are exactly the color-boundaries of this cyclic word.  Thus
the total defect splits into disjoint even blocks, one block for each trade
cycle that changes color.

\begin{definition}[Defect cycle]
A defect cycle of a root-preserving trade is a cycle of \(\tau\) containing at
least one vertex of \(A\triangle B\).
\end{definition}

\begin{conjecture}[Atomic trade reduction]
If a centered counterexample exists, then there is one with a
root-preserving centered trade having a single defect cycle.
\end{conjecture}

This is not yet proved, but it is a promising reduction target.  If a trade
has several defect cycles, the edge-count charge already balances on each
cycle separately.  One hopes to either recenter using a vertex on one cycle to
obtain descent, or modify the trade to reduce the number of defect cycles.  In
an atomic counterexample, all positive and negative defect vertices lie on one
alternating trade cycle; the problem is then to show that at least one
boundary vertex on that cycle has small star error.

\begin{conjecture}[Cycle-local descent]
Let \(\Gamma\) be a defect cycle, and let
\[
  d_\Gamma=|\Gamma\cap(A\triangle B)|.
\]
Unless the full extensions are already isomorphic, some boundary vertex
\(u\in \Gamma\cap(A\triangle B)\) admits a root-preserving recentering whose
star error is strictly smaller than \(d_\Gamma\).
\end{conjecture}

Cycle-local descent would imply one-step defect descent immediately, since
\(d_\Gamma\leq |A\triangle B|\).  It also specializes to the single-switch
base case: then the unique defect cycle has two boundary vertices, so strict
cycle-local descent again means zero star error.

\begin{conjecture}[No nontrivial centered trade]
There is no centered trade \((C,A)\leadsto(C,B)\) with \(A\) and \(B\) in
distinct \(\Aut(C)\)-orbits that is minimal under recentering.
\end{conjecture}

This is the clean obstruction.  The moving-hole machinery is an attempt to
take a nontrivial trade and manufacture a smaller one by replacing the center
with one of the traded cards.

\section{Minimal-defect proof sketch}

Assume for contradiction that the reconstruction conjecture fails.  Choose a
counterexample \(G,H\) with the same deck and \(G \not\cong H\).  Among all
choices of matched cards \(G-x \cong H-y\), and all identifications of the
matched card with a graph \(C\), choose a centered presentation
\[
  G \cong \Ext(C,A),
  \qquad
  H \cong \Ext(C,B)
\]
with \(|A\triangle B|\) minimal.

The defect is nonzero, because zero defect gives a root-preserving isomorphism.
The needed local theorem is not that an arbitrary defect vertex works, but
that some defect vertex admits a good match.  By the one-step defect descent
lemma, choose
\[
  u\in A\triangle B,\qquad v\in V(C),
\]
and a matched punctured-card isomorphism
\[
  \Ext(C-u,A\setminus\{u\})
  \cong
  \Ext(C-v,B\setminus\{v\})
\]
whose recentering improves the chosen potential.

Apply the moving-hole analysis to this good isomorphism.

\begin{itemize}[leftmargin=2em]
  \item If the isomorphism moves root to root, the star-error formula says
  that the new defect is exactly the one-star graph error of the associated
  bijection \(f_u\), with \(f_u(A)=B\).  The descent choice makes this smaller than the
  original attachment defect, contradicting minimality.

  \item If the isomorphism moves root to non-root, the root-moving success lemma
  either gives a full isomorphism \(G \cong H\), contradiction, or decreases
  the secondary root-ambiguity part of the potential.
\end{itemize}

Therefore no positive-defect counterexample exists.  Hence every same-deck
pair is isomorphic.

\section{Ehrenfeucht-Fraisse interpretation}

The same idea can be read as an Ehrenfeucht-Fraisse strategy.  A card
alignment is a partial isomorphism with one hidden hole.  Spoiler tries to
pick a vertex in the defect set; Duplicator responds by moving the hole to a
new matched card before the defect becomes fatal.  A defect-minimizing repair
lemma would imply that Duplicator can maintain coherence through all \(n\)
rounds.

However, the direct centered-extension formulation seems cleaner for proof
work.  For finite graphs, an \(n\)-round EF win is essentially equivalent to
isomorphism, so the EF game does not reduce the difficulty; it organizes it.
The core mathematical claim remains the centered trade theorem.

\section{Lean-facing milestones}

The following milestones are ordered so that failed attempts still produce
useful formal infrastructure.

\begin{enumerate}[leftmargin=2em]
  \item Define \(\Ext(C,A)\) in Lean as a graph on \(V(C) \oplus \mathbf{1}\).

  \item Define \(\PDeck_C(A)\) as the multiset of isomorphism classes of
  \(\Ext(C-u,A\setminus\{u\})\).

  \item Prove the centered reduction:
  same deck plus one matched card gives equality of punctured extension decks.

  \item Define defect \(A\triangle B\), positive defect \(A\setminus B\), and
  negative defect \(B\setminus A\).

  \item Prove the aligned edge-count defect lemma.

  \item Generalize edge-count shadows to induced-subgraph-count shadows for
  copies using the root.

  \item Formalize moving-hole recentering: from an isomorphism of punctured
  extensions, construct the new centered presentation.

  \item Prove the star-error lemma: a root-preserving punctured-card
  isomorphism extends to a color-preserving bijection whose only graph errors
  are incident to the deleted vertex.

  \item Prove the parity lemma for star error.

  \item Formalize the single-switch base case infrastructure: degree balance,
  rooted triangle counts \(T_w\), rooted two-step counts \(P_w\), witness sets
  \(W^-\) and \(W^+\), and twin collapse.

  \item Prove single-switch root-link profile balance by subtracting
  root-avoiding subgraph counts in the common card \(C\) from deck-reconstructed
  total subgraph counts.

  \item Formalize absorbers, absorber constraints, and the natural shadow
  isolation criterion.

  \item Formalize the rooted card-match graph and the direct-edge selection
  problem.

  \item Formalize neighborhood selection as an \(\Aut(C-p,S)\)-orbit problem.

  \item Define the fixed-host subset deck
  \(\Deck_{K,S}(U)\), fixed-host trades, and the fixed-host star-error
  lemma: every matched deleted colored card extends to a color-preserving
  bijection whose graph error is concentrated in one star.

  \item Prove the fixed-host rehosting lemma, converting a fixed-host orbit
  failure into a centered attachment-defect problem over a smaller colored
  card.

  \item Formalize the fixed-host singleton switch: moving-vertex profile
  balance, low-order singleton balances, absorber constraints, and the
  second-color two-step witness charge formula.

  \item Prove the direct-edge size-drop lemma for the fixed-host singleton
  switch, and package it with no closed singleton flow into an induction step.

  \item Formalize the two-slice type-flow decomposition and the rectangular
  second-boundary identity
  \[
    \partial_O([H_b]-[H_a])
    =
    \partial_T([C_a^b]-[C_b^a]).
  \]

  \item Attack the no nonzero rectangular kernel lemma.  This is now the
  most concrete fixed-host singleton target: a proof gives the direct
  cancellation needed for the size-drop induction, while a counterexample
  gives a precise obstruction to the current campaign.  The first subtarget is
  active isolation and support escalation: every unbalanced active two-hole
  type must either produce a smaller move or become terminal exposure.  The
  name-defect case split reduces active isolation to zero-defect side cycles,
  one-sided path chords, and mixed-defect locked matches.

  \item Prove the two-hole splice lemma, or find a locked two-hole match.
  This is the local step inside the rectangular kernel argument: every
  matched two-hole card should either give a smaller obstruction, a direct
  cancellation, or a chord shortening one of the two type-flow paths.

  \item Define the balanced residual rectangular potential \(\Phi\), choose a
  \(\Phi\)-minimal residual matching, and prove the minimum-error residual
  matching lemma: no essential locked match survives after all exposed terms
  have been handled and all partner swaps that reduce
  \((\Phi,\#\text{locks})\) have been performed.

  \item Define the root-anchored centered subproblem and prove that direct
  recentering preserves punctured deck equality.

  \item Prove cycle balance for arbitrary rooted numerical invariants, with
  edge-count balance and color-boundary trade cycles as the first corollaries.

  \item Attack the single-switch theorem via detour selection/isolation plus
  neighborhood selection.

  \item Formalize transition maps between two star-error bijections, and prove
  the two-hole error-support bound.

  \item Attack one-step defect descent, with two-hole rigidity as the main
  local obstruction.

  \item Attack the root-moving success lemma as the secondary ambiguity case.
\end{enumerate}

\section{Current gap}

The proof currently hinges on four increasingly local claims.

\begin{enumerate}[leftmargin=2em]
  \item Root-moving matches either give an isomorphism or decrease a
  root-ambiguity potential, allowing us to extract a root-preserving trade in
  a minimal counterexample.

  \item In a root-preserving trade, defect lies on color-boundaries of trade
  cycles.  The desired local theorem is cycle-local descent: some boundary
  vertex on a defect cycle has star error smaller than the number of defect
  vertices on that cycle.

  \item The base case of cycle-local descent is the single-switch case.  The
  target has split into direct-edge selection and neighborhood selection.

  \item In the hard single-switch case, all proper root-link profiles agree.
  A remaining counterexample must therefore be a full-size orbit failure in
  the root-anchored two-color state.  The fixed-host trade analysis sharpens
  this: such a failure is a coherent family of color-preserving one-star
  near-automorphisms with no zero-error member.
\end{enumerate}

The next English task is to attack this sharpened fixed-host obstruction.  The
most local target is now the one-slice orbit theorem for the fixed-host
singleton switch.  The restricted Kelly lemma proves the first half:
one-slice equality forces the endpoint card type \(H_a\) to appear on the
other side of the low slice.  The remaining problem is endpoint detour
interpretation: show that if this occurrence routes through old marked
vertices \(t\in T\), then the resulting transport path integrates to an
automorphism carrying \(T\cup\{a\}\) to \(T\cup\{b\}\), or else descends to a
smaller centered obstruction.

Computationally, the six-vertex atlas scan supports the clean factorization
\[
  \text{low-slice equality}
  \Longleftrightarrow
  \text{full fixed-host deck equality}
  \Longleftrightarrow
  T\cup\{a\}\sim_{\operatorname{Aut}(K,S)}T\cup\{b\},
\]
and every state in this common class has both direct cancellations.  This is
the best current route, but the direct-cancellation clause is a small-case
bonus: a connected nine-vertex cyclic example shows that orbit success need
not imply low direct cancellation.  Orbit success itself is enough to solve
the singleton step.

The rectangular formulation remains the fallback if one-slice orbit
reconstruction fails.  It sharpens the same obstruction to a rectangular
kernel problem:
\[
  \partial_O([H_b]-[H_a])
  =
  \partial_T([C_a^b]-[C_b^a]).
\]
If the rectangular boundary forces one endpoint difference to vanish, the
fixed-host singleton theorem follows by size-drop induction.  If not, the
approach has found its likely break point: a genuine nonzero rectangular
kernel for the colored deletion operator, rather than merely a missing
low-order invariant.  The rectangular kernel has itself been localized to
named-deletion naturality.  The name-defect split says that zero-defect matches
should be removable side cycles, one-sided defects should be chords when their
corresponding extension errors vanish, while nonzero one-sided extension
errors require partial localization plus a cover/completion argument before
they inherit smaller centered-extension obstructions.  Only mixed defects
should enter the locked residual analysis.  A productive two-hole match gives
a direct cancellation, a smaller fixed-host singleton state, or a chord in one
of the two shortest type-flow paths.  Thus the immediate make-or-break target
is the essential locked case: a mixed-defective two-hole match whose four ports
stay internal to \(T\) and \(O\), whose low and high extension errors are
nonzero or color-incompatible, and which survives after all exposed
rectangular terms are handled and the residual matching minimizes
\[
  (\Phi,\#\text{essential locks}).
\]

Equivalently, the English proof has been reduced to one primary make-or-break
lemma, with four fallback rectangular lemmas behind it:
\[
  \text{one-slice equality}\Rightarrow\text{fixed-host orbit success}.
\]
If this is true, the singleton step closes directly.  If it is false, the
fallback rectangular proof requires the older four local lemmas: active
exposure isolation, support escalation to productive splice or terminal
exposure, one-sided localization with cover/completion, and exclusion of
essential mixed locked matches in a minimum-error residual transport plan.  A
failure should now appear concretely as either a one-slice orbit failure, a
terminal exposed active term, a one-sided coverage gap, or an essential locked
residual match whose correction cycles are neutral.

\end{document}
